\input{preamble}

% OK, start here.
%
\begin{document}

\title{Propriétés des espaces algébriques}


\maketitle

\phantomsection
\label{section-phantom}

\tableofcontents

\section{Introduction}
\label{section-introduction}

\noindent
Veuillez consulter
Espaces, section \ref{spaces-section-introduction},
pour une brève introduction aux espaces algébriques, et lire aussi
une partie de ce chapitre pour nos définitions et conventions de base
concernant les espaces algébriques. Dans le présent chapitre, nous introduisons
quelques notions et propriétés élémentaires des espaces algébriques. Une référence
fondamentale pour le cas des espaces algébriques quasi-séparés est
\cite{Kn}.

\medskip\noindent
L'exposé est parfois quelque peu malaisé, car nous avons choisi
de traiter d'abord des propriétés des espaces algébriques
en eux-mêmes, et seulement plus tard de celles des morphismes d'espaces
algébriques. Nous dérogeons à cette règle pour les
{\it morphismes étales} d'espaces algébriques, que nous introduisons à la
section \ref{section-etale-morphisms}.
Mais, jusqu'à cette section, lorsque nous disons qu'un morphisme possède une propriété donnée,
cela signifie automatiquement que sa source est un schéma (ou peut-être que
le morphisme est représentable).

\medskip\noindent
Certains résultats de ce chapitre (notamment ceux qui concernent les points)
seront améliorés dans le chapitre consacré aux espaces algébriques décents.


\section{Conventions}
\label{section-conventions}

\noindent
On suppose une fois pour toutes que tous les schémas appartiennent à
un grand site fppf $\Sch_{fppf}$. De plus, tous les anneaux $A$ considérés
sont tels que $\Spec(A)$ soit (isomorphe à) un
objet de ce grand site.

\medskip\noindent
Soient $S$ un schéma et $X$ un espace algébrique sur $S$.
Dans ce chapitre et le suivant, nous noterons $X \times_S X$
le produit de $X$ par lui-même (dans la catégorie des espaces
algébriques sur $S$), au lieu de $X \times X$. Nous voulons ainsi
éviter toute confusion lorsque nous changeons de schéma de base, comme dans
Espaces, section \ref{spaces-section-change-base-scheme}.


\section{Axiomes de séparation}
\label{section-separation}

\noindent
Dans cette section, nous réunissons toutes les conditions de séparation « absolues »
des espaces algébriques. Puisque, selon nos conventions, tout espace algébrique est un
espace algébrique sur un schéma de base bien déterminé, toute propriété absolue
de $X$ sur $S$ correspond à une condition imposée à $X$ considéré
comme espace algébrique sur $\Spec(\mathbf{Z})$. En voici la
formulation précise.

\begin{definition}
\label{definition-separated}
(Comparer avec Espaces, définition \ref{spaces-definition-separated}.)
Considérons un grand site fppf
$\Sch_{fppf} = (\Sch/\Spec(\mathbf{Z}))_{fppf}$.
Soit $X$ un espace algébrique sur
$\Spec(\mathbf{Z})$. Soit $\Delta : X \to X \times X$
le morphisme diagonal.
\begin{enumerate}
\item On dit que $X$ est {\it séparé} si $\Delta$ est une immersion fermée.
\item On dit que $X$ est {\it localement séparé}\footnote{Dans la
littérature, cela désigne souvent des espaces algébriques quasi-séparés et
localement séparés.} si $\Delta$ est une
immersion.
\item On dit que $X$ est {\it quasi-séparé} si $\Delta$ est quasi-compact.
\item On dit que $X$ est {\it localement quasi-séparé pour la topologie de Zariski}\footnote{
Cette notion a été proposée par B.\ Conrad.} s'il
existe un recouvrement de Zariski $X = \bigcup_{i \in I} X_i$ (voir Espaces,
définition \ref{spaces-definition-Zariski-open-covering}) tel que
chaque $X_i$ soit quasi-séparé.
\end{enumerate}
Soit $S$ un schéma appartenant à $\Sch_{fppf}$, et soit
$X$ un espace algébrique sur $S$. On dit alors que $X$ est {\it séparé},
{\it localement séparé}, {\it quasi-séparé} ou
{\it localement quasi-séparé pour la topologie de Zariski}
si $X$, considéré comme espace algébrique sur $\Spec(\mathbf{Z})$ (voir
Espaces, définition \ref{spaces-definition-base-change}),
possède la propriété correspondante.
\end{definition}

\noindent
Un espace algébrique $X$ sur $S$ qui est séparé
(au sens absolu ci-dessus) est bien séparé sur $S$ (et il en va de même
pour les autres propriétés absolues de séparation ci-dessus). Ce point sera étudié
en détail dans
Morphismes d'espaces, section \ref{spaces-morphisms-section-separation-axioms}.
Nous verrons au
lemme \ref{lemma-quasi-separated-quasi-compact-pieces}
que la propriété d'être localement quasi-séparé pour la topologie de Zariski est indépendante du schéma de base
(et donc équivalente à la notion absolue).

\begin{lemma}
\label{lemma-trivial-implications}
Soit $S$ un schéma.
Soit $X$ un espace algébrique sur $S$.
On a les implications suivantes entre les axiomes de séparation
de la définition \ref{definition-separated} :
\begin{enumerate}
\item la propriété d'être séparé entraîne toutes les autres,
\item la propriété d'être quasi-séparé entraîne celle d'être localement quasi-séparé pour la topologie de Zariski.
\end{enumerate}
\end{lemma}

\begin{proof}
La démonstration est omise.
\end{proof}

\begin{lemma}
\label{lemma-characterize-quasi-separated}
Soit $S$ un schéma. Soit $X$ un espace algébrique sur $S$.
Les conditions suivantes sont équivalentes :
\begin{enumerate}
\item $X$ est un espace algébrique quasi-séparé,
\item pour $U \to X$, $V \to X$, où $U$ et $V$ sont des schémas quasi-compacts,
le produit fibré $U \times_X V$ est quasi-compact,
\item pour $U \to X$, $V \to X$, où $U$ et $V$ sont affines,
le produit fibré $U \times_X V$ est quasi-compact.
\end{enumerate}
\end{lemma}

\begin{proof}
Le lemme
\ref{spaces-lemma-category-of-spaces-over-smaller-base-scheme}
d'Espaces permet de supposer que $S = \Spec(\mathbf{Z})$.
Puisque $U \times_X V = X \times_{X \times X} (U \times V)$
et que $U \times V$ est quasi-compact lorsque $U$ et $V$ le sont,
on voit que (1) entraîne (2). Il est clair que (2) entraîne (3).
Supposons (3). Choisissons un schéma $W$ et un morphisme étale surjectif
$W \to X$. Alors $W \times W \to X \times X$ est étale surjectif.
Il suffit donc de montrer que
$$
j : W \times_X W = X \times_{(X \times X)} (W \times W) \to W \times W
$$
est quasi-compact, voir Espaces, lemme
\ref{spaces-lemma-descent-representable-transformations-property}.
Si $U \subset W$ et $V \subset W$ sont des ouverts affines, alors
$j^{-1}(U \times V) = U \times_X V$ est quasi-compact par hypothèse.
Puisque les ouverts affines $U \times V$ forment un recouvrement ouvert affine de
$W \times W$ (Schémas, lemme \ref{schemes-lemma-affine-covering-fibre-product}),
on conclut par
Schémas, lemme \ref{schemes-lemma-quasi-compact-affine}.
\end{proof}

\begin{lemma}
\label{lemma-characterize-separated}
Soit $S$ un schéma. Soit $X$ un espace algébrique sur $S$.
Les conditions suivantes sont équivalentes :
\begin{enumerate}
\item $X$ est un espace algébrique séparé,
\item pour $U \to X$, $V \to X$, où $U$ et $V$ sont affines,
le produit fibré $U \times_X V$ est affine et
$$
\mathcal{O}(U) \otimes_\mathbf{Z} \mathcal{O}(V)
\longrightarrow
\mathcal{O}(U \times_X V)
$$
est surjectif.
\end{enumerate}
\end{lemma}

\begin{proof}
Le lemme
\ref{spaces-lemma-category-of-spaces-over-smaller-base-scheme}
d'Espaces permet de supposer que $S = \Spec(\mathbf{Z})$.
Puisque $U \times_X V = X \times_{X \times X} (U \times V)$
et que $U \times V$ est affine lorsque $U$ et $V$ le sont,
on voit que (1) entraîne (2).
Supposons (2). Choisissons un schéma $W$ et un morphisme étale surjectif
$W \to X$. Alors $W \times W \to X \times X$ est étale surjectif.
Il suffit donc de montrer que
$$
j : W \times_X W = X \times_{(X \times X)} (W \times W) \to W \times W
$$
est une immersion fermée, voir Espaces, lemme
\ref{spaces-lemma-descent-representable-transformations-property}.
Si $U \subset W$ et $V \subset W$ sont des ouverts affines, alors
$j^{-1}(U \times V) = U \times_X V$ est affine par hypothèse
et le morphisme $U \times_X V \to U \times V$ est une immersion fermée,
car l'homomorphisme d'anneaux correspondant est surjectif.
Puisque les ouverts affines $U \times V$ forment un recouvrement ouvert affine
de $W \times W$
(Schémas, lemme \ref{schemes-lemma-affine-covering-fibre-product}),
on conclut par
Morphismes, lemme \ref{morphisms-lemma-closed-immersion}.
\end{proof}







\section{Points des espaces algébriques}
\label{section-points}

\noindent
Comme il ressort clairement de
Espaces, exemple \ref{spaces-example-affine-line-translation},
il ne convient pas de définir un point d'un espace algébrique comme un monomorphisme
ayant pour source le spectre d'un corps.
Nous définissons plutôt les points comme des classes d'équivalence de morphismes de spectres
de corps, exactement comme expliqué dans
Schémas, section \ref{schemes-section-points}.

\medskip\noindent
Soit $S$ un schéma.
Soit $F$ un préfaisceau sur $(\Sch/S)_{fppf}$.
Soit $K$ un corps. Considérons un morphisme
$$
\Spec(K) \longrightarrow F.
$$
D'après le lemme de Yoneda, celui-ci est donné par un
élément $p \in F(\Spec(K))$. On dit que deux tels
couples $(\Spec(K), p)$ et $(\Spec(L), q)$
sont {\it équivalents} s'il existe
un troisième corps $\Omega$ et un diagramme commutatif
$$
\xymatrix{
\Spec(\Omega) \ar[r] \ar[d] &
\Spec(L) \ar[d]^q \\
\Spec(K) \ar[r]^p &
F.
}
$$
Autrement dit, il existe des extensions de corps
$K \to \Omega$ et $L \to \Omega$ telles que
les images de $p$ et de $q$ coïncident dans
$F(\Spec(\Omega))$. Nous omettons la vérification du fait que ceci
définit une relation d'équivalence.

\begin{definition}
\label{definition-points}
Soit $S$ un schéma. Soit $X$ un espace algébrique sur $S$.
Un {\it point} de $X$ est une classe d'équivalence de morphismes
de spectres de corps vers $X$.
L'ensemble des points de $X$ est noté $|X|$.
\end{definition}

\noindent
Notons que, si $f : X \to Y$ est un morphisme d'espaces algébriques
sur $S$, alors on dispose d'une application induite $|f| : |X| \to |Y|$ qui
envoie un représentant $x : \Spec(K) \to X$ sur le représentant
$f \circ x : \Spec(K) \to Y$.

\begin{lemma}
\label{lemma-scheme-points}
Soit $S$ un schéma. Soit $X$ un schéma sur $S$.
Les points de $X$ en tant que schéma sont canoniquement en bijection
avec les points de $X$ en tant qu'espace algébrique.
\end{lemma}

\begin{proof}
C'est Schémas, lemme \ref{schemes-lemma-characterize-points}.
\end{proof}

\begin{lemma}
\label{lemma-points-cartesian}
Soit $S$ un schéma. Soit
$$
\xymatrix{
Z \times_Y X \ar[r] \ar[d] & X \ar[d] \\
Z \ar[r] & Y
}
$$
un diagramme cartésien d'espaces algébriques sur $S$. Alors l'application entre les ensembles
de points
$$
|Z \times_Y X|
\longrightarrow
|Z| \times_{|Y|} |X|
$$
est surjective.
\end{lemma}

\begin{proof}
En effet, soient des corps $K$, $L$ et des morphismes
$\Spec(K) \to X$, $\Spec(L) \to Z$. L'hypothèse
que leurs images dans $|Y|$ coïncident signifie qu'il
existe une extension commune $M/K$ et $M/L$
telle que les morphismes
$\Spec(M) \to \Spec(K) \to X \to Y$ et
$\Spec(M) \to \Spec(L) \to Z \to Y$ coïncident.
C'est précisément la condition qui assure l'existence d'un
morphisme $\Spec(M) \to Z \times_Y X$.
\end{proof}

\begin{lemma}
\label{lemma-characterize-surjective}
Soit $S$ un schéma.
Soit $X$ un espace algébrique sur $S$.
Soit $f : T \to X$ un morphisme d'un schéma vers $X$.
Les assertions suivantes sont équivalentes :
\begin{enumerate}
\item $f : T \to X$ est surjectif (au sens de
Espaces, définition \ref{spaces-definition-relative-representable-property}),
et
\item $|f| : |T| \to |X|$ est surjectif.
\end{enumerate}
\end{lemma}

\begin{proof}
Supposons (1). Soit $x : \Spec(K) \to X$ un morphisme
du spectre d'un corps vers $X$. Par hypothèse, le morphisme de
schémas $\Spec(K) \times_X T \to \Spec(K)$ est surjectif.
Il existe donc une extension de corps $K'/K$ et un morphisme
$\Spec(K') \to \Spec(K) \times_X T$ tels que le carré
de gauche du diagramme
$$
\xymatrix{
\Spec(K') \ar[r] \ar[d] &
\Spec(K) \times_X T \ar[d] \ar[r] &
T \ar[d] \\
\Spec(K) \ar@{=}[r] &
\Spec(K) \ar[r]^-x & X
}
$$
soit commutatif. Cela montre que $|f| : |T| \to |X|$ est surjectif.

\medskip\noindent
Supposons (2). Soit $Z \to X$ un morphisme, où $Z$ est
un schéma. Il faut montrer que le morphisme de schémas $Z \times_X T \to Z$
est surjectif, c'est-à-dire que $|Z \times_X T| \to |Z|$ est surjectif.
Cela résulte de (2) et du
lemme \ref{lemma-points-cartesian}.
\end{proof}

\begin{lemma}
\label{lemma-points-presentation}
Soit $S$ un schéma.
Soit $X$ un espace algébrique sur $S$.
Soit $X = U/R$ une présentation de $X$, voir
Espaces, définition \ref{spaces-definition-presentation}.
Alors l'image de $|R| \to |U| \times |U|$ est une relation d'équivalence
et $|X|$ est le quotient de $|U|$ par cette relation d'équivalence.
\end{lemma}

\begin{proof}
L'hypothèse signifie que $U$ est un schéma, que $p : U \to X$ est un morphisme
étale surjectif, que $R = U \times_X U$ est un schéma et définit une relation
d'équivalence étale sur $U$ telle que $X = U/R$ comme faisceaux. D'après le
lemme \ref{lemma-characterize-surjective},
on voit que $|U| \to |X|$ est surjective. D'après le
lemme \ref{lemma-points-cartesian},
l'application
$$
|R| \longrightarrow |U| \times_{|X|} |U|
$$
est surjective. L'image de $|R| \to |U| \times |U|$ est donc
exactement l'ensemble des couples $(u_1, u_2) \in |U| \times |U|$
tels que $u_1$ et $u_2$ aient la même image dans $|X|$.
En réunissant ces deux assertions, on obtient le résultat du lemme.
\end{proof}

\begin{lemma}
\label{lemma-topology-points}
Soit $S$ un schéma. Il existe une unique topologie sur les ensembles de points
des espaces algébriques sur $S$ qui possède les propriétés suivantes :
\begin{enumerate}
\item si $X$ est un schéma sur $S$, la topologie sur $|X|$ est la topologie usuelle
(via l'identification du lemme \ref{lemma-scheme-points}),
\item pour tout morphisme d'espaces algébriques $X \to Y$ sur $S$,
l'application $|X| \to |Y|$ est continue, et
\item pour tout morphisme étale $U \to X$, où $U$ est un schéma,
l'application d'espaces topologiques $|U| \to |X|$ est continue et ouverte.
\end{enumerate}
\end{lemma}

\begin{proof}
Soit $X$ un espace algébrique sur $S$. Soit $p : U \to X$ un
morphisme étale surjectif, où $U$ est un schéma sur $S$.
On déclare que $W \subset |X|$ est ouvert si et seulement si $|p|^{-1}(W)$
est une partie ouverte de $|U|$. Cela définit une topologie sur $|X|$
(c'est la topologie quotient sur $|X|$, voir
Topologie, lemme \ref{topology-lemma-quotient}).

\medskip\noindent
Montrons que la topologie est indépendante du choix de
la présentation. Pour cela, il suffit de montrer que, si $U'$ est un schéma
et $U' \to X$ un morphisme étale, alors l'application $|U'| \to |X|$
(la topologie sur $|X|$ étant définie à l'aide de $U \to X$ comme ci-dessus)
est ouverte et continue, ce qui démontrera en outre (3).
Posons $U'' = U \times_X U'$ ; on a alors le diagramme commutatif
$$
\xymatrix{
U'' \ar[r] \ar[d] & U' \ar[d] \\
U \ar[r] & X
}
$$
Puisque $U \to X$ et $U' \to X$ sont étales, on voit que
$U'' \to U$ et $U'' \to U'$ sont tous deux des morphismes étales de schémas.
De plus, $U'' \to U'$ est surjectif. On obtient donc
un diagramme commutatif d'applications d'ensembles
$$
\xymatrix{
|U''| \ar[r] \ar[d] & |U'| \ar[d] \\
|U| \ar[r] & |X|
}
$$
La flèche horizontale inférieure est surjective (voir le
lemme \ref{lemma-characterize-surjective}
ou le
lemme \ref{lemma-points-presentation})
et continue par définition de la topologie sur $|X|$.
La flèche horizontale supérieure est surjective, continue et ouverte d'après
Morphismes, lemme \ref{morphisms-lemma-etale-open}.
La flèche verticale gauche est également continue et ouverte d'après
Morphismes, lemme \ref{morphisms-lemma-etale-open}.
Il en résulte formellement que la flèche verticale droite
est continue et ouverte.

\medskip\noindent
Pour achever la démonstration, prouvons (2).
Soit $a : X \to Y$ un morphisme d'espaces algébriques. D'après
Espaces, lemme \ref{spaces-lemma-lift-morphism-presentations},
on peut trouver un diagramme
$$
\xymatrix{
U \ar[d]_p \ar[r]_\alpha & V \ar[d]^q \\
X \ar[r]^a & Y
}
$$
où $U$ et $V$ sont des schémas, et où $p$ et $q$ sont étales et surjectifs.
On en déduit le diagramme
$$
\xymatrix{
|U| \ar[d]_p \ar[r]_\alpha & |V| \ar[d]^q \\
|X| \ar[r]^a & |Y|
}
$$
où toutes les flèches sauf l'horizontale inférieure sont continues et
où les deux flèches verticales sont surjectives et ouvertes. Il s'ensuit que la
flèche horizontale inférieure est continue, comme voulu.
\end{proof}

\begin{definition}
\label{definition-topological-space}
Soit $S$ un schéma. Soit $X$ un espace algébrique sur $S$.
L'{\it espace topologique sous-jacent} à $X$ est l'ensemble des points
$|X|$ muni de la topologie construite au
lemme \ref{lemma-topology-points}.
\end{definition}

\noindent
Il se trouve que cet espace topologique porte la même information
que le petit site de Zariski $X_{Zar}$ de
Espaces, définition \ref{spaces-definition-small-Zariski-site}.

\begin{lemma}
\label{lemma-open-subspaces}
Soit $S$ un schéma.
Soit $X$ un espace algébrique sur $S$.
\begin{enumerate}
\item La règle $X' \mapsto |X'|$ définit une bijection respectant les inclusions
entre les sous-espaces ouverts $X'$ (voir
Espaces, définition \ref{spaces-definition-immersion})
de $X$ et les ouverts de l'espace topologique $|X|$.
\item Une famille $\{X_i \subset X\}_{i \in I}$ de sous-espaces ouverts de $X$
est un recouvrement de Zariski (voir
Espaces, définition \ref{spaces-definition-Zariski-open-covering})
si et seulement si $|X| = \bigcup |X_i|$.
\end{enumerate}
Autrement dit, le petit site de Zariski $X_{Zar}$ de $X$ s'identifie canoniquement
à un site associé à l'espace topologique $|X|$ (voir
Sites, exemple \ref{sites-example-site-topological}).
\end{lemma}

\begin{proof}
Pour démontrer (1), construisons l'inverse de cette règle.
Supposons donc que $W \subset |X|$ soit ouvert. Choisissons une présentation
$X = U/R$ correspondant au morphisme étale surjectif
$p : U \to X$ et aux morphismes étales $s, t : R \to U$.
Par construction, on voit que $|p|^{-1}(W)$ est un
ouvert de $U$. Notons $W' \subset U$ le sous-schéma ouvert correspondant.
Il est clair que $R' = s^{-1}(W') = t^{-1}(W')$ est un ouvert de Zariski
de $R$ qui définit une relation d'équivalence étale sur $W'$.
D'après Espaces, lemme \ref{spaces-lemma-finding-opens}, le morphisme
$X' = W'/R' \to X$ est une immersion ouverte. Ainsi, $X'$ est un espace algébrique
d'après Espaces, lemme \ref{spaces-lemma-representable-over-space}.
Par construction, $|X'| = W$, c'est-à-dire que $X'$ est un sous-espace de $X$
correspondant à $W$. Cela démontre (1).

\medskip\noindent
Pour démontrer (2), notons que, si $\{X_i \subset X\}_{i \in I}$ est une famille
de sous-espaces ouverts, elle est un recouvrement de Zariski si et seulement si
$U = \bigcup U \times_X X_i$ est un recouvrement ouvert. Cela résulte de
la définition d'un recouvrement de Zariski et du fait que le morphisme
$U \to X$ est surjectif comme morphisme de préfaisceaux sur $(\Sch/S)_{fppf}$.
D'autre part, on voit que $|X| = \bigcup |X_i|$ si et seulement si
$U = \bigcup U \times_X X_i$ d'après le lemme \ref{lemma-points-presentation}
(et le fait que les projections $U \times_X X_i \to X_i$ sont surjectives
et étales). L'équivalence de (2) en résulte.
\end{proof}

\begin{lemma}
\label{lemma-factor-through-open-subspace}
Soit $S$ un schéma.
Soient $X$, $Y$ des espaces algébriques sur $S$.
Soit $X' \subset X$ un sous-espace ouvert.
Soit $f : Y \to X$ un morphisme d'espaces algébriques sur $S$.
Alors $f$ se factorise par $X'$ si et seulement si $|f| : |Y| \to |X|$
se factorise par $|X'| \subset |X|$.
\end{lemma}

\begin{proof}
D'après Espaces, lemme \ref{spaces-lemma-base-change-immersions},
on voit que $Y' = Y \times_X X' \to Y$ est une immersion ouverte.
Si $|f|(|Y|) \subset |X'|$, alors $|Y'| = |Y|$ manifestement. Ainsi, $Y' = Y$ d'après le
lemme \ref{lemma-open-subspaces}.
\end{proof}

\begin{lemma}
\label{lemma-etale-image-open}
Soit $S$ un schéma. Soit $X$ un espace algébrique sur $S$.
Soit $U$ un schéma et soit $f : U \to X$ un morphisme étale.
Soit $X' \subset X$ le sous-espace ouvert correspondant à
l'ouvert $|f|(|U|) \subset |X|$ par le
lemme \ref{lemma-open-subspaces}.
Alors $f$ se factorise par un morphisme étale surjectif $f' : U \to X'$.
En outre, si $R = U \times_X U$, alors $R = U \times_{X'} U$ et $X'$ admet
la présentation $X' = U/R$.
\end{lemma}

\begin{proof}
L'existence de la factorisation résulte du
lemme \ref{lemma-factor-through-open-subspace}.
Le morphisme $f'$ est surjectif d'après le
lemme \ref{lemma-characterize-surjective}.
Pour voir que $f'$ est étale, soit $T \to X'$ un morphisme,
où $T$ est un schéma. Alors $T \times_X U = T \times_{X'} U$,
car $X' \to X$ est un monomorphisme de faisceaux. Ainsi, la projection
$T \times_{X'} U \to T$ est étale puisque $f$ est supposé étale.
On a $U \times_X U = U \times_{X'} U$, car $X' \to X$ est un monomorphisme.
Alors $X' = U/R$ résulte de
Espaces, lemme \ref{spaces-lemma-space-presentation}.
\end{proof}

\begin{lemma}
\label{lemma-equivalence-class-point-monomorphism}
Soit $S$ un schéma. Soit $X$ un espace algébrique sur $S$.
Soient $p : \Spec(K) \to X$ et $q : \Spec(L) \to X$
des morphismes, où $K$ et $L$ sont des corps. Supposons que $p$ et $q$
déterminent le même point de $|X|$ et que $p$ soit un monomorphisme.
Alors $q$ se factorise de manière unique par $p$.
\end{lemma}

\begin{proof}
Puisque $p$ et $q$ définissent le même point de $|X|$, on voit que le schéma
$$
Y = \Spec(K) \times_{p, X, q} \Spec(L)
$$
est non vide. Comme tout changement de base d'un monomorphisme est un monomorphisme,
cela signifie que le morphisme de projection $Y \to \Spec(L)$
est un monomorphisme. Ainsi, $Y = \Spec(L)$, voir
Schémas, lemme \ref{schemes-lemma-mono-towards-spec-field}.
On en conclut que $q$ se factorise par $p$. L'unicité
résulte du fait que $p$ est un monomorphisme.
\end{proof}

\begin{lemma}
\label{lemma-points-monomorphism}
Soit $S$ un schéma. Soit $X$ un espace algébrique sur $S$.
Considérons l'application
$$
\{\Spec(k) \to X \text{ monomorphisme où }k\text{ est un corps}\}
\longrightarrow
|X|
$$
Cette application est injective.
\end{lemma}

\begin{proof}
Cela résulte du lemme \ref{lemma-equivalence-class-point-monomorphism}.
\end{proof}

\noindent
Nous verrons dans Espaces algébriques décents,
lemme \ref{decent-spaces-lemma-decent-points-monomorphism},
que l'application du lemme \ref{lemma-points-monomorphism}
est bijective lorsque $X$ est décent.

















\section{Espaces quasi-compacts}
\label{section-quasi-compact}

\begin{definition}
\label{definition-quasi-compact}
Soit $S$ un schéma.
Soit $X$ un espace algébrique sur $S$.
On dit que $X$ est {\it quasi-compact} s'il existe un morphisme
étale surjectif $U \to X$, où $U$ est un schéma quasi-compact.
\end{definition}

\begin{lemma}
\label{lemma-quasi-compact-space}
Soit $S$ un schéma.
Soit $X$ un espace algébrique sur $S$.
Alors $X$ est quasi-compact si et seulement si $|X|$ est quasi-compact.
\end{lemma}

\begin{proof}
Choisissons un schéma $U$ et un morphisme étale surjectif $U \to X$.
Nous utiliserons le lemme \ref{lemma-characterize-surjective}.
Si $U$ est quasi-compact, alors, puisque $|U| \to |X|$ est surjective,
on en conclut que $|X|$ est quasi-compact.
Si $|X|$ est quasi-compact, alors, puisque $|U| \to |X|$ est ouverte,
on voit qu'il existe un ouvert quasi-compact $U' \subset U$
tel que $|U'| \to |X|$ soit surjective (et encore étale).
D'où le résultat.
\end{proof}

\begin{lemma}
\label{lemma-finite-disjoint-quasi-compact}
Une somme disjointe finie d'espaces algébriques quasi-compacts est
un espace algébrique quasi-compact.
\end{lemma}

\begin{proof}
Cela résulte immédiatement du
lemme \ref{lemma-quasi-compact-space}
et du fait topologique correspondant.
\end{proof}

\begin{example}
\label{example-quasi-compact-not-very-reasonable}
L'espace $\mathbf{A}^1_{\mathbf{Q}}/\mathbf{Z}$ est un espace algébrique
quasi-compact.
\end{example}

\begin{lemma}
\label{lemma-space-locally-quasi-compact}
Soit $S$ un schéma.
Soit $X$ un espace algébrique sur $S$.
Tout point de $|X|$ admet un système fondamental de voisinages
ouverts quasi-compacts.
En particulier, $|X|$ est localement quasi-compact au sens de
Topologie, définition \ref{topology-definition-locally-quasi-compact}.
\end{lemma}

\begin{proof}
Cela découle formellement du fait qu'il existe un schéma
$U$ et une application surjective, ouverte et continue
$U \to |X|$ entre espaces topologiques. Plus précisément, si
$u \in U$ a pour image $x \in |X|$, alors les images des voisinages
affines de $u$ forment un système fondamental de voisinages ouverts
quasi-compacts de $x$.
\end{proof}







\section{Recouvrements spéciaux}
\label{section-special-coverings}

\noindent
Dans cette section, nous réunissons quelques lemmes immédiats sur l'existence
de recouvrements étales surjectifs d'espaces algébriques.

\begin{lemma}
\label{lemma-cover-by-union-affines}
Soit $S$ un schéma.
Soit $X$ un espace algébrique sur $S$.
Il existe un morphisme étale surjectif $U \to X$, où
$U$ est une somme disjointe de schémas affines.
On peut en outre supposer que chacun de ces schémas affines
est envoyé dans un ouvert affine de $S$.
\end{lemma}

\begin{proof}
Soit $V \to X$ un morphisme étale surjectif.
Soit $V = \bigcup_{i \in I} V_i$ un recouvrement ouvert de Zariski
tel que chaque $V_i$ soit envoyé dans un ouvert affine de $S$.
Posons alors $U = \coprod_{i \in I} V_i$, muni du morphisme induit
$U \to V \to X$. Celui-ci est étale et surjectif, car il est composé
de transformations représentables de foncteurs qui sont étales et
surjectives (d'après le principe général de
Espaces, lemme
\ref{spaces-lemma-composition-representable-transformations-property}
et de
Morphismes, lemmes \ref{morphisms-lemma-composition-surjective} et
\ref{morphisms-lemma-composition-etale}).
\end{proof}

\begin{lemma}
\label{lemma-union-of-quasi-compact}
Soit $S$ un schéma.
Soit $X$ un espace algébrique sur $S$.
Il existe un recouvrement de Zariski $X = \bigcup X_i$
tel que chaque espace algébrique $X_i$ admette un recouvrement
étale surjectif par un schéma affine. On peut en outre supposer
que chaque $X_i$ soit envoyé dans un ouvert affine de $S$.
\end{lemma}

\begin{proof}
D'après le lemme \ref{lemma-cover-by-union-affines}, on peut trouver un morphisme
étale surjectif $U = \coprod U_i \to X$, où chaque $U_i$ est affine et est envoyé
dans un ouvert affine de $S$. Soit $X_i \subset X$ le sous-espace ouvert
de $X$ tel que $U_i \to X$ se factorise par un morphisme étale surjectif
$U_i \to X_i$, voir le
lemme \ref{lemma-etale-image-open}.
Puisque $U = \bigcup U_i$, on voit que $X = \bigcup X_i$.
Comme $U_i \to X_i$ est surjectif, il s'ensuit que $X_i \to S$ se factorise par
un ouvert affine de $S$.
\end{proof}

\begin{lemma}
\label{lemma-quasi-compact-affine-cover}
Soit $S$ un schéma.
Soit $X$ un espace algébrique sur $S$.
Alors $X$ est quasi-compact si et seulement s'il
existe un morphisme étale surjectif $U \to X$
où $U$ est un schéma affine.
\end{lemma}

\begin{proof}
S'il existe un morphisme étale surjectif $U \to X$ où $U$
est affine, alors $X$ est quasi-compact par la définition \ref{definition-quasi-compact}.
Réciproquement, si $X$ est quasi-compact, alors $|X|$ est quasi-compact.
Soit $U = \coprod_{i \in I} U_i$ une somme disjointe de schémas affines
munie d'un morphisme étale surjectif $\varphi : U \to X$
(lemme \ref{lemma-cover-by-union-affines}).
Alors $|X| = \bigcup \varphi(|U_i|)$ et,
par quasi-compacité, il existe des indices $i_1, \ldots, i_n$
tels que $|X| = \bigcup \varphi(|U_{i_j}|)$. Par conséquent,
$U_{i_1} \cup \ldots \cup U_{i_n}$ est un schéma affine muni d'un
morphisme étale surjectif vers $X$.
\end{proof}

\noindent
Le lemme suivant deviendra caduc lorsque nous traiterons des
morphismes séparés dans le chapitre consacré aux morphismes d'espaces algébriques.

\begin{lemma}
\label{lemma-separated-cover}
Soit $S$ un schéma.
Soit $X$ un espace algébrique sur $S$.
Soit $U$ un schéma séparé et soit $U \to X$ étale.
Alors $U \to X$ est séparé, et $R = U \times_X U$ est un schéma séparé.
\end{lemma}

\begin{proof}
Soit $X' \subset X$ le sous-espace ouvert tel que $U \to X$ se factorise
par un morphisme étale surjectif $U \to X'$, voir le
lemme \ref{lemma-etale-image-open}.
Si $U \to X'$ est séparé, alors $U \to X$ l'est aussi, voir
Espaces, lemme
\ref{spaces-lemma-composition-representable-transformations-property}
(car l'immersion ouverte $X' \to X$ est séparée d'après
Espaces, lemme
\ref{spaces-lemma-representable-transformations-property-implication}
et
Schémas, lemme \ref{schemes-lemma-immersions-monomorphisms}).
De plus, puisque $U \times_{X'} U = U \times_X U$, il suffit
de démontrer le résultat après avoir remplacé $X$ par $X'$, c'est-à-dire qu'on peut
supposer $U \to X$ surjectif.
Considérons le diagramme commutatif
$$
\xymatrix{
R = U \times_X U \ar[r] \ar[d] & U \ar[d] \\
U \ar[r] & X
}
$$
Dans la démonstration de
Espaces, lemme \ref{spaces-lemma-properties-diagonal},
nous avons vu que $j : R \to U \times_S U$ est séparé.
Le morphisme de schémas $U \to S$ est séparé, puisque $U$ est un schéma
séparé, voir
Schémas, lemme \ref{schemes-lemma-compose-after-separated}.
Par conséquent, $U \times_S U \to U$ est séparé en tant que changement de base, voir
Schémas, lemme \ref{schemes-lemma-separated-permanence}.
Le schéma $U \times_S U$ est donc séparé (d'après le même lemme).
Puisque $j$ est séparé, on voit de même que $R$ est séparé.
Ainsi, $R \to U$ est un morphisme séparé (d'après
Schémas, lemme \ref{schemes-lemma-compose-after-separated},
là encore). Par conséquent, d'après
Espaces, lemme \ref{spaces-lemma-representable-morphisms-spaces-property}
et le diagramme ci-dessus, on conclut que $U \to X$ est séparé.
\end{proof}

\begin{lemma}
\label{lemma-quasi-separated}
Soit $S$ un schéma. Soit $X$ un espace algébrique sur $S$.
S'il existe un schéma quasi-séparé $U$ et un morphisme
étale surjectif $U \to X$ tel que l'une ou l'autre des projections
$U \times_X U \to U$ soit quasi-compacte, alors $X$ est quasi-séparé.
\end{lemma}

\begin{proof}
On peut considérer $X$ comme un espace algébrique sur $\mathbf{Z}$.
Considérons le diagramme cartésien
$$
\xymatrix{
U \times_X U \ar[r] \ar[d]_j & X \ar[d]^\Delta \\
U \times U \ar[r] & X \times X
}
$$
Puisque $U$ est quasi-séparé, la projection $U \times U \to U$ est
quasi-séparée (elle s'obtient par changement de base à partir d'un morphisme
quasi-séparé de schémas, voir Schémas, lemme \ref{schemes-lemma-separated-permanence}).
L'hypothèse du lemme implique donc que $j$ est quasi-compact d'après
Schémas, lemme \ref{schemes-lemma-quasi-compact-permanence}.
D'après Espaces, lemme \ref{spaces-lemma-representable-morphisms-spaces-property},
on voit que le morphisme $\Delta$ est quasi-compact, comme voulu.
\end{proof}

\begin{lemma}
\label{lemma-quasi-separated-quasi-compact-pieces}
Soit $S$ un schéma.
Soit $X$ un espace algébrique sur $S$.
Les assertions suivantes sont équivalentes :
\begin{enumerate}
\item $X$ est localement quasi-séparé pour la topologie de Zariski sur $S$,
\item $X$ est localement quasi-séparé pour la topologie de Zariski,
\item il existe un recouvrement ouvert de Zariski $X = \bigcup X_i$
tel que, pour tout $i$, il existe un schéma affine
$U_i$ et un morphisme étale, surjectif et quasi-compact
$U_i \to X_i$, et
\item il existe un recouvrement ouvert de Zariski $X = \bigcup X_i$
tel que, pour tout $i$, il existe un schéma affine
$U_i$ qui est envoyé dans un ouvert affine de $S$ et un morphisme
étale, surjectif et quasi-compact $U_i \to X_i$.
\end{enumerate}
\end{lemma}

\begin{proof}
Supposons que les morphismes $U_i \to X_i \subset X$ soient comme en (3). Pour démontrer (4),
choisissons, pour tout $i$, un recouvrement ouvert affine fini $U_i =
U_{i1} \cup \ldots \cup U_{in_i}$ tel que chaque $U_{ij}$ soit envoyé
dans un ouvert affine de $S$. Les composés
$U_{ij} \to U_i \to X_i$ sont étales et quasi-compacts (voir
Espaces, lemme
\ref{spaces-lemma-composition-representable-transformations-property}).
Soit $X_{ij} \subset X_i$ le sous-espace ouvert correspondant à
l'image de $|U_{ij}| \to |X_i|$, voir le
lemme \ref{lemma-etale-image-open}.
Notons que $U_{ij} \to X_{ij}$ est quasi-compact, car $X_{ij} \subset X_i$
est un monomorphisme et que $U_{ij} \to X$ est quasi-compact.
Alors $X = \bigcup X_{ij}$ est un recouvrement comme en (4).
L'implication (4) $\Rightarrow$ (3) est immédiate.

\medskip\noindent
Supposons (4). Pour montrer que $X$ est localement quasi-séparé pour la topologie de Zariski sur $S$,
il suffit de montrer que $X_i$ est quasi-séparé sur $S$.
On peut donc supposer qu'il existe un schéma affine $U$ envoyé dans
un ouvert affine de $S$ et un morphisme étale, surjectif et quasi-compact
$U \to X$. Considérons le carré cartésien
$$
\xymatrix{
U \times_X U \ar[r] \ar[d] & U \times_S U \ar[d] \\
X \ar[r]^-{\Delta_{X/S}} & X \times_S X
}
$$
La flèche verticale droite est étale et surjective (voir
Espaces, lemme
\ref{spaces-lemma-product-representable-transformations-property})
et $U \times_S U$ est affine (car $U$ est envoyé dans un ouvert affine de $S$, voir
Schémas, section \ref{schemes-section-fibre-products}),
et $U \times_X U$ est quasi-compact,
car la projection $U \times_X U \to U$ est quasi-compacte en tant que
changement de base de $U \to X$. Il résulte de
Espaces, lemme \ref{spaces-lemma-representable-morphisms-spaces-property}
que $\Delta_{X/S}$ est quasi-compact, comme voulu.

\medskip\noindent
Supposons (1). Pour démontrer (3), on se ramène immédiatement au cas
où $X$ est quasi-séparé sur $S$. D'après le
lemme \ref{lemma-union-of-quasi-compact},
on peut trouver un recouvrement ouvert de Zariski $X = \bigcup X_i$ tel que chaque
$X_i$ soit envoyé dans un ouvert affine de $S$ et qu'il existe des schémas affines
$U_i$ et des morphismes étales surjectifs $U_i \to X_i$.
Puisque $U_i \to S$ se factorise par un ouvert affine de $S$, on voit que
$U_i \times_S U_i$ est affine, voir
Schémas, section \ref{schemes-section-fibre-products}.
Comme $X$ est quasi-séparé sur $S$, les morphismes
$$
R_i = U_i \times_{X_i} U_i = U_i \times_X U_i
\longrightarrow
U_i \times_S U_i
$$
sont quasi-compacts en tant que changements de base de $\Delta_{X/S}$. On en conclut
que $R_i$ est un schéma quasi-compact. Cela implique à son tour que chaque
projection $R_i \to U_i$ est quasi-compacte. Par conséquent, en appliquant
Espaces, lemme \ref{spaces-lemma-representable-morphisms-spaces-property}
au recouvrement $U_i \to X_i$ et au morphisme $U_i \to X_i$,
on conclut que les morphismes $U_i \to X_i$ sont quasi-compacts, comme voulu.

\medskip\noindent
On voit maintenant que (1), (3) et (4) sont équivalentes. Puisque (3) ne
fait pas intervenir le schéma de base, on en conclut qu'elles sont également équivalentes
à (2).
\end{proof}

\noindent
Le lemme suivant se révélera très utile.

\begin{lemma}
\label{lemma-finite-fibres-presentation}
Soit $S$ un schéma.
Soit $X$ un espace algébrique sur $S$.
Soit $U$ un schéma. Soit $\varphi : U \to X$ un morphisme étale tel que
les projections $R = U \times_X U \to U$ soient quasi-compactes ; c'est par exemple le cas si
$\varphi$ est quasi-compact. Alors les fibres de
$$
|U| \to |X|
\quad\text{et}\quad
|R| \to |X|
$$
sont finies.
\end{lemma}

\begin{proof}
Notons $R = U \times_X U$, et $s, t : R \to U$ les projections.
Soit $u \in U$ un point, et soit $x \in |X|$ son image.
La fibre de $|U| \to |X|$ au-dessus de $x$ est égale à
$s(t^{-1}(\{u\}))$ d'après le
lemme \ref{lemma-points-cartesian},
et la fibre de $|R| \to |X|$ au-dessus de $x$ est $t^{-1}(s(t^{-1}(\{u\})))$.
Puisque $t : R \to U$ est étale et quasi-compact, ses fibres sont finies
(ses fibres sont des sommes disjointes de spectres de corps d'après
Morphismes, lemme \ref{morphisms-lemma-etale-over-field},
et elles sont quasi-compactes). D'où le résultat.
\end{proof}








\section{Propriétés des espaces définies par des propriétés des schémas}
\sectionmark{Propriétés des espaces héritées des schémas}
\label{section-types-properties}

\noindent
Toute propriété de schémas qui est locale pour la topologie \'etale donne lieu à une
propriété correspondante des espaces algébriques grâce au lemme suivant.

\begin{lemma}
\label{lemma-type-property}
Soit $S$ un schéma.
Soit $X$ un espace algébrique sur $S$.
Soit $\mathcal{P}$ une propriété de schémas locale pour la topologie \'etale,
voir
Descente, définition \ref{descent-definition-property-local}.
Les assertions suivantes sont équivalentes :
\begin{enumerate}
\item il existe un schéma $U$ et un morphisme \'etale surjectif $U \to X$
tels que le schéma $U$ possède la propriété $\mathcal{P}$, et
\item pour tout schéma $U$ et tout morphisme \'etale $U \to X$,
le schéma $U$ possède la propriété $\mathcal{P}$.
\end{enumerate}
Si $X$ est représentable, ces conditions équivalent à $\mathcal{P}(X)$.
\end{lemma}

\begin{proof}
L'implication (2) $\Rightarrow$ (1) est immédiate.
Réciproquement, choisissons un morphisme \'etale surjectif $U \to X$
où $U$ est un schéma qui possède $\mathcal{P}$, et soit $V$ un schéma \'etale sur
$X$. Alors $U \times_X V \rightarrow V$ est un morphisme \'etale surjectif
de schémas ; ainsi, $V$ hérite de $\mathcal{P}$ de $U \times_X V$, lequel
hérite à son tour de $\mathcal{P}$ de $U$ (voir la discussion qui suit
Descente, définition \ref{descent-definition-property-local}).
La dernière assertion résulte clairement de (1) et de
Descente, définition \ref{descent-definition-property-local}.
\end{proof}

\begin{definition}
\label{definition-type-property}
Soit $\mathcal{P}$ une propriété de schémas qui est
locale pour la topologie \'etale.
Soit $S$ un schéma.
Soit $X$ un espace algébrique sur $S$.
On dit que $X$ {\it possède la propriété $\mathcal{P}$}
si l'une des conditions équivalentes du
lemme \ref{lemma-type-property}
est satisfaite.
\end{definition}

\begin{remark}
\label{remark-list-properties-local-etale-topology}
Voici une liste de propriétés qui sont locales pour la topologie \'etale
(rappelons que les topologies fpqc, fppf, syntomique et lisse sont
plus fines que la topologie \'etale) :
\begin{enumerate}
\item localement noethérien, voir
Descente, lemme \ref{descent-lemma-Noetherian-local-fppf},
\item de Jacobson, voir
Descente, lemme \ref{descent-lemma-Jacobson-local-fppf},
\item localement noethérien et vérifiant $(S_k)$, voir
Descente, lemme \ref{descent-lemma-Sk-local-syntomic},
\item Cohen-Macaulay, voir
Descente, lemme \ref{descent-lemma-CM-local-syntomic},
\item de Gorenstein, voir
Dualité pour les schémas, lemme \ref{duality-lemma-gorenstein-local-syntomic},
\item réduit, voir
Descente, lemme \ref{descent-lemma-reduced-local-smooth},
\item normal, voir
Descente, lemme \ref{descent-lemma-normal-local-smooth},
\item localement noethérien et vérifiant $(R_k)$, voir
Descente, lemme \ref{descent-lemma-Rk-local-smooth},
\item régulier, voir
Descente, lemme \ref{descent-lemma-regular-local-smooth},
\item de Nagata, voir
Descente, lemme \ref{descent-lemma-Nagata-local-smooth}.
\end{enumerate}
\end{remark}

\noindent
Toute propriété de germes de schémas locale pour la topologie \'etale donne lieu à une
propriété correspondante des espaces algébriques. Voici le lemme de rigueur.

\begin{lemma}
\label{lemma-local-source-target-at-point}
Soit $\mathcal{P}$ une propriété de germes de schémas locale pour la topologie \'etale, voir
Descente, définition \ref{descent-definition-local-at-point}.
Soit $S$ un schéma.
Soit $X$ un espace algébrique sur $S$.
Soit $x \in |X|$ un point de $X$.
Considérons les morphismes \'etales $a : U \to X$, où $U$ est un schéma.
Les assertions suivantes sont équivalentes :
\begin{enumerate}
\item pour tout $U \to X$ comme ci-dessus et tout $u \in U$ tel que $a(u) = x$, on a
$\mathcal{P}(U, u)$, et
\item il existe $U \to X$ comme ci-dessus et $u \in U$ tel que $a(u) = x$ et que l'on ait
$\mathcal{P}(U, u)$.
\end{enumerate}
Si $X$ est représentable, ces conditions équivalent à $\mathcal{P}(X, x)$.
\end{lemma}

\begin{proof}
Démonstration omise.
\end{proof}

\begin{definition}
\label{definition-property-at-point}
Soit $S$ un schéma. Soit $X$ un espace algébrique sur $S$.
Soit $x \in |X|$. Soit $\mathcal{P}$ une propriété de germes de schémas qui est
locale pour la topologie \'etale.
On dit que $X$ {\it possède la propriété $\mathcal{P}$ en $x$} si l'une des
conditions équivalentes du
lemme \ref{lemma-local-source-target-at-point}
est satisfaite.
\end{definition}

\begin{remark}
\label{remark-list-properties-local-ring-local-etale-topology}
Soit $P$ une propriété d'anneaux locaux. Supposons que, pour tout
homomorphisme \'etale d'anneaux $A \to B$ et tout idéal premier $\mathfrak q$ de $B$ au-dessus de
l'idéal premier $\mathfrak p$ de $A$, on ait
$P(A_\mathfrak p) \Leftrightarrow P(B_\mathfrak q)$.
On obtient alors une propriété locale pour la topologie \'etale des germes $(U, u)$ de schémas
en posant $\mathcal{P}(U, u) = P(\mathcal{O}_{U, u})$.
Dans cette situation, nous emploierons l'expression
``l'anneau local de $X$ en $x$ possède $P$'' pour signifier que
$X$ possède la propriété $\mathcal{P}$ en $x$.
Voici une liste de telles propriétés $P$ :
\begin{enumerate}
\item noethérien, voir
Compléments d'algèbre, lemme \ref{more-algebra-lemma-Noetherian-etale-extension},
\item de dimension $d$, voir
Compléments d'algèbre, lemme \ref{more-algebra-lemma-dimension-etale-extension},
\item régulier, voir
Compléments d'algèbre, lemme \ref{more-algebra-lemma-regular-etale-extension},
\item anneau de valuation discrète, cela résulte de (2), (3) et du
lemme \ref{algebra-lemma-characterize-dvr} d'Algèbre,
\item réduit, voir
Compléments d'algèbre, lemme \ref{more-algebra-lemma-henselization-reduced},
\item normal, voir
Compléments d'algèbre, lemme \ref{more-algebra-lemma-henselization-normal},
\item noethérien et de profondeur $k$, voir
Compléments d'algèbre, lemme \ref{more-algebra-lemma-henselization-depth},
\item noethérien et Cohen-Macaulay, voir
Compléments d'algèbre, lemme \ref{more-algebra-lemma-henselization-CM},
\item noethérien et de Gorenstein, voir
Complexes dualisants, lemme \ref{dualizing-lemma-flat-under-gorenstein}.
\end{enumerate}
Il existe d'autres propriétés pour lesquelles ceci vaut, par exemple celles d'être un anneau G ou un
anneau de Nagata. Si nous en avons un jour besoin, nous les ajouterons ici,
ainsi que des références à des démonstrations détaillées des faits
d'algèbre correspondants.
\end{remark}







\section{Ensembles constructibles}
\label{section-constructible}

\begin{lemma}
\label{lemma-locally-constructible}
Soit $S$ un schéma. Soit $X$ un espace algébrique sur $S$.
Soit $E \subset |X|$ une partie. Les assertions suivantes sont équivalentes :
\begin{enumerate}
\item pour tout morphisme \'etale $U \to X$, où $U$ est un
schéma, l'image réciproque de $E$ dans $U$ est une partie localement constructible
de $U$,
\item pour tout morphisme \'etale $U \to X$, où $U$ est un
schéma affine, l'image réciproque de $E$ dans $U$ est une partie constructible
de $U$,
\item il existe un morphisme \'etale surjectif $U \to X$, où $U$ est un
schéma, tel que l'image réciproque de $E$ dans $U$ soit une partie localement constructible
de $U$.
\end{enumerate}
\end{lemma}

\begin{proof}
D'après Propriétés, lemme \ref{properties-lemma-locally-constructible},
on voit que (1) et (2) sont équivalentes. Il est immédiat que (1)
implique (3). Supposons donc donné un morphisme \'etale surjectif
$\varphi : U \to X$, où $U$ est un schéma tel que $\varphi^{-1}(E)$
soit localement constructible. Soit $\varphi' : U' \to X$ un autre
morphisme \'etale, où $U'$ est un schéma. On a alors
$$
E'' = \text{pr}_1^{-1}(\varphi^{-1}(E)) = \text{pr}_2^{-1}((\varphi')^{-1}(E))
$$
où $\text{pr}_1 : U \times_X U' \to U$ et
$\text{pr}_2 : U \times_X U' \to U'$ sont les projections.
D'après Morphismes, lemme \ref{morphisms-lemma-inverse-image-constructible},
on voit que $E''$ est localement constructible dans $U \times_X U'$.
Soit $W' \subset U'$ un ouvert affine. Puisque $\text{pr}_2$ est
\'etale et donc ouverte, on peut choisir un ouvert quasi-compact
$W'' \subset U \times_X U'$ tel que $\text{pr}_2(W'') = W'$.
Alors $\text{pr}_2|_{W''} : W'' \to W'$ est quasi-compact.
On a $W' \cap (\varphi')^{-1}(E) = \text{pr}_2(E'' \cap W'')$
puisque $\varphi$ est surjectif, voir le lemme \ref{lemma-points-cartesian}.
Ainsi, $W' \cap (\varphi')^{-1}(E) = \text{pr}_2(E'' \cap W'')$
est localement constructible d'après
Morphismes, théorème \ref{morphisms-theorem-chevalley}, comme voulu.
\end{proof}

\begin{definition}
\label{definition-locally-constructible}
Soit $S$ un schéma. Soit $X$ un espace algébrique sur $S$.
Soit $E \subset |X|$ une partie. On dit que $E$ est
{\it localement constructible pour la topologie \'etale} si les conditions
équivalentes du lemme \ref{lemma-locally-constructible} sont satisfaites.
\end{definition}

\noindent
Bien entendu, si $X$ est représentable, c'est-à-dire si $X$ est un schéma,
cela signifie simplement que $E$ est une partie localement constructible
de l'espace topologique sous-jacent.







\section{Dimension en un point}
\label{section-dimension}

\noindent
On peut utiliser
Descente, lemme \ref{descent-lemma-dimension-at-point-local},
pour définir la dimension d'un espace
algébrique $X$ en un point $x$. On obtient ainsi une notion différente de la
notion topologique (c'est-à-dire la dimension de $|X|$ en $x$).

\begin{definition}
\label{definition-dimension-at-point}
Soit $S$ un schéma.
Soit $X$ un espace algébrique sur $S$.
Soit $x \in |X|$ un point de $X$.
On définit la {\it dimension de $X$ en $x$} comme
l'élément $\dim_x(X) \in \{0, 1, 2, \ldots, \infty\}$
tel que $\dim_x(X) = \dim_u(U)$ pour tout (ou, de manière équivalente, pour un)
couple $(a : U \to X, u)$ constitué d'un morphisme \'etale $a : U \to X$
d'un schéma vers $X$ et d'un point $u \in U$ tel que $a(u) = x$.
Voir la
définition \ref{definition-property-at-point}, le
lemme \ref{lemma-local-source-target-at-point} et
Descente, lemme \ref{descent-lemma-dimension-at-point-local}.
\end{definition}

\noindent
Attention : en général, on n'a {\bf pas} $\dim_x(X) = \dim_x(|X|)$.
Un contre-exemple est l'espace algébrique $X$ de
Espaces, exemple \ref{spaces-example-infinite-product}.
Plus précisément, soit $x \in |X|$ un point distinct du point générique $x_0$
de $|X|$. Alors $\dim_x(X) = 0$, mais $\dim_x(|X|) = 1$.
En particulier, la dimension de $X$ (telle qu'elle est définie
ci-dessous) est différente de la dimension de $|X|$.

\begin{definition}
\label{definition-dimension}
Soit $S$ un schéma. Soit $X$ un espace algébrique sur $S$.
La {\it dimension} $\dim(X)$ de $X$ est définie par la formule
$$
\dim(X) = \sup\nolimits_{x \in |X|} \dim_x(X)
$$
\end{definition}

\noindent
D'après
Propriétés, lemme \ref{properties-lemma-dimension},
on voit que c'est la notion usuelle lorsque $X$ est un schéma.
Il existe un autre invariant qui mesure la dimension d'un schéma
en un point, à savoir la dimension de l'anneau local. Cet invariant
est lui aussi compatible avec les morphismes \'etales, voir la
section \ref{section-dimension-local-ring}.






\section{Dimension des anneaux locaux}
\label{section-dimension-local-ring}

\noindent
La dimension de l'anneau local d'un espace algébrique est une notion bien
définie.

\begin{lemma}
\label{lemma-pre-dimension-local-ring}
Soit $S$ un schéma.
Soit $X$ un espace algébrique sur $S$.
Soit $x \in |X|$ un point.
Soit $d \in \{0, 1, 2, \ldots, \infty\}$.
Les assertions suivantes sont équivalentes :
\begin{enumerate}
\item il existe un schéma $U$, un morphisme \'etale $a : U \to X$ et un point
$u \in U$ tel que $a(u) = x$ et $\dim(\mathcal{O}_{U, u}) = d$,
\item pour tout schéma $U$, tout morphisme \'etale $a : U \to X$ et tout point
$u \in U$ tel que $a(u) = x$, on a $\dim(\mathcal{O}_{U, u}) = d$.
\end{enumerate}
Si $X$ est un schéma, ces conditions équivalent à $\dim(\mathcal{O}_{X, x}) = d$.
\end{lemma}

\begin{proof}
Combiner le
lemme \ref{lemma-local-source-target-at-point} et
Descente, lemme \ref{descent-lemma-dimension-local-ring-local}.
\end{proof}

\begin{definition}
\label{definition-dimension-local-ring}
Soit $S$ un schéma. Soit $X$ un espace algébrique sur $S$. Soit $x \in |X|$
un point. La {\it dimension de l'anneau local de $X$ en $x$} est
l'élément $d \in \{0, 1, 2, \ldots, \infty\}$ qui satisfait les conditions
équivalentes du lemme \ref{lemma-pre-dimension-local-ring}. Dans ce cas, on
dira aussi que {\it $x$ est un point de codimension $d$ dans $X$}.
\end{definition}

\noindent
Outre le lemme ci-dessous, signalons aussi au lecteur les
lemmes \ref{lemma-dimension-local-ring} et
\ref{lemma-dimension-decent-invariant-under-etale}.

\begin{lemma}
\label{lemma-dimension}
Soit $S$ un schéma. Soit $X$ un espace algébrique sur $S$.
Les quantités suivantes sont égales :
\begin{enumerate}
\item La dimension de $X$.
\item La borne supérieure des dimensions des anneaux locaux de $X$.
\item La borne supérieure des $\dim_x(X)$ pour $x \in |X|$.
\end{enumerate}
\end{lemma}

\begin{proof}
Les nombres de (1) et (3) sont égaux par la définition \ref{definition-dimension}.
Soit $U \to X$ un morphisme \'etale surjectif issu d'un schéma $U$.
La borne supérieure des $\dim_x(X)$ pour $x \in |X|$ est égale à la
borne supérieure des $\dim_u(U)$ pour les points $u$ de $U$ par définition.
Celle-ci est égale à la borne supérieure des $\dim(\mathcal{O}_{U, u})$ d'après
Propriétés, lemme \ref{properties-lemma-dimension}. Ce dernier
est à son tour égal à (2) par définition.
\end{proof}




\section{Points génériques}
\label{section-generic-points}

\noindent
Soit $T$ un espace topologique. D'après la deuxième édition de
l'EGA I, un {\it point maximal de $T$} est un point générique d'une composante
irréductible de $T$. Si $T = |X|$ est l'espace topologique associé à
un espace algébrique $X$, il existe au moins deux notions de points maximaux :
on peut considérer les points maximaux de $T$ vu comme espace topologique, ou
les images des points maximaux de $U$, où $U \to X$ est un
morphisme \'etale et $U$ un schéma. La seconde notion correspond à
l'ensemble des points de codimension $0$ (lemme \ref{lemma-codimension-0-points}).
Les points de codimension $0$ sont plus faciles
à manier pour les espaces algébriques généraux ; les deux notions
coïncident pour les espaces algébriques quasi-séparés et, plus généralement, décents
(Espaces algébriques décents, lemme \ref{decent-spaces-lemma-decent-generic-points}).

\begin{lemma}
\label{lemma-codimension-0-points}
Soit $S$ un schéma et soit $X$ un espace algébrique sur $S$.
Soit $x \in |X|$. Considérons les morphismes \'etales $a : U \to X$, où
$U$ est un schéma. Les assertions suivantes sont équivalentes :
\begin{enumerate}
\item $x$ est un point de codimension $0$ dans $X$,
\item il existe $U \to X$ comme ci-dessus et $u \in U$ tel que $a(u) = x$,
et le point $u$ est le point générique d'une composante irréductible de $U$,
\item pour tout $U \to X$ comme ci-dessus et tout $u \in U$ ayant pour image $x$,
le point $u$ est le point générique d'une composante irréductible de $U$.
\end{enumerate}
Si $X$ est représentable, ces conditions équivalent à ce que $x$ soit un point générique
d'une composante irréductible de $|X|$.
\end{lemma}

\begin{proof}
Remarquons qu'un point $u$ d'un schéma $U$ est le point générique d'une
composante irréductible de $U$ si et seulement si $\dim(\mathcal{O}_{U, u}) = 0$
(Propriétés, lemme \ref{properties-lemma-generic-point}).
Le résultat découle donc de la définition de la codimension d'un
point dans $X$ (définition \ref{definition-dimension-local-ring}).
\end{proof}

\begin{lemma}
\label{lemma-codimension-0-points-dense}
Soit $S$ un schéma et soit $X$ un espace algébrique sur $S$.
L'ensemble des points de codimension $0$ de $X$ est dense dans $|X|$.
\end{lemma}

\begin{proof}
Si $U$ est un schéma, l'ensemble des points génériques de ses composantes
irréductibles est dense dans $U$ (cela vaut pour tout espace topologique quasi-sobre).
Ainsi, si $U \to X$ est un morphisme \'etale surjectif, l'ensemble
des points de codimension $0$ de $X$ est l'image d'une partie dense de
$|U|$ (lemme \ref{lemma-codimension-0-points}). Puisque $|X|$ possède la
topologie quotient induite par $|U| \to |X|$, on conclut.
\end{proof}





\section{Espaces réduits}
\label{section-reduced}

\noindent
Nous avons déjà défini les espaces algébriques réduits à la
section \ref{section-types-properties}.
Nous démontrons ici seulement quelques lemmes simples sur les espaces
algébriques réduits.

\begin{lemma}
\label{lemma-subspace-induced-topology}
Soit $S$ un schéma. Soit $Z \to X$ une immersion d'espaces algébriques.
Alors $|Z| \to |X|$ est un homéomorphisme de $|Z|$ sur une partie localement fermée
de $|X|$.
\end{lemma}

\begin{proof}
Soit $U$ un schéma et soit $U \to X$ un morphisme \'etale surjectif.
Alors $Z \times_X U \to U$ est une immersion de schémas, donc induit un
homéomorphisme de $|Z \times_X U|$ sur une partie localement fermée $T'$
de $|U|$. D'après le lemme \ref{lemma-points-cartesian}, la partie
$T'$ est l'image réciproque de l'image $T$ de $|Z| \to |X|$.
L'application $|Z| \to |X|$ est injective, car la transformation de
foncteurs $Z \to X$ est injective, voir
Espaces, section \ref{spaces-section-Zariski}. D'après
Topologie, lemme \ref{topology-lemma-open-morphism-quotient-topology}
on voit que $T$ est localement fermée dans $|X|$. De plus, l'application continue
$|Z| \to T$ est un homéomorphisme, puisque l'application $|Z \times_X U| \to T'$
est un homéomorphisme et que $|Z \times_X U| \to |Z|$ est submersive.
\end{proof}

\noindent
Le lemme suivant nous aidera à construire des sous-espaces
(localement) fermés.

\begin{lemma}
\label{lemma-subspaces-presentation}
Soit $S$ un schéma. Soit $j : R \to U \times_S U$ une relation d'équivalence \'etale.
Soit $X = U/R$ l'espace algébrique associé
(Espaces, théorème \ref{spaces-theorem-presentation}). Il existe une
bijection canonique
$$
\begin{gathered}
\text{Sous-schémas localement fermés stables par }R : Z'\text{ de }U\\[2pt]
\leftrightarrow\\[2pt]
\text{sous-espaces localement fermés }Z\text{ de }X
\end{gathered}
$$
De plus, $Z \to X$ est fermé (resp.\ ouvert) si et seulement si
$Z' \to U$ est fermé (resp.\ ouvert).
\end{lemma}

\begin{proof}
Notons $\varphi : U \to X$ l'application canonique. La bijection envoie
$Z \to X$ sur $Z' = Z \times_X U \to U$. Il résulte immédiatement de la définition
que $Z' \to U$ est une immersion, resp.\ une immersion fermée, resp.\ une
immersion ouverte, si tel est le cas de $Z \to X$. Il est aussi clair que $Z'$ est stable par $R$
(voir Groupoïdes, définition \ref{groupoids-definition-invariant-open}).

\medskip\noindent
Réciproquement, supposons que $Z' \to U$ soit une immersion stable par $R$.
Soit $R'$ la restriction de $R$ à $Z'$, voir
Groupoïdes, définition \ref{groupoids-definition-restrict-groupoid}.
Puisque $R' = R \times_{s, U} Z' = Z' \times_{U, t} R$ dans ce cas,
on voit que $R'$ est une relation d'équivalence \'etale sur $Z'$. D'après
Espaces, théorème \ref{spaces-theorem-presentation}, on voit que
$Z = Z'/R'$ est un espace algébrique. Par construction, on a
$U \times_X Z = Z'$ ; ainsi, $U \times_X Z \to U$ est une immersion.
Notons que la propriété ``immersion'' est préservée par changement de base
et locale pour la topologie fppf sur la base (voir Espaces, section \ref{spaces-section-lists}).
De plus, les immersions sont séparées et localement quasi-finies (voir
Schémas, lemme \ref{schemes-lemma-immersions-monomorphisms}
et
Morphismes, lemme \ref{morphisms-lemma-immersion-locally-quasi-finite}).
Par conséquent, d'après Compléments sur les morphismes, lemme
\ref{more-morphisms-lemma-separated-locally-quasi-finite-morphisms-fppf-descend}
les immersions satisfont à la descente pour les recouvrements fppf. Toutes les hypothèses de
Espaces,
lemme \ref{spaces-lemma-morphism-sheaves-with-P-effective-descent-etale}
sont donc satisfaites pour $Z \to X$, $\mathcal{P}=$``immersion'',
et le morphisme \'etale surjectif $U \to X$. On conclut que $Z \to X$
est représentable et est une immersion, ce qui est la
définition d'un sous-espace (voir
Espaces, définition \ref{spaces-definition-immersion}).

\medskip\noindent
Il est clair que ces constructions sont inverses l'une de l'autre, d'où le résultat.
\end{proof}

\begin{lemma}
\label{lemma-reduced-closed-subspace}
Soit $S$ un schéma.
Soit $X$ un espace algébrique sur $S$.
Soit $T \subset |X|$ une partie fermée.
Il existe un unique sous-espace fermé $Z \subset X$ possédant
les propriétés suivantes : (a) on a $|Z| = T$, et (b) $Z$ est réduit.
\end{lemma}

\begin{proof}
Soit $U \to X$ un morphisme \'etale surjectif, où $U$ est un schéma.
Posons $R = U \times_X U$, de sorte que $X = U/R$, voir
Espaces, lemme \ref{spaces-lemma-space-presentation}.
Comme d'habitude, on note $s, t : R \to U$ les deux morphismes de projection.
D'après le lemme \ref{lemma-points-presentation},
on voit que $T$ correspond à une partie fermée $T' \subset |U|$ telle
que $s^{-1}(T') = t^{-1}(T')$.
Soit $Z' \subset U$ le sous-schéma porté par $T'$ et muni de la structure réduite induite.
Dans ce cas, les produits fibrés
$Z' \times_{U, t} R$ et $Z' \times_{U, s} R$ sont des sous-schémas fermés
de $R$
(Schémas, lemme \ref{schemes-lemma-base-change-immersion})
qui sont \'etales sur $Z'$
(Morphismes, lemme \ref{morphisms-lemma-base-change-etale}),
et sont donc réduits
(car la propriété d'être réduit est locale pour la topologie \'etale, voir la
remarque \ref{remark-list-properties-local-etale-topology}).
Comme ils ont le même espace topologique sous-jacent (voir ci-dessus),
on conclut que $Z' \times_{U, t} R = Z' \times_{U, s} R$.
On peut donc appliquer le lemme \ref{lemma-subspaces-presentation}
pour obtenir un sous-espace fermé $Z \subset X$ dont l'image réciproque sur $U$ est $Z'$.
Par construction, $|Z| = T$ et $Z$ est réduit. Cela prouve l'existence.
Nous omettons la démonstration de l'unicité.
\end{proof}

\begin{lemma}
\label{lemma-map-into-reduction}
Soit $S$ un schéma.
Soient $X$, $Y$ des espaces algébriques sur $S$.
Soit $Z \subset X$ un sous-espace fermé.
Supposons $Y$ réduit.
Un morphisme $f : Y \to X$ se factorise par $Z$ si et seulement si
$f(|Y|) \subset |Z|$.
\end{lemma}

\begin{proof}
Supposons $f(|Y|) \subset |Z|$. Choisissons un diagramme
$$
\xymatrix{
V \ar[d]_b \ar[r]_h & U \ar[d]^a \\
Y \ar[r]^f & X
}
$$
où $U$, $V$ sont des schémas et les flèches verticales sont surjectives et
\'etales. Le schéma $V$ est réduit, voir le
lemme \ref{lemma-type-property}.
Par conséquent, $h$ se factorise par $a^{-1}(Z)$ d'après
Schémas, lemme \ref{schemes-lemma-map-into-reduction}.
Ainsi, $a \circ h$ se factorise par $Z$.
Puisque $Z \subset X$ est un sous-faisceau et que $V \to Y$ est une surjection de faisceaux
sur $(\Sch/S)_{fppf}$, on conclut que $Y \to X$ se factorise
par $Z$.
\end{proof}

\begin{definition}
\label{definition-reduced-induced-space}
Soit $S$ un schéma et soit $X$ un espace algébrique sur $S$.
Soit $Z \subset |X|$ une partie fermée.
Une {\it structure d'espace algébrique sur $Z$} est donnée par un sous-espace fermé
$Z'$ de $X$ tel que $|Z'|$ soit égal à $Z$.
La {\it structure d'espace algébrique réduite induite}
sur $Z$ est celle qui a été construite dans le
lemme \ref{lemma-reduced-closed-subspace}.
La {\it réduction $X_{red}$ de $X$} est la structure d'espace
algébrique réduite induite sur $|X|$.
\end{definition}











\section{Le lieu schématique}
\label{section-schematic}

\noindent
Tout espace algébrique possède un plus grand sous-espace ouvert qui est un
schéma ; cela est plus ou moins clair, mais nous en donnons aussi la preuve ci-dessous.
Bien entendu, ce sous-espace peut être vide, par exemple si
$X = \mathbf{A}^1_{\mathbf{Q}}/\mathbf{Z}$ (le
contre-exemple universel). En revanche, si $X$ est par exemple quasi-séparé,
alors ce plus grand sous-espace ouvert est en fait dense dans $X$ !

\begin{lemma}
\label{lemma-subscheme}
Soit $S$ un schéma.
Soit $X$ un espace algébrique sur $S$.
Il existe un plus grand sous-espace ouvert $X' \subset X$ qui est un schéma.
\end{lemma}

\begin{proof}
Soit $U \to X$ un morphisme \'etale surjectif, où $U$ est un schéma.
Soit $R = U \times_X U$. Les sous-espaces ouverts de $X$ correspondent biunivoquement
avec les sous-schémas ouverts de $U$ qui sont stables par $R$. Il existe donc un
ensemble de tels sous-espaces. Soient $X_i$, $i \in I$, les sous-espaces ouverts
de $X$ qui sont des schémas, c'est-à-dire qui sont représentables. Considérons le
sous-espace ouvert $X' \subset X$ dont l'ensemble de points sous-jacent est
l'ouvert $\bigcup |X_i|$ de $|X|$. D'après le
lemme \ref{lemma-characterize-surjective},
on voit que
$$
\coprod X_i \longrightarrow X'
$$
est une application surjective de faisceaux sur $(\Sch/S)_{fppf}$.
Mais puisque chaque $X_i \to X'$ est représentable par des immersions ouvertes,
on voit qu'en fait l'application est surjective pour la topologie de Zariski.
En effet, si $T \to X'$ est un morphisme d'un schéma
vers $X'$, alors $X_i \times_{X'} T$ est un sous-schéma ouvert de $T$.
On peut donc appliquer
Schémas, lemme \ref{schemes-lemma-glue-functors}
pour voir que $X'$ est un schéma.
\end{proof}

\noindent
Dans la suite de cette section, on dit qu'un sous-espace ouvert
$X'$ d'un espace algébrique $X$ est {\it dense} si la partie ouverte
correspondante $|X'| \subset |X|$ est dense.

\begin{lemma}
\label{lemma-quasi-separated-finite-etale-cover-dense-open-scheme}
Soit $S$ un schéma. Soit $X$ un espace algébrique sur $S$.
S'il existe un morphisme fini, \'etale et surjectif
$U \to X$ où $U$ est un schéma quasi-séparé, alors
il existe un sous-espace ouvert dense $X'$ de $X$ qui est un schéma.
Plus précisément, tout point $x \in |X|$ de codimension $0$ dans $X$
appartient à $X'$.
\end{lemma}

\begin{proof}
Soit $X' \subset X$ le sous-espace ouvert maximal qui est un schéma
(lemme \ref{lemma-subscheme}).
Soit $x \in |X|$ un point de codimension $0$ dans $X$.
D'après le lemme \ref{lemma-codimension-0-points-dense},
il suffit de montrer que $x \in X'$.
Soit $U \to X$ comme dans l'énoncé du lemme.
Écrivons $R = U \times_X U$ et notons $s, t : R \to U$ les projections usuelles.
Notons que $s, t$ sont surjectifs, finis et \'etales.
D'après le lemme \ref{lemma-finite-fibres-presentation},
la fibre de $|U| \to |X|$ au-dessus de $x$ est finie, disons
$\{\eta_1, \ldots, \eta_n\}$. D'après le lemme \ref{lemma-codimension-0-points},
chaque $\eta_i$ est le point générique d'une composante irréductible de $U$.
D'après Propriétés, lemme \ref{properties-lemma-maximal-points-affine},
on peut trouver un ouvert affine $W \subset U$ contenant
$\{\eta_1, \ldots, \eta_n\}$
(c'est ici que l'on utilise le fait que $U$ est quasi-séparé). D'après
Groupoïdes, lemme \ref{groupoids-lemma-find-invariant-affine}
on peut supposer que $W$ est stable par $R$.
Puisque $W \subset U$ est un ouvert affine stable par $R$, la restriction
$R_W$ de $R$ à $W$ est égale à $R_W = s^{-1}(W) = t^{-1}(W)$ (voir
Groupoïdes, définition \ref{groupoids-definition-invariant-open}
et la discussion qui la suit). En particulier, les morphismes $R_W \to W$ sont
eux aussi finis et \'etales. Il en résulte que $R_W$ est affine.
On voit donc que $W/R_W$ est un schéma, d'après
Groupoïdes, proposition \ref{groupoids-proposition-finite-flat-equivalence}.
D'autre part, $W/R_W$ est un sous-espace ouvert de $X$ d'après
Espaces, lemme \ref{spaces-lemma-finding-opens}, et il contient
$x$ par construction.
\end{proof}

\noindent
Nous étendrons la proposition suivante au cas des
espaces algébriques décents dans
Espaces algébriques décents, théorème
\ref{decent-spaces-theorem-decent-open-dense-scheme}.

\begin{proposition}
\label{proposition-locally-quasi-separated-open-dense-scheme}
Soit $S$ un schéma. Soit $X$ un espace algébrique sur $S$. Si $X$ est
localement quasi-séparé pour la topologie de Zariski (par exemple si $X$ est quasi-séparé), alors
il existe un sous-espace ouvert dense $X'$ de $X$ qui est un schéma.
Plus précisément, tout point $x \in |X|$ de codimension $0$ dans $X$
appartient à $X'$.
\end{proposition}

\begin{proof}
La question est locale sur $X$ d'après le lemme \ref{lemma-subscheme}.
Ainsi, d'après le lemme \ref{lemma-quasi-separated-quasi-compact-pieces},
on peut supposer qu'il existe un schéma affine $U$ et un morphisme
surjectif, quasi-compact et \'etale $U \to X$.
De plus, $U \to X$ est séparé (lemme \ref{lemma-separated-cover}).
Posons $R = U \times_X U$ et notons $s, t : R \to U$ les projections
usuelles. Alors $s, t$ sont surjectifs, quasi-compacts, séparés et
\'etales. Par conséquent, $s, t$ sont aussi quasi-finis et ont des fibres finies
(Morphismes,
lemmes \ref{morphisms-lemma-etale-locally-quasi-finite},
\ref{morphisms-lemma-quasi-finite-locally-quasi-compact}, et
\ref{morphisms-lemma-quasi-finite}).
D'après Morphismes, lemme \ref{morphisms-lemma-generically-finite},
pour tout $\eta \in U$ qui est le point générique d'une
composante irréductible de $U$, il existe un voisinage ouvert
$V \subset U$ de $\eta$ tel que $s^{-1}(V) \to V$ soit fini. D'après
Descente, lemme \ref{descent-lemma-descending-property-finite},
la propriété d'être fini est locale pour la topologie fpqc (et en particulier \'etale) sur le but.
On peut donc appliquer
Compléments sur les groupoïdes, lemme \ref{more-groupoids-lemma-property-invariant},
qui affirme que le plus grand ouvert $W \subset U$ au-dessus duquel $s$ est
fini est stable par $R$. D'après ce qui précède, $W$ contient tout
point générique d'une composante irréductible de $U$.
La restriction $R_W$ de $R$ à $W$ est égale à $R_W = s^{-1}(W) = t^{-1}(W)$
(voir Groupoïdes, définition \ref{groupoids-definition-invariant-open}
et la discussion qui la suit).
Par construction, $s_W, t_W : R_W \to W$ sont finis \'etales.
Considérons le sous-espace ouvert $X' = W/R_W \subset X$ (voir
Espaces, lemme \ref{spaces-lemma-finding-opens}).
Par construction, l'application d'inclusion $X' \to X$
induit une bijection sur les points de codimension $0$.
On est ainsi ramené au
lemme \ref{lemma-quasi-separated-finite-etale-cover-dense-open-scheme}.
\end{proof}




\section{Obtention d'un schéma}
\label{section-getting-a-scheme}

\noindent
Nous avons utilisé dans la section précédente le fait que le quotient $U/R$ d'un
schéma affine $U$ par une relation d'équivalence $R$ est un schéma si les
morphismes $s, t : R \to U$ sont finis \'etales. C'est un cas particulier
du résultat suivant.

\begin{proposition}
\label{proposition-finite-flat-equivalence-global}
Soit $S$ un schéma.
Soit $(U, R, s, t, c)$ un groupoïde en schémas sur $S$.
Supposons que
\begin{enumerate}
\item $s, t : R \to U$ soient finis localement libres,
\item $j = (t, s)$ soit une relation d'équivalence, et
\item toute partie fermée non vide $Z$ de $U$ contienne un point
$u$ dont la classe d'équivalence pour $R$, $t(s^{-1}(\{u\}))$, soit contenue
dans un ouvert affine de $U$\footnote{Soit $E \subset U$ la
partie des points $u$ tels que $t(s^{-1}(\{u\}))$
soit contenu dans un ouvert affine de $U$. La condition (3) est satisfaite
si $E = U$, ou si tout point de type fini de $U$ appartient à $E$, ou
si tout $u \in U$ se spécialise en un point de $E$.}.
\end{enumerate}
Alors il existe un morphisme fini localement libre $U \to M$
de schémas sur $S$ tel que $R = U \times_M U$ et tel que $M$
représente le faisceau quotient $U/R$ pour la topologie fppf.
\end{proposition}

\begin{proof}
Par l'hypothèse (3) et
Groupoïdes, lemme \ref{groupoids-lemma-find-invariant-affine}
nous pouvons trouver un recouvrement ouvert $U = \bigcup U_i$ tel que chaque $U_i$
soit stable par $R$ et ouvert affine dans $U$. Posons $R_i = R|_{U_i}$.
Considérons les faisceaux fppf $F = U/R$ et $F_i = U_i/R_i$.
D'après Espaces, lemme \ref{spaces-lemma-finding-opens}, les morphismes
$F_i \to F$ sont représentables et sont des immersions ouvertes.
D'après Groupoïdes, proposition \ref{groupoids-proposition-finite-flat-equivalence},
les faisceaux $F_i$ sont représentables par des schémas affines.
Si $T$ est un schéma et $T \to F$ un morphisme, alors $V_i = F_i \times_F T$
est ouvert dans $T$, et nous affirmons que $T = \bigcup V_i$. En effet,
localement pour la topologie fppf sur $T$, nous pouvons relever $T \to F$ en un morphisme
$f : T \to U$ et, dans ce cas, $f^{-1}(U_i) \subset V_i$.
Nous en concluons que $F$ est représentable par un schéma, voir
Schémas, lemme \ref{schemes-lemma-glue-functors}.
\end{proof}

\noindent
Par exemple, si $U$ est isomorphe à un sous-schéma localement fermé d'un
schéma affine ou à un sous-schéma localement fermé de
$\text{Proj}(A)$ pour un anneau gradué $A$, alors la troisième hypothèse est satisfaite d'après
Propriétés, lemme \ref{properties-lemma-ample-finite-set-in-affine}.
En particulier, nous pouvons appliquer ce résultat aux actions libres de groupes finis et de
schémas en groupes finis sur des schémas quasi-affines ou quasi-projectifs.
Par exemple, le quotient $X/G$ d'une variété quasi-projective
$X$ par l'action libre d'un groupe fini $G$ est un schéma. Voici un
énoncé détaillé.

\begin{lemma}
\label{lemma-quotient-scheme}
Soit $S$ un schéma. Soit $G \to S$ un schéma en groupes. Soit $X \to S$ un
morphisme de schémas. Soit $a : G \times_S X \to X$ une action. Supposons que
\begin{enumerate}
\item $G \to S$ soit fini localement libre,
\item l'action $a$ soit libre,
\item $X \to S$ soit affine, ou quasi-affine, ou projectif, ou
quasi-projectif, ou que $X$ soit isomorphe à un sous-schéma ouvert d'un
schéma affine, ou que $X$ soit isomorphe à un sous-schéma ouvert de $\text{Proj}(A)$
pour un anneau gradué $A$, ou que $G \to S$ soit radiciel.
\end{enumerate}
Alors le faisceau quotient $X/G$ pour la topologie fppf est un schéma et $X \to X/G$
est un torseur sous $G$ pour la topologie fppf.
\end{lemma}

\begin{proof}
Montrons d'abord que $X/G$ est un schéma. Puisque l'action est libre, le morphisme
$j = (a, \text{pr}) : G \times_S X \to X \times_S X$
est un monomorphisme et donc une relation d'équivalence, voir
Groupoïdes, lemme \ref{groupoids-lemma-free-action}. Les morphismes
$s, t : G \times_S X \to X$
sont finis localement libres, puisque nous avons supposé que $G \to S$ est fini localement
libre. Pour conclure, il suffit maintenant de vérifier la dernière hypothèse de la
proposition \ref{proposition-finite-flat-equivalence-global}.
Puisque l'action de $G$ est définie sur $S$, il suffit de montrer que
toute partie finie d'une fibre de $X \to S$ est contenue dans un
ouvert affine de $X$. Si $X$ est isomorphe à un sous-schéma ouvert d'un
schéma affine ou à un sous-schéma ouvert de $\text{Proj}(A)$
pour un anneau gradué $A$, cela résulte de
Propriétés, lemme \ref{properties-lemma-ample-finite-set-in-affine}.
Si $X \to S$ est affine, ou quasi-affine, ou projectif, ou
quasi-projectif, nous pouvons remplacer $S$ par un ouvert affine et nous
nous ramenons au cas que nous venons de traiter. Si $G \to S$ est radiciel,
alors les orbites des points de $X$ sous l'action de $G$ sont réduites à un point
et la condition est trivialement satisfaite. Nous omettons quelques détails.

\medskip\noindent
Pour voir que $X \to X/G$ est un torseur sous $G$ pour la topologie fppf
(Groupoïdes, définition \ref{groupoids-definition-principal-homogeneous-space})
il faut montrer que $G \times_S X \to X \times_{X/G} X$
est un isomorphisme et que $X \to X/G$ admet localement des sections pour la topologie fppf.
La seconde assertion résulte du fait que $X \to X/G$ est surjectif
comme morphisme de faisceaux fppf (par construction). La première résulte de
l'isomorphisme $R = U \times_M U$ dans la conclusion de la
proposition \ref{proposition-finite-flat-equivalence-global}
(notons que $R = G \times_S X$ dans notre cas).
\end{proof}

\begin{lemma}
\label{lemma-quotient-separated}
Notations et hypothèses comme dans la
proposition \ref{proposition-finite-flat-equivalence-global}. Alors
\begin{enumerate}
\item si $U$ est quasi-séparé sur $S$, alors $U/R$ est quasi-séparé
sur $S$,
\item si $U$ est quasi-séparé, alors $U/R$ est quasi-séparé,
\item si $U$ est séparé sur $S$, alors $U/R$ est séparé sur $S$,
\item si $U$ est séparé, alors $U/R$ est séparé, et
\item d'autres énoncés pourront être ajoutés ici.
\end{enumerate}
Des résultats analogues valent dans la situation du Lemme \ref{lemma-quotient-scheme}.
\end{lemma}

\begin{proof}
Puisque $M$ représente le faisceau quotient, nous avons un diagramme cartésien
$$
\xymatrix{
R \ar[r]_-j \ar[d] & U \times_S U \ar[d] \\
M \ar[r] & M \times_S M
}
$$
de schémas. Puisque $U \times_S U \to M \times_S M$ est surjectif, fini et localement
libre, pour montrer que $M \to M \times_S M$ est quasi-compact, resp.\ une
immersion fermée, il suffit de montrer que $j : R \to U \times_S U$ est
quasi-compact, resp.\ une immersion fermée, voir
Descente, lemmes \ref{descent-lemma-descending-property-quasi-compact} et
\ref{descent-lemma-descending-property-closed-immersion}.
Puisque $j : R \to U \times_S U$ est un morphisme sur $U$ et que
$R$ est fini sur $U$, nous voyons que $j$ est quasi-compact dès que
le morphisme de projection $U \times_S U \to U$ est quasi-séparé
(Schémas, lemme \ref{schemes-lemma-quasi-compact-permanence}).
Puisque $j$ est un monomorphisme localement de type fini, nous voyons que
$j$ est une immersion fermée dès qu'il est propre
(Morphismes \'etales, lemme \ref{etale-lemma-characterize-closed-immersion})
ce qui est le cas dès que le morphisme de projection
$U \times_S U \to U$ est séparé
(Morphismes, lemme \ref{morphisms-lemma-image-proper-scheme-closed}).
Ceci prouve (1) et (3). Pour prouver (2) et (4), nous remplaçons $S$ par
$\Spec(\mathbf{Z})$, voir définition \ref{definition-separated}.
Puisque le Lemme \ref{lemma-quotient-scheme} se démontre en appliquant la
proposition \ref{proposition-finite-flat-equivalence-global},
le dernier énoncé est lui aussi clair.
\end{proof}




\section{Points des espaces quasi-séparés}
\label{section-points-quasi-separated}

\noindent
Les points peuvent se comporter très mal pour les espaces algébriques, au degré de généralité adopté
dans le Projet Champs. Cependant, pour les espaces quasi-séparés, leur comportement
est essentiellement semblable à celui des points des schémas. Nous démontrons quelques résultats
à ce sujet dans cette section ; le chapitre sur les espaces décents en contient bien davantage,
voir par exemple
Espaces algébriques décents, section \ref{decent-spaces-section-points}.

\begin{lemma}
\label{lemma-quasi-separated-sober}
Soit $S$ un schéma. Soit $X$ un espace algébrique qui soit localement
quasi-séparé pour la topologie de Zariski et défini sur $S$. Alors l'espace topologique $|X|$ est sobre (voir
Topologie, définition \ref{topology-definition-generic-point}).
\end{lemma}

\begin{proof}
En combinant
Topologie, lemme \ref{topology-lemma-sober-local}
et le
lemme \ref{lemma-quasi-separated-quasi-compact-pieces}
on voit que l'on peut supposer qu'il existe un schéma affine $U$
et un morphisme \'etale, surjectif et quasi-compact $U \to X$.
Posons $R = U \times_X U$, avec pour projections $s, t : R \to U$. En appliquant le
lemme \ref{lemma-finite-fibres-presentation}
on voit que les fibres de $s, t$ sont finies. Il s'ensuit que toutes les hypothèses de
Topologie, lemme \ref{topology-lemma-quotient-kolmogorov}
sont satisfaites, et l'on conclut que $|X|$ est un espace de Kolmogoroff\footnote{
En fait, nous utilisons aussi ici le
Schémas, lemme \ref{schemes-lemma-scheme-sober} (sobriété des schémas),
Morphismes, les lemmes \ref{morphisms-lemma-etale-flat}
et \ref{morphisms-lemma-generalizations-lift-flat} (les générisations
se relèvent le long des morphismes \'etales), le
lemme \ref{lemma-points-presentation} (les points d'un espace algébrique en
termes d'une présentation), et le
lemme \ref{lemma-topology-points} (l'application quotient est ouverte).}.

\medskip\noindent
Il reste à montrer que toute partie fermée irréductible
$T \subset |X|$ possède un point générique. D'après le
lemme \ref{lemma-reduced-closed-subspace}
il existe un sous-espace fermé $Z \subset X$ tel que $|Z| = T$.
Notons que $U \times_X Z \to Z$ est un morphisme \'etale, surjectif et
quasi-compact d'un schéma affine vers $Z$ ; ainsi, $Z$ est localement
quasi-séparé pour la topologie de Zariski d'après le
lemme \ref{lemma-quasi-separated-quasi-compact-pieces}.
Par la
proposition \ref{proposition-locally-quasi-separated-open-dense-scheme}
nous voyons qu'il existe un sous-espace ouvert dense $Z' \subset Z$
qui est un schéma. Cela signifie que $|Z'| \subset T$ est ouvert et dense.
L'espace topologique $|Z'|$ est donc irréductible, ce qui signifie que
$Z'$ est un schéma irréductible. D'après
Schémas, lemme \ref{schemes-lemma-scheme-sober}
nous concluons que $|Z'|$ est l'adhérence d'un unique point
$\eta \in |Z'| \subset T$ et donc aussi que $T = \overline{\{\eta\}}$, ce qui conclut.
\end{proof}

\begin{lemma}
\label{lemma-quasi-compact-quasi-separated-spectral}
Soit $S$ un schéma. Soit $X$ un espace algébrique quasi-compact et
quasi-séparé sur $S$. L'espace topologique $|X|$ est un espace spectral.
\end{lemma}

\begin{proof}
D'après Topologie, définition \ref{topology-definition-spectral-space},
nous devons vérifier que $|X|$ est sobre et quasi-compact, qu'il possède une base
d'ouverts quasi-compacts et que l'intersection de deux
ouverts quasi-compacts quelconques est quasi-compacte. D'après le
Lemme \ref{lemma-quasi-separated-sober}, nous voyons que $|X|$ est sobre.
D'après le Lemme \ref{lemma-quasi-compact-space}, nous voyons que $|X|$ est quasi-compact.
D'après le Lemme \ref{lemma-quasi-compact-affine-cover}, il existe un schéma affine
$U$ et un morphisme \'etale surjectif $f : U \to X$.
Puisque $|f| : |U| \to |X|$ est ouverte et continue, et que $|U|$ possède
une base d'ouverts quasi-compacts, nous concluons que $|X|$ possède une base
d'ouverts quasi-compacts. Enfin, supposons que
$A, B \subset |X|$ soient des ouverts quasi-compacts. Alors $A = |X'|$ et $B = |X''|$
pour des sous-espaces ouverts $X', X'' \subset X$ (Lemme \ref{lemma-open-subspaces}),
et nous pouvons choisir des schémas affines $V$ et $W$ et des morphismes
\'etales surjectifs $V \to X'$ et $W \to X''$
(Lemme \ref{lemma-quasi-compact-affine-cover}).
Alors $A \cap B$ est l'image de
$|V \times_X W| \to |X|$ (Lemme \ref{lemma-points-cartesian}).
Puisque $V \times_X W$ est quasi-compact car $X$ est quasi-séparé
(Lemme \ref{lemma-characterize-quasi-separated})
nous concluons que $A \cap B$ est quasi-compacte, ce qui achève la démonstration.
\end{proof}

\noindent
Le lemme suivant permet de prouver qu'un
espace algébrique est isomorphe au spectre d'un corps.

\begin{lemma}
\label{lemma-point-like-spaces}
Soit $S$ un schéma. Soit $k$ un corps.
Soit $X$ un espace algébrique sur $S$ et supposons qu'il existe
un morphisme \'etale surjectif $\Spec(k) \to X$.
Si $X$ est quasi-séparé, alors $X \cong \Spec(k')$,
où $k/k'$ est une extension finie séparable.
\end{lemma}

\begin{proof}
Posons $R = \Spec(k) \times_X \Spec(k)$ ; nous avons donc un
diagramme de produit fibré
$$
\xymatrix{
R \ar[r]_-s \ar[d]_-t & \Spec(k) \ar[d] \\
\Spec(k) \ar[r] & X
}
$$
D'après
Espaces, lemme \ref{spaces-lemma-space-presentation}
nous savons que $X = \Spec(k)/R$ est le faisceau quotient.
Comme $\Spec(k) \to X$ est \'etale, les morphismes $s$ et $t$ sont
\'etales. Ainsi, $R = \coprod_{i \in I} \Spec(k_i)$ est une somme
disjointe de spectres de corps, et $s$ et $t$
induisent tous deux des extensions finies séparables $s, t : k \subset k_i$, voir
Morphismes, lemme \ref{morphisms-lemma-etale-over-field}.
Comme
$$
R = \Spec(k) \times_X \Spec(k)
= (\Spec(k) \times_S \Spec(k)) \times_{X \times_S X, \Delta} X
$$
et que $\Delta$ est quasi-compact par hypothèse, nous concluons que
$R \to \Spec(k) \times_S \Spec(k)$ est quasi-compact.
Ainsi, $R$ est quasi-compact puisque $\Spec(k) \times_S \Spec(k)$ est
affine. Nous concluons que $I$ est fini. Il s'ensuit
que $s$ et $t$ sont des morphismes finis localement libres. D'après
Groupoïdes, proposition \ref{groupoids-proposition-finite-flat-equivalence}
nous concluons que $\Spec(k)/R$ est
représenté par $\Spec(k')$, l'inclusion $k' \subset k$ étant finie localement libre,
où
$$
k' = \{x \in k \mid s_i(x) = t_i(x)\text{ pour tout }i \in I\}
$$
Il est facile de voir que $k'$ est un corps.
\end{proof}

\begin{remark}
\label{remark-cannot-decide-yet}
Le Lemme \ref{lemma-point-like-spaces} vaut pour les espaces algébriques décents, voir
Espaces algébriques décents, lemme \ref{decent-spaces-lemma-decent-point-like-spaces}.
En fait, un espace algébrique décent à un seul point est un schéma, voir
Espaces algébriques décents, lemme \ref{decent-spaces-lemma-when-field}.
Cela vaut aussi lorsque $X$ est localement séparé, car un
espace algébrique localement séparé est décent, voir
Espaces algébriques décents, lemme \ref{decent-spaces-lemma-locally-separated-decent}.
\end{remark}







\section{Morphismes \'etales d'espaces algébriques}
\label{section-etale-morphisms}

\noindent
Cette section devrait en réalité figurer dans le chapitre consacré aux morphismes d'espaces
algébriques, mais nous avons besoin de la notion d'espace algébrique \'etale sur un autre
afin de définir le petit site \'etale d'un espace algébrique.
Nous devons donc effectuer quelques travaux préliminaires sur les morphismes \'etales de schémas vers des
espaces algébriques et sur les morphismes \'etales entre espaces algébriques.
Pour davantage de résultats sur les morphismes \'etales d'espaces algébriques, voir
Morphismes d'espaces, section \ref{spaces-morphisms-section-etale}.

\begin{lemma}
\label{lemma-etale-over-space}
Soit $S$ un schéma.
Soit $X$ un espace algébrique sur $S$.
Soient $U$, $U'$ des schémas sur $S$.
\begin{enumerate}
\item Si $U \to U'$ est un morphisme \'etale de schémas, et
si $U' \to X$ est un morphisme \'etale de $U'$ vers $X$, alors la
composée $U \to X$ est un morphisme \'etale de $U$ vers $X$.
\item Si $\varphi : U \to X$ et $\varphi' : U' \to X$ sont des
morphismes \'etales vers $X$, et si $\chi : U \to U'$ est un
morphisme de schémas tel que $\varphi = \varphi' \circ \chi$,
alors $\chi$ est un morphisme \'etale de schémas.
\item Si $\chi : U \to U'$ est un morphisme \'etale surjectif
de schémas et si $\varphi' : U' \to X$ est un morphisme tel que
$\varphi = \varphi' \circ \chi$ soit \'etale, alors $\varphi'$
est \'etale.
\end{enumerate}
\end{lemma}

\begin{proof}
Rappelons que la définition d'un morphisme \'etale d'un schéma vers un
espace algébrique est celle de
Espaces, définition
\ref{spaces-definition-relative-representable-property}
puisque tout morphisme d'un schéma vers un espace algébrique
est représentable.

\medskip\noindent
L'assertion (1) du lemme résulte de ceci, du fait que
les morphismes \'etales sont stables par composition
(Morphismes, lemme
\ref{morphisms-lemma-composition-etale})
et des lemmes
d'Espaces
\ref{spaces-lemma-composition-representable-transformations-property} et
\ref{spaces-lemma-morphism-schemes-gives-representable-transformation-property}
(qui sont formels).

\medskip\noindent
Pour prouver l'assertion (2), choisissons un schéma $W$ sur $S$ et un
morphisme \'etale surjectif $W \to X$. Considérons le changement de base
$\chi_W : W \times_X U \to W \times_X U'$ de $\chi$.
Comme $W \times_X U$ et $W \times_X U'$ sont \'etales sur $W$, nous concluons que
$\chi_W$ est \'etale d'après
Morphismes, lemme \ref{morphisms-lemma-etale-permanence}.
D'autre part, dans le diagramme commutatif
$$
\xymatrix{
W \times_X U \ar[r] \ar[d] & W \times_X U' \ar[d] \\
U \ar[r] & U'
}
$$
les deux flèches verticales sont \'etales et surjectives.
Ainsi, d'après
Descente, lemme \ref{descent-lemma-syntomic-smooth-etale-permanence}
nous concluons que $U \to U'$ est \'etale.

\medskip\noindent
Pour prouver l'assertion (3), choisissons un schéma $W$ sur $S$ et un morphisme $W \to X$.
Comme ci-dessus, considérons le diagramme
$$
\xymatrix{
W \times_X U \ar[r] \ar[d] & W \times_X U' \ar[d] \ar[r] & W \ar[d] \\
U \ar[r] & U' \ar[r] & X
}
$$
Nous savons que $W \times_X U \to W \times_X U'$ est \'etale surjectif
(comme changement de base de $U \to U'$)
et que $W \times_X U \to W$ est \'etale. Ainsi, $W \times_X U' \to W$
est \'etale d'après Descente, lemme
\ref{descent-lemma-syntomic-smooth-etale-permanence}. Par définition,
cela signifie que $\varphi'$ est \'etale.
\end{proof}

\begin{definition}
\label{definition-etale}
Soit $S$ un schéma.
Un morphisme $f : X \to Y$ entre espaces algébriques sur $S$ est
dit {\it \'etale} si et seulement si, pour tout morphisme \'etale
$\varphi : U \to X$, où $U$ est un schéma, la composée
$f \circ \varphi$ est elle aussi \'etale.
\end{definition}

\noindent
Si $X$ et $Y$ sont des schémas, cela coïncide avec la notion usuelle de
morphisme \'etale de schémas. En fait, dès que $X \to Y$ est un morphisme
représentable d'espaces algébriques, cela coïncide avec la notion définie dans
Espaces, définition \ref{spaces-definition-relative-representable-property}.
Cela résulte en combinant le lemme \ref{lemma-etale-local} ci-dessous avec
Espaces, lemme \ref{spaces-lemma-representable-morphisms-spaces-property}.

\begin{lemma}
\label{lemma-etale-local}
Soit $S$ un schéma.
Soit $f : X \to Y$ un morphisme d'espaces algébriques sur $S$.
Les assertions suivantes sont équivalentes :
\begin{enumerate}
\item $f$ est \'etale,
\item il existe un morphisme \'etale surjectif $\varphi : U \to X$,
où $U$ est un schéma, tel que la composée $f \circ \varphi$ soit
\'etale (en tant que morphisme d'espaces algébriques),
\item il existe un morphisme \'etale surjectif $\psi : V \to Y$,
où $V$ est un schéma, tel que le changement de base $V \times_Y X \to V$
soit \'etale (en tant que morphisme d'espaces algébriques),
\item il existe un diagramme commutatif
$$
\xymatrix{
U \ar[d] \ar[r] & V \ar[d] \\
X \ar[r] & Y
}
$$
où $U$, $V$ sont des schémas, les flèches verticales sont \'etales, la
flèche verticale gauche est surjective et la flèche horizontale est \'etale.
\end{enumerate}
\end{lemma}

\begin{proof}
Montrons que (4) implique (1). Supposons donné un diagramme comme en (4).
Soit $W \to X$ un morphisme \'etale, où $W$ est un schéma. Alors
$W \times_X U \to U$ est \'etale. Ainsi, $W \times_X U \to V$ est \'etale
comme composée des morphismes \'etales de schémas $W \times_X U \to U$
et $U \to V$. Par conséquent, $W \times_X U \to Y$ est \'etale d'après le
Lemme \ref{lemma-etale-over-space} (1). Comme, de plus,
la projection $W \times_X U \to W$ est surjective et \'etale, nous concluons
du Lemme \ref{lemma-etale-over-space} (3) que $W \to Y$ est \'etale.

\medskip\noindent
Montrons que (1) implique (4). Supposons (1). Choisissons un diagramme commutatif
$$
\xymatrix{
U \ar[d] \ar[r] & V \ar[d] \\
X \ar[r] & Y
}
$$
où $U \to X$ et $V \to Y$ sont surjectifs et \'etales, voir
Espaces, lemme \ref{spaces-lemma-lift-morphism-presentations}.
Par hypothèse, le morphisme $U \to Y$ est \'etale ;
donc $U \to V$ est \'etale d'après le Lemme \ref{lemma-etale-over-space} (2).

\medskip\noindent
Nous omettons la preuve que (2) et (3) sont aussi équivalents à (1).
\end{proof}

\begin{lemma}
\label{lemma-composition-etale}
La composée de deux morphismes \'etales d'espaces algébriques
est \'etale.
\end{lemma}

\begin{proof}
Cela résulte immédiatement de la définition.
\end{proof}

\begin{lemma}
\label{lemma-base-change-etale}
Le changement de base d'un morphisme \'etale d'espaces algébriques
par un morphisme quelconque d'espaces algébriques est \'etale.
\end{lemma}

\begin{proof}
Soit $X \to Y$ un morphisme \'etale d'espaces algébriques sur $S$.
Soit $Z \to Y$ un morphisme d'espaces algébriques.
Choisissons un schéma $U$ et un morphisme \'etale surjectif $U \to X$.
Choisissons un schéma $W$ et un morphisme \'etale surjectif $W \to Z$.
Alors $U \to Y$ est \'etale ; ainsi, dans le diagramme
$$
\xymatrix{
W \times_Y U \ar[d] \ar[r] & W \ar[d] \\
Z \times_Y X \ar[r] & Z
}
$$
la flèche horizontale supérieure est \'etale.
De plus, la flèche verticale gauche est surjective
et \'etale (nous omettons la vérification). Nous concluons donc que la flèche
horizontale inférieure est \'etale d'après le Lemme \ref{lemma-etale-local}.
\end{proof}

\begin{lemma}
\label{lemma-etale-permanence}
Soit $S$ un schéma. Soient $X, Y, Z$ des espaces algébriques.
Soient $g : X \to Z$, $h : Y \to Z$ des morphismes \'etales, et soit
$f : X \to Y$ un morphisme tel que $h \circ f = g$.
Alors $f$ est \'etale.
\end{lemma}

\begin{proof}
Choisissons un diagramme commutatif
$$
\xymatrix{
U \ar[d] \ar[r]_\chi & V \ar[d] \\
X \ar[r] & Y
}
$$
où $U \to X$ et $V \to Y$ sont surjectifs et \'etales, voir
Espaces, lemme \ref{spaces-lemma-lift-morphism-presentations}.
Par hypothèse, les morphismes $\varphi : U \to X \to Z$ et
$\psi : V \to Y \to Z$ sont \'etales. De plus, $\psi \circ \chi = \varphi$
par notre hypothèse sur $f, g, h$.
Ainsi, $U \to V$ est \'etale d'après le Lemme \ref{lemma-etale-over-space},
assertion (2).
\end{proof}

\begin{lemma}
\label{lemma-etale-open}
Soit $S$ un schéma.
Si $X \to Y$ est un morphisme \'etale d'espaces algébriques sur $S$,
alors l'application associée $|X| \to |Y|$ entre espaces topologiques
est ouverte.
\end{lemma}

\begin{proof}
Cela résulte clairement du diagramme du
lemme \ref{lemma-etale-local}
et du
lemme \ref{lemma-topology-points}.
\end{proof}

\noindent
Enfin, voici un joli lemme. Il est faux qu'un espace algébrique
muni d'un morphisme \'etale vers un schéma soit nécessairement un schéma, voir
Espaces, exemple \ref{spaces-example-non-representable-descent}.
Mais cela est vrai si le but est le spectre d'un corps.

\begin{lemma}
\label{lemma-etale-over-field-scheme}
Soit $S$ un schéma. Soit $X \to \Spec(k)$
un morphisme \'etale sur $S$, où $k$ est un corps.
Alors $X$ est un schéma.
\end{lemma}

\begin{proof}
Soit $U$ un schéma affine et soit $U \to X$ un morphisme \'etale. D'après la
définition \ref{definition-etale}
nous voyons que $U \to \Spec(k)$ est un morphisme \'etale.
Ainsi, $U = \coprod_{i = 1, \ldots, n} \Spec(k_i)$
est une somme disjointe finie de spectres d'extensions finies séparables
$k_i$ de $k$, voir
Morphismes, lemme \ref{morphisms-lemma-etale-over-field}.
Le morphisme $R = U \times_X U \to U \times_{\Spec(k)} U$ est un monomorphisme,
et $U \times_{\Spec(k)} U$ est aussi une somme disjointe finie de
spectres d'extensions finies séparables de $k$. Ainsi, d'après
Schémas, lemme \ref{schemes-lemma-mono-towards-spec-field}
nous voyons que $R$ est de même une somme disjointe finie de
spectres d'extensions finies séparables de $k$.
Ainsi, $U$ et $R$ sont affines, et
les deux projections $R \to U$ sont finies localement libres.
Par conséquent, $U/R$ est un schéma d'après
Groupoïdes, proposition \ref{groupoids-proposition-finite-flat-equivalence}.
D'après
Espaces, lemme \ref{spaces-lemma-finding-opens}
il est aussi un sous-espace ouvert de $X$. D'après le
lemme \ref{lemma-subscheme}
nous concluons que $X$ est un schéma.
\end{proof}













\section{Espaces et recouvrements fpqc}
\label{section-fpqc}

\noindent
Soit $S$ un schéma.
Un espace algébrique sur $S$ est défini comme un faisceau pour la topologie fppf muni de
propriétés supplémentaires. Il n'est donc pas immédiatement clair qu'il vérifie
la condition de faisceau pour la topologie fpqc (voir
Topologies, définition \ref{topologies-definition-sheaf-property-fpqc}).
Dans cette section, nous donnons l'argument de Gabber qui le démontre.
Toutefois, lorsque nous disons que l'espace algébrique $X$ vérifie la
condition de faisceau pour la topologie fpqc, nous ne considérons en réalité que les
recouvrements fpqc $\{f_i : T_i \to T\}_{i \in I}$ tels que $T, T_i$ soient des
objets du grand site $(\Sch/S)_{fppf}$ (conformément à
nos conventions, voir la section \ref{section-conventions}).

\begin{proposition}[Gabber]
\label{proposition-sheaf-fpqc}
Soit $S$ un schéma. Soit $X$ un espace algébrique sur $S$. Alors
$X$ vérifie la condition de faisceau pour la topologie fpqc.
\end{proposition}

\begin{proof}
Puisque $X$ est un faisceau pour la topologie de Zariski, il suffit de montrer
ce qui suit. Étant donné un morphisme plat surjectif de schémas affines
$f : T' \to T$, on a :
$X(T)$ est le noyau du couple d'applications $X(T') \to X(T' \times_T T')$.
Voir Topologies, lemme \ref{topologies-lemma-sheaf-property-fpqc}
(nous omettons ici un petit argument, car le lemme cité
est formulé pour des foncteurs définis sur la catégorie de tous les schémas).

\medskip\noindent
Soient $a, b : T \to X$ deux morphismes tels que $a \circ f = b \circ f$.
Il faut montrer que $a = b$. Considérons le produit fibré
$$
E = X \times_{\Delta_{X/S}, X \times_S X, (a, b)} T.
$$
D'après Espaces, lemme \ref{spaces-lemma-properties-diagonal},
le morphisme $\Delta_{X/S}$ est un monomorphisme représentable. Ainsi,
$E \to T$ est un monomorphisme de schémas. Notre hypothèse
$a \circ f = b \circ f$ implique que $T' \to T$ se factorise (uniquement) par $E$.
Considérons le diagramme commutatif
$$
\xymatrix{
T' \times_T E \ar[r] \ar[d] & E \ar[d] \\
T' \ar[r] \ar@/^5ex/[u] \ar[ru] & T
}
$$
Puisque la projection $T' \times_T E \to T'$ est un monomorphisme
admettant une section, nous concluons que c'est un isomorphisme. Nous en déduisons que
$E \to T$ est un isomorphisme d'après
Descente, lemme \ref{descent-lemma-descending-property-isomorphism}.
Cela signifie que $a = b$, comme voulu.

\medskip\noindent
Soit ensuite $c : T' \to X$ un morphisme tel que les deux composées
$T' \times_T T' \to T' \to X$ soient égales. Il faut trouver un morphisme
$a : T \to X$ dont la composée avec $T' \to T$ soit $c$. Choisissons un
schéma affine $U$ et un morphisme \'etale $U \to X$ tels que l'image de
$|U| \to |X|$ contienne l'image de $|c| : |T'| \to |X|$.
C'est possible d'après les lemmes \ref{lemma-topology-points} et
\ref{lemma-cover-by-union-affines}, le fait qu'une somme disjointe finie de
schémas affines est affine, et le fait que $|T'|$ est quasi-compact
(nous omettons un petit argument). Puisque $U \to X$ est séparé
(lemme \ref{lemma-separated-cover}), on voit que
$$
V = U \times_{X, c} T' \longrightarrow T'
$$
est un morphisme de schémas \'etale, surjectif et séparé
(pour voir qu'il est surjectif, utiliser le lemme \ref{lemma-points-cartesian}
et notre choix de $U \to X$). L'égalité
$c \circ \text{pr}_0 = c \circ \text{pr}_1$ signifie que nous obtenons une
donnée de descente sur $V/T'/T$
(Descente, définition \ref{descent-definition-descent-datum}),
car
\begin{align*}
V \times_{T'} (T' \times_T T')
& =
U \times_{X, c \circ \text{pr}_0} (T' \times_T T') \\
& =
(T' \times_T T') \times_{c \circ \text{pr}_1, X} U \\
& =
(T' \times_T T') \times_{T'} V
\end{align*}
Le morphisme $V \to T'$ est ind-quasi-affine d'après
Compléments sur les morphismes, lemme
\ref{more-morphisms-lemma-etale-separated-ind-quasi-affine}
(car les morphismes \'etales sont localement quasi-finis, voir
Morphismes, lemme \ref{morphisms-lemma-etale-locally-quasi-finite}).
D'après Compléments sur les groupoïdes, lemme \ref{more-groupoids-lemma-ind-quasi-affine},
la donnée de descente est effective. Soit $W \to T$ un morphisme
tel qu'il existe un isomorphisme $\alpha : T' \times_T W \to V$
compatible avec la donnée de descente donnée sur $V$ et la donnée de descente
canonique sur $T' \times_T W$. Alors $W \to T$ est surjectif et \'etale
(Descente, lemmes \ref{descent-lemma-descending-property-surjective} et
\ref{descent-lemma-descending-property-etale}).
Considérons la composée
$$
b' : T' \times_T W \longrightarrow V = U \times_{X, c} T' \longrightarrow U
$$
Les deux composées
$b' \circ (\text{pr}_0, 1), 
b' \circ (\text{pr}_1, 1) :
(T' \times_T T') \times_T W \to T' \times_T W \to U$
coïncident par notre choix de $\alpha$ et la propriété correspondante de $c$
(calcul omis). Ainsi, $b'$ descend en un morphisme $b : W \to U$ d'après
Descente, lemme \ref{descent-lemma-fpqc-universal-effective-epimorphisms}.
Le diagramme
$$
\xymatrix{
T' \times_T W \ar[r] \ar[d] & W \ar[r]_b & U \ar[d] \\
T' \ar[rr]^c &  & X
}
$$
est commutatif. Cela signifie que nous avons démontré l'existence
d'une solution $a$ localement pour la topologie \'etale sur $T$, c'est-à-dire d'un $a' : W \to X$.
Cependant, puisque nous avons démontré l'unicité
dans le premier paragraphe, nous constatons que cette solution locale pour la topologie \'etale
satisfait à la condition de recollement, c'est-à-dire que
$\text{pr}_0^*a' = \text{pr}_1^*a'$ comme éléments de $X(W \times_T W)$.
Puisque $X$ est un faisceau pour la topologie \'etale, il existe un unique $a \in X(T)$ dont la restriction
est $a'$ sur $W$.
\end{proof}









\section{Le site \'etale d'un espace algébrique}
\label{section-etale-site}

\noindent
Dans cette section, nous définissons le petit site \'etale d'un espace algébrique.
C'est l'analogue du petit site \'etale $S_\etale$ d'un schéma.
Le lemme \ref{lemma-etale-over-space} implique que, dans la définition ci-dessous,
tout morphisme entre objets du site \'etale de $X$ est \'etale, et que
tout schéma \'etale sur un objet de $X_\etale$ est encore un objet de
$X_\etale$.

\begin{definition}
\label{definition-etale-site}
Soit $S$ un schéma.
Soit $\Sch_{fppf}$ un grand site fppf contenant $S$,
et soit $\Sch_\etale$ le grand site \'etale correspondant
(c'est-à-dire ayant la même catégorie sous-jacente).
Soit $X$ un espace algébrique sur $S$.
Le {\it petit site \'etale $X_\etale$} de $X$ est défini comme suit :
\begin{enumerate}
\item un objet de $X_\etale$ est un morphisme $\varphi : U \to X$
où $U \in \Ob((\Sch/S)_\etale)$ est un schéma et
$\varphi$ est un morphisme \'etale,
\item un morphisme $(\varphi : U \to X) \to (\varphi' : U' \to X)$
est donné par un morphisme de schémas $\chi : U \to U'$ tel que
$\varphi = \varphi' \circ \chi$, et
\item une famille de morphismes $\{(U_i \to X) \to (U \to X)\}_{i \in I}$
de $X_\etale$ est un recouvrement si et seulement si $\{U_i \to U\}_{i \in I}$
est un recouvrement de $(\Sch/S)_\etale$.
\end{enumerate}
\end{definition}

\noindent
Une conséquence de notre choix est que le site \'etale d'un espace algébrique
n'a en général pas d'objet final ! En revanche, si $X$ est
un schéma, la définition ci-dessus coïncide avec
Topologies, définition \ref{topologies-definition-big-small-etale}.

\medskip\noindent
Le site précédent est celui que nous employons par défaut, mais nous utiliserons aussi
deux variantes. On peut considérer tous les {\it espaces algébriques}
$U$ \'etales sur $X$, ce qui donne le site
$X_{spaces, \etale}$ défini ci-dessous, ou tous les {\it schémas affines}
$U$ \'etales sur $X$, ce qui donne le site
$X_{affine, \etale}$ défini ci-dessous. La première de ces deux notions
sert à étudier la fonctorialité du petit site \'etale, voir le
lemme \ref{lemma-functoriality-etale-site}.

\begin{definition}
\label{definition-spaces-etale-site}
Soit $S$ un schéma.
Soit $\Sch_{fppf}$ un grand site fppf contenant $S$,
et soit $\Sch_\etale$ le grand site \'etale correspondant
(c'est-à-dire ayant la même catégorie sous-jacente).
Soit $X$ un espace algébrique sur $S$.
Le site {\it $X_{spaces, \etale}$} de $X$ est défini comme suit :
\begin{enumerate}
\item un objet de $X_{spaces, \etale}$ est un morphisme
$\varphi : U \to X$, où $U$ est un espace algébrique sur $S$ et
$\varphi$ est un morphisme \'etale d'espaces algébriques sur $S$,
\item un morphisme $(\varphi : U \to X) \to (\varphi' : U' \to X)$ de
$X_{spaces, \etale}$ est donné par un morphisme d'espaces algébriques
$\chi : U \to U'$ tel que $\varphi = \varphi' \circ \chi$, et
\item une famille de morphismes
$\{\varphi_i : (U_i \to X) \to (U \to X)\}_{i \in I}$
de $X_{spaces, \etale}$ est un recouvrement si et seulement si
$|U| = \bigcup \varphi_i(|U_i|)$.
\end{enumerate}
Comme d'habitude, nous choisissons un ensemble de recouvrements de ce type, contenant au moins
les recouvrements de $X_\etale$, comme dans
Ensembles, lemme \ref{sets-lemma-coverings-site},
afin de faire de $X_{spaces, \etale}$ un site.
\end{definition}

\noindent
Puisque le morphisme identité de $X$ est \'etale, il est clair que
$X_{spaces, \etale}$ possède un objet final.
Montrons immédiatement que le topos correspondant est équivalent au
petit topos \'etale de $X$.

\begin{lemma}
\label{lemma-compare-etale-sites}
Le foncteur
$$
X_\etale \longrightarrow X_{spaces, \etale}, \quad
U/X \longmapsto U/X
$$
est un foncteur cocontinu spécial
(Sites, définition \ref{sites-definition-special-cocontinuous-functor})
et induit donc une équivalence de topos
$\Sh(X_\etale) \to \Sh(X_{spaces, \etale})$.
\end{lemma}

\begin{proof}
Il faut montrer que le foncteur satisfait aux hypothèses (1) -- (5) du
Sites, lemme \ref{sites-lemma-equivalence}.
Il est clair que le foncteur est continu et cocontinu, ce qui
prouve les hypothèses (1) et (2).
Les hypothèses (3) et (4) sont satisfaites simplement parce que le foncteur est pleinement fidèle.
L'hypothèse (5) est satisfaite parce que, par définition, un espace algébrique admet
un recouvrement par un schéma.
\end{proof}

\begin{remark}
\label{remark-explain-equivalence}
Expliquons le sens du lemme \ref{lemma-compare-etale-sites}.
Soit $S$ un schéma et soit $X$ un espace algébrique sur $S$.
Soit $\mathcal{F}$ un faisceau sur le petit site \'etale $X_\etale$ de
$X$. Le lemme affirme qu'il existe un unique faisceau $\mathcal{F}'$ sur
$X_{spaces, \etale}$ dont la restriction est $\mathcal{F}$ sur la
sous-catégorie $X_\etale$. Si $U \to X$ est un morphisme \'etale
d'espaces algébriques, comment calculer $\mathcal{F}'(U)$ ? Par définition
d'un espace algébrique, il existe un schéma $U'$ et un morphisme
\'etale surjectif $U' \to U$. Alors $\{U' \to U\}$ est un recouvrement de
$X_{spaces, \etale}$, et l'on obtient donc un diagramme exact
$$
\xymatrix{
\mathcal{F}'(U) \ar[r] &
\mathcal{F}(U') \ar@<1ex>[r] \ar@<-1ex>[r] &
\mathcal{F}(U' \times_U U').
}
$$
Notons que $U' \times_U U'$ est un schéma ; nous pouvons donc
écrire $\mathcal{F}$ et non $\mathcal{F}'$.
Nous voyons ainsi comment calculer $\mathcal{F}'$
à partir du faisceau $\mathcal{F}$.
\end{remark}

\begin{definition}
\label{definition-affine-etale-site}
Soit $S$ un schéma.
Soit $\Sch_{fppf}$ un grand site fppf contenant $S$,
et soit $\Sch_\etale$ le grand site \'etale correspondant
(c'est-à-dire ayant la même catégorie sous-jacente).
Soit $X$ un espace algébrique sur $S$.
Le site {\it $X_{affine, \etale}$} de $X$ est défini comme suit :
\begin{enumerate}
\item un objet de $X_{affine, \etale}$ est un morphisme
$\varphi : U \to X$, où $U \in \Ob((\Sch/S)_\etale)$ est un schéma affine et
$\varphi$ est un morphisme \'etale,
\item un morphisme $(\varphi : U \to X) \to (\varphi' : U' \to X)$ de
$X_{affine, \etale}$ est donné par un morphisme de schémas
$\chi : U \to U'$ tel que $\varphi = \varphi' \circ \chi$, et
\item une famille de morphismes
$\{\varphi_i : (U_i \to X) \to (U \to X)\}_{i \in I}$
de $X_{affine, \etale}$ est un recouvrement si et seulement si
$\{U_i \to U\}$ est un recouvrement \'etale standard, voir
Topologies, définition \ref{topologies-definition-standard-etale}.
\end{enumerate}
Comme d'habitude, nous choisissons un ensemble de recouvrements de ce type, comme dans
Ensembles, lemme \ref{sets-lemma-coverings-site},
afin de faire de $X_{affine, \etale}$ un site.
\end{definition}

\begin{lemma}
\label{lemma-alternative}
Soit $S$ un schéma. Soit $X$ un espace algébrique sur $S$.
Le foncteur $X_{affine, \etale} \to X_\etale$
est cocontinu spécial et induit une équivalence de topos de
$\Sh(X_{affine, \etale})$ vers $\Sh(X_\etale)$.
\end{lemma}

\begin{proof}
Omis. Indication : comparer avec la preuve de
Topologies, lemme \ref{topologies-lemma-affine-big-site-etale}.
\end{proof}

\begin{definition}
\label{definition-etale-topos}
Soit $S$ un schéma. Soit $X$ un espace algébrique sur $S$.
Le {\it topos \'etale} de $X$, ou plus précisément le
{\it petit topos \'etale} de $X$, est la catégorie
$\Sh(X_\etale)$
des faisceaux d'ensembles sur $X_\etale$.
\end{definition}

\noindent
D'après le
lemme \ref{lemma-compare-etale-sites},
on a
$\Sh(X_\etale) = \Sh(X_{spaces, \etale})$,
de sorte que l'on peut aussi voir ce topos comme la catégorie des faisceaux d'ensembles sur
$X_{spaces, \etale}$. De même, d'après le
lemme \ref{lemma-alternative},
on voit que
$\Sh(X_\etale) = \Sh(X_{affine, \etale})$.
Ce topos est fonctoriel par rapport aux morphismes
d'espaces algébriques. Voici un énoncé précis.

\begin{lemma}
\label{lemma-functoriality-etale-site}
Soit $S$ un schéma.
Soit $f : X \to Y$ un morphisme d'espaces algébriques sur $S$.
\begin{enumerate}
\item Le foncteur continu
$$
Y_{spaces, \etale} \longrightarrow X_{spaces, \etale}, \quad
V \longmapsto X \times_Y V
$$
induit un morphisme de sites
$$
f_{spaces, \etale} :
X_{spaces, \etale}
\to
Y_{spaces, \etale}.
$$
\item La règle $f \mapsto f_{spaces, \etale}$ est compatible aux
compositions, autrement dit $(f \circ g)_{spaces, \etale}
= f_{spaces, \etale} \circ g_{spaces, \etale}$ (voir
Sites, définition \ref{sites-definition-composition-morphisms-sites}).
\item Le morphisme de topos associé à $f_{spaces, \etale}$
induit, via le lemme \ref{lemma-compare-etale-sites}, un morphisme de topos
$f_{small} : \Sh(X_\etale) \to \Sh(Y_\etale)$
dont la construction est compatible aux compositions.
\item Si $f$ est un morphisme représentable d'espaces algébriques,
alors $f_{small}$ provient d'un morphisme de sites
$X_\etale \to Y_\etale$,
correspondant au foncteur continu $V \mapsto X \times_Y V$.
\end{enumerate}
\end{lemma}

\begin{proof}
Montrons que le foncteur décrit en (1) satisfait aux hypothèses
de Sites, proposition \ref{sites-proposition-get-morphism}.
Il faut donc montrer que
$Y_{spaces, \etale}$ possède un objet final (à savoir $Y$) et que
le foncteur transforme celui-ci en un objet final de $X_{spaces, \etale}$
(à savoir $X$). C'est clair, car $X \times_Y Y = X$ dans toute catégorie.
Il faut ensuite montrer que $Y_{spaces, \etale}$ possède des produits fibrés.
C'est vrai puisque la catégorie des espaces algébriques possède des produits fibrés,
et puisque $V \times_Y V'$ est \'etale sur $Y$ si $V$ et $V'$ sont \'etales
sur $Y$ (voir les lemmes \ref{lemma-composition-etale} et
\ref{lemma-base-change-etale} ci-dessus).
La proposition s'applique donc et fournit un morphisme
de sites comme en (1).

\medskip\noindent
L'assertion (2) s'obtient en explicitant les définitions.
L'assertion (3) est claire en utilisant les équivalences pour $X$ et $Y$
du lemme \ref{lemma-compare-etale-sites} ci-dessus.
L'assertion (4) en résulte car, si $f$ est représentable, les
foncteurs ci-dessus s'insèrent dans un diagramme commutatif
$$
\xymatrix{
X_\etale \ar[r] &
X_{spaces, \etale} \\
Y_\etale \ar[r] \ar[u] &
Y_{spaces, \etale} \ar[u]
}
$$
de catégories.
\end{proof}

\noindent
On peut préciser quelque peu le lemme précédent lorsqu'on décrit
la relation entre les faisceaux sur $X$ et ceux sur $Y$.
Plus précisément, on peut la formuler en termes de $f$-applications ; comparer
Faisceaux, définition \ref{sheaves-definition-f-map}, comme suit.

\begin{definition}
\label{definition-f-map}
Soit $S$ un schéma.
Soit $f : X \to Y$ un morphisme d'espaces algébriques sur $S$.
Soit $\mathcal{F}$ un faisceau d'ensembles sur $X_\etale$ et
soit $\mathcal{G}$ un faisceau d'ensembles sur $Y_\etale$.
Une {\it $f$-application $\varphi : \mathcal{G} \to \mathcal{F}$}
est une collection d'applications
$\varphi_{(U, V, g)} : \mathcal{G}(V) \to \mathcal{F}(U)$
indexée par les diagrammes commutatifs
$$
\xymatrix{
U \ar[d]_g \ar[r] & X \ar[d]^f \\
V \ar[r] & Y
}
$$
où $U \in X_\etale$, $V \in Y_\etale$, telle que, pour tout
diagramme prolongé
$$
\xymatrix{
U' \ar[r] \ar[d]_{g'} & U \ar[d]_g \ar[r] & X \ar[d]^f \\
V' \ar[r] & V \ar[r] & Y
}
$$
où $V' \to V$ et $U' \to U$ sont des morphismes \'etales de schémas, le diagramme
$$
\xymatrix{
\mathcal{G}(V)
\ar[rr]_{\varphi_{(U, V, g)}}
\ar[d]_{\text{restriction de }\mathcal{G}} & &
\mathcal{F}(U)
\ar[d]^{\text{restriction de }\mathcal{F}} \\
\mathcal{G}(V')
\ar[rr]^{\varphi_{(U', V', g')}} & &
\mathcal{F}(U')
}
$$
est commutatif.
\end{definition}

\begin{lemma}
\label{lemma-f-map}
Soit $S$ un schéma.
Soit $f : X \to Y$ un morphisme d'espaces algébriques sur $S$.
Soit $\mathcal{F}$ un faisceau d'ensembles sur $X_\etale$ et
soit $\mathcal{G}$ un faisceau d'ensembles sur $Y_\etale$.
Il existe des bijections canoniques entre les trois ensembles suivants :
\begin{enumerate}
\item l'ensemble des morphismes $\mathcal{G} \to f_{small, *}\mathcal{F}$ ;
\item l'ensemble des morphismes $f_{small}^{-1}\mathcal{G} \to \mathcal{F}$ ;
\item l'ensemble des $f$-applications $\varphi : \mathcal{G} \to \mathcal{F}$.
\end{enumerate}
\end{lemma}

\begin{proof}
Notons que (1) et (2) coïncident parce que les foncteurs $f_{small, *}$
et $f_{small}^{-1}$ forment un couple de foncteurs adjoints.
Supposons que $\alpha : f_{small}^{-1}\mathcal{G} \to \mathcal{F}$
soit un morphisme de faisceaux sur $X_\etale$. Considérons un diagramme
$$
\xymatrix{
U \ar[d]_g \ar[r]_{j_U} & X \ar[d]^f \\
V \ar[r]^{j_V} & Y
}
$$
comme dans la définition \ref{definition-f-map}.
La commutativité du diagramme fournit également un morphisme
$g_{small}^{-1}(j_V)^{-1}\mathcal{G} \to (j_U)^{-1}\mathcal{F}$
(comparer Sites, section \ref{sites-section-localize}, pour la
description des foncteurs de localisation). Nous obtenons donc en particulier
un morphisme
$\varphi_{(U, V, g)} :
\mathcal{G}(V) = (j_V)^{-1}\mathcal{G}(V)
\to
(j_U)^{-1}\mathcal{F}(U) = \mathcal{F}(U)$.
Nous omettons de vérifier que cette règle est compatible aux
restrictions ultérieures et définit une $f$-application de $\mathcal{G}$ vers
$\mathcal{F}$.

\medskip\noindent
Réciproquement, supposons donnée une $f$-application
$\varphi = (\varphi_{(U, V, g)})$.
Désignons par $\mathcal{G}'$ (resp.\ $\mathcal{F}'$) le prolongement de
$\mathcal{G}$ (resp.\ $\mathcal{F}$) à $Y_{spaces, \etale}$
(resp.\ $X_{spaces, \etale}$), voir le
lemme \ref{lemma-compare-etale-sites}.
Il faut alors construire un morphisme de faisceaux
$$
\mathcal{G}' \longrightarrow (f_{spaces, \etale})_*\mathcal{F}'
$$
Pour cela, soit $V \to Y$ un morphisme \'etale d'espaces algébriques.
Il faut construire une application d'ensembles
$$
\mathcal{G}'(V) \to \mathcal{F}'(X \times_Y V)
$$
Choisissons un morphisme \'etale surjectif $V' \to V$ avec $V'$ schéma,
puis un morphisme \'etale surjectif
$U' \to X \times_Y V'$ avec $U'$ schéma. Posons $U'' = U' \times_{X \times_Y V} U'$
et $V'' = V' \times_V V'$. Nous obtenons un morphisme de schémas $g' : U' \to V'$ ainsi qu'un morphisme de schémas
$$
g'' : U' \times_{X \times_Y V} U' \longrightarrow V' \times_V V'
$$
Considérons le diagramme suivant
$$
\xymatrix{
\mathcal{F}'(X \times_Y V) \ar[r] &
\mathcal{F}(U') \ar@<1ex>[r] \ar@<-1ex>[r] &
\mathcal{F}(U' \times_{X \times_Y V} U') \\
\mathcal{G}'(V) \ar[r] \ar@{..>}[u] &
\mathcal{G}(V') \ar@<1ex>[r] \ar@<-1ex>[r] \ar[u]_{\varphi_{(U', V', g')}} &
\mathcal{G}(V' \times_V V') \ar[u]_{\varphi_{(U'', V'', g'')}}
}
$$
La compatibilité des applications $\varphi_{...}$
aux restrictions montre que les deux carrés de droite sont commutatifs.
La définition des recouvrements de $X_{spaces, \etale}$ montre que
les lignes horizontales sont des diagrammes exacts. On obtient donc
la flèche pointillée. Nous laissons au lecteur le soin de montrer que ces
flèches sont compatibles aux applications de restriction.
\end{proof}

\noindent
Si le morphisme d'espaces algébriques $X \to Y$ est \'etale, alors le morphisme
de topos $\Sh(X_\etale) \to \Sh(Y_\etale)$
est un morphisme de localisation. Voici un énoncé précis.

\begin{lemma}
\label{lemma-etale-morphism-topoi}
Soit $S$ un schéma, et soit $f : X \to Y$ un morphisme d'espaces algébriques
sur $S$. Supposons $f$ \'etale. Il existe alors un foncteur
$$
j : X_\etale \to Y_\etale, \quad
(\varphi : U \to X) \mapsto (f \circ \varphi : U \to Y)
$$
qui est cocontinu. Le morphisme de topos $f_{small}$ est le
morphisme de topos associé à $j$, voir
Sites, lemme \ref{sites-lemma-cocontinuous-morphism-topoi}.
De plus, $j$ est aussi continu ; par conséquent,
Sites, lemme \ref{sites-lemma-when-shriek},
s'applique. En particulier, $f_{small}^{-1}\mathcal{G}(U) = \mathcal{G}(jU)$
pour tout faisceau $\mathcal{G}$ sur $Y_\etale$.
\end{lemma}

\begin{proof}
Notons que notre définition même d'un morphisme \'etale d'espaces algébriques
(définition \ref{definition-etale}) assure
que la règle donnée définit bien le foncteur $j$ indiqué.
Il est clair que $j$ est cocontinu et continu, simplement parce qu'un recouvrement
$\{U_i \to U\}$ de $j(\varphi : U \to X)$ dans $Y_\etale$ est
la même chose qu'un recouvrement de $(\varphi : U \to X)$ dans $X_\etale$. Il
reste à montrer que $j$ induit le même morphisme de topos que $f_{small}$.
Considérons pour cela le diagramme
$$
\xymatrix{
X_\etale \ar[r] \ar[d]^j &
X_{spaces, \etale} \ar@/_/[d]_{j_{spaces}} \\
Y_\etale \ar[r] &
Y_{spaces, \etale} \ar@/_/[u]_{v : V \mapsto X \times_Y V}
}
$$
de catégories. Ici, le foncteur $j_{spaces}$ est le prolongement évident de $j$
à la catégorie $X_{spaces, \etale}$. Le carré intérieur est donc
commutatif. En fait, $j_{spaces}$ s'identifie au
foncteur de localisation
$j_X : Y_{spaces, \etale}/X \to Y_{spaces, \etale}$
étudié dans
Sites, section \ref{sites-section-localize}.
Par conséquent, d'après
Sites, lemme \ref{sites-lemma-localize-given-products},
le foncteur cocontinu $j_{spaces}$ et le foncteur $v$ du diagramme
induisent le même morphisme de topos. D'après
Sites, lemme \ref{sites-lemma-composition-cocontinuous},
la commutativité du carré intérieur (formé de foncteurs cocontinus
entre sites) fournit un diagramme commutatif des morphismes de topos associés.
La construction de $f_{small}$ dans le
lemme \ref{lemma-functoriality-etale-site} achève donc la preuve.
\end{proof}

\noindent
Le lemme précédent dit que l'image réciproque de $\mathcal{G}$ par un morphisme \'etale
$f : X \to Y$ d'espaces algébriques est simplement la restriction de $\mathcal{G}$
à la catégorie $X_\etale$. Nous emploierons souvent l'abréviation
\begin{equation}
\label{equation-restrict}
\mathcal{G}|_{X_\etale} = f_{small}^{-1}\mathcal{G}
\end{equation}
pour le signaler. Notons que le foncteur
$j : X_\etale \to Y_\etale$
du lemme est alors fidèle, mais n'est pas pleinement fidèle en
général. Nous l'étudierons plus techniquement dans la
section \ref{section-localize}.

\begin{lemma}
\label{lemma-pushforward-etale-base-change}
Soit $S$ un schéma. Soit
$$
\xymatrix{
X' \ar[r] \ar[d]_{f'} & X \ar[d]^f \\
Y' \ar[r]^g & Y
}
$$
un carré cartésien d'espaces algébriques sur $S$. Soit
$\mathcal{F}$ un faisceau sur $X_\etale$. Si $g$ est \'etale, alors
\begin{enumerate}
\item $f'_{small, *}(\mathcal{F}|_{X'}) = (f_{small, *}\mathcal{F})|_{Y'}$
dans $\Sh(Y'_\etale)$\footnote{On a aussi
$(f')_{small}^{-1}(\mathcal{G}|_{Y'}) = (f_{small}^{-1}\mathcal{G})|_{X'}$
en raison de la commutativité du diagramme et de (\ref{equation-restrict})}, et
\item si $\mathcal{F}$ est un faisceau abélien, alors
$R^if'_{small, *}(\mathcal{F}|_{X'}) = (R^if_{small, *}\mathcal{F})|_{Y'}$.
\end{enumerate}
\end{lemma}

\begin{proof}
Considérons le diagramme de foncteurs suivant
$$
\xymatrix{
X'_{spaces, \etale} \ar[r]_j &
X_{spaces, \etale} \\
Y'_{spaces, \etale} \ar[r]^j \ar[u]^{V' \mapsto V' \times_{Y'} X'} &
Y_{spaces, \etale} \ar[u]_{V \mapsto V \times_Y X}
}
$$
Les flèches horizontales sont des foncteurs de localisation et les flèches verticales induisent
des morphismes de sites. La dernière assertion de
Sites, lemme \ref{sites-lemma-localize-morphism},
donne donc (1). Pour obtenir (2), appliquons (1) à une résolution injective de $\mathcal{F}$
et utilisons que la restriction est exacte et préserve les injectifs (voir
Cohomologie sur les sites, lemme \ref{sites-cohomology-lemma-cohomology-of-open}).
\end{proof}

\noindent
Le lemme suivant dit que l'on peut considérer un faisceau sur le petit
site \'etale d'un espace algébrique comme une famille compatible de faisceaux
sur les petits sites \'etales des schémas \'etales sur cet espace. Notons bien
que tous les morphismes de comparaison $c_f$ du lemme sont des isomorphismes,
ce qui est compatible avec
Topologies, lemme \ref{topologies-lemma-characterize-sheaf-big-etale},
et avec le fait que tous les morphismes entre objets de $X_\etale$
sont \'etales.

\begin{lemma}
\label{lemma-characterize-sheaf-small-etale}
Soit $S$ un schéma. Soit $X$ un espace algébrique sur $S$.
Un faisceau $\mathcal{F}$ sur $X_\etale$ équivaut aux données suivantes :
\begin{enumerate}
\item pour tout $U \in \Ob(X_\etale)$, un faisceau
$\mathcal{F}_U$ sur $U_\etale$ ;
\item pour tout $f : U' \to U$ dans $X_\etale$, un isomorphisme
$c_f : f_{small}^{-1}\mathcal{F}_U \to \mathcal{F}_{U'}$.
\end{enumerate}
Ces données sont soumises à la condition que, pour tous $f : U' \to U$
et $g : U'' \to U'$ dans $X_\etale$, la composée
$c_g \circ g_{small}^{-1} c_f$ soit égale à $c_{f \circ g}$.
\end{lemma}

\begin{proof}
On peut interpréter $g_{small}^{-1}$ comme dans le lemme \ref{lemma-etale-morphism-topoi}.
Le lemme résulte alors d'un fait général sur les sites, voir
Sites, lemme \ref{sites-lemma-glue-sheaves-absolute}.
\end{proof}

\noindent
Soit $S$ un schéma. Soit $X$ un espace algébrique sur $S$.
Soit $X = U/R$ une présentation de $X$ provenant d'un morphisme
\'etale surjectif quelconque $\varphi : U \to X$, voir
Espaces, définition \ref{spaces-definition-presentation}.
On obtient en particulier un groupoïde $(U, R, s, t, c, e, i)$ tel que
$j = (t, s) : R \to U \times_S U$, voir
Groupoïdes, lemme \ref{groupoids-lemma-equivalence-groupoid}.

\begin{lemma}
\label{lemma-descent-sheaf}
Avec $S$, $\varphi : U \to X$ et $(U, R, s, t, c, e, i)$ comme ci-dessus,
pour tout faisceau $\mathcal{F}$ sur $X_\etale$, le
faisceau\footnote{Dans ce lemme
et sa preuve, nous écrivons simplement $\varphi^{-1}$ au lieu de $\varphi_{small}^{-1}$,
et de même pour toutes les autres images réciproques.}
$\mathcal{G} = \varphi^{-1}\mathcal{F}$ est muni d'un
isomorphisme canonique
$$
\alpha :
t^{-1}\mathcal{G}
\longrightarrow
s^{-1}\mathcal{G}
$$
tel que le diagramme
$$
\xymatrix{
& \text{pr}_1^{-1}t^{-1}\mathcal{G} \ar[r]_-{\text{pr}_1^{-1}\alpha} &
\text{pr}_1^{-1}s^{-1}\mathcal{G} \ar@{=}[rd] & \\
\text{pr}_0^{-1}s^{-1}\mathcal{G} \ar@{=}[ru] & & &
c^{-1}s^{-1}\mathcal{G} \\
&
\text{pr}_0^{-1}t^{-1}\mathcal{G} \ar[lu]^{\text{pr}_0^{-1}\alpha} \ar@{=}[r] &
c^{-1}t^{-1}\mathcal{G} \ar[ru]_{c^{-1}\alpha}
}
$$
est commutatif. Le foncteur $\mathcal{F} \mapsto (\mathcal{G}, \alpha)$
définit une équivalence de catégories entre les faisceaux sur
$X_\etale$ et les couples $(\mathcal{G}, \alpha)$ ci-dessus.
\end{lemma}

\begin{proof}[Première preuve du lemme \ref{lemma-descent-sheaf}]
Soit $\mathcal{C} = X_{spaces, \etale}$. D'après le
lemme \ref{lemma-etale-morphism-topoi}
et sa preuve, on a $U_{spaces, \etale} = \mathcal{C}/U$
et le foncteur image réciproque $\varphi^{-1}$ est simplement le foncteur de restriction.
De plus, $\{U \to X\}$ est un recouvrement du site $\mathcal{C}$ et
$R = U \times_X U$. L'isomorphisme $\alpha$ est simplement l'identification
canonique
$$
\left(\mathcal{F}|_{\mathcal{C}/U}\right)|_{\mathcal{C}/U \times_X U}
=
\left(\mathcal{F}|_{\mathcal{C}/U}\right)|_{\mathcal{C}/U \times_X U}
$$
et la commutativité du diagramme est la condition de cocycle pour les données
de recollement. Ce lemme est donc un cas particulier du recollement des faisceaux, voir
Sites, section \ref{sites-section-glueing-sheaves}.
\end{proof}

\begin{proof}[Seconde preuve du lemme \ref{lemma-descent-sheaf}]
L'existence de $\alpha$ vient de l'égalité
$\varphi \circ t = \varphi \circ s$ et de la fonctorialité de l'image réciproque
par rapport au morphisme, voir le
lemme \ref{lemma-functoriality-etale-site}.
De la même manière, c'est-à-dire par fonctorialité de l'image réciproque, on voit
que l'isomorphisme $\alpha$ s'insère dans le diagramme commutatif.
La construction $\mathcal{F} \mapsto (\varphi^{-1}\mathcal{F}, \alpha)$
est manifestement fonctorielle en le faisceau $\mathcal{F}$.
On obtient ainsi le foncteur.

\medskip\noindent
Réciproquement, supposons que $(\mathcal{G}, \alpha)$ soit un couple.
Soit $V \to X$ un objet de $X_\etale$.
Alors le morphisme $V' = U \times_X V \to V$ est un morphisme \'etale surjectif
de schémas, et $\{V' \to V\}$ est donc un recouvrement \'etale
de $V$. Posons $\mathcal{G}' = (V' \to V)^{-1}\mathcal{G}$.
Puisque $R = U \times_X U$, avec $t = \text{pr}_0$
et $s = \text{pr}_1$, on a $V' \times_V V' = R \times_X V$,
les applications de projection $s', t' : V' \times_V V' \to V'$ étant obtenues par changement de base
à partir de $t$ et $s$. Ainsi, $\alpha$ s'induit en un isomorphisme
$\alpha' : (t')^{-1}\mathcal{G}' \to (s')^{-1}\mathcal{G}'$. Cela étant posé,
nous définissons simplement
$$
\xymatrix{
\mathcal{F}(V) \ar@{=}[r] &
\text{Noyau du couple}(\mathcal{G}(V') \ar@<1ex>[r] \ar@<-1ex>[r] &
\mathcal{G}(V' \times_V V')).
}
$$
Nous omettons de vérifier que cela définit un faisceau. Pour voir que
$\mathcal{G}(V) = \mathcal{F}(V)$ s'il existe un morphisme $V \to U$,
notons que le noyau du couple est alors
$H^0(\{V' \to V\}, \mathcal{G}) = \mathcal{G}(V)$.
\end{proof}







\section{Points du petit site \'etale}
\label{section-points-small-etale-site}

\noindent
Cette section est l'analogue de
Cohomologie \'etale, section
\ref{etale-cohomology-section-stalks}.

\begin{definition}
\label{definition-geometric-point}
Soit $S$ un schéma. Soit $X$ un espace algébrique sur $S$.
\begin{enumerate}
\item Un {\it point géométrique} de $X$ est un morphisme
$\overline{x} : \Spec(k) \to X$, où $k$ est un corps algébriquement
clos. Nous commettons souvent un abus de notation en
écrivant $\overline{x} = \Spec(k)$.
\item À tout point géométrique $\overline{x}$ correspond le point
« image » $x \in |X|$. On dit que $\overline{x}$ est un
{\it point géométrique au-dessus de $x$}.
\end{enumerate}
\end{definition}

\noindent
On peut prendre les fibres des faisceaux sur $X_\etale$
en des points géométriques exactement comme dans le cas
du petit site \'etale d'un schéma. Pour cela, nous
définissons comme suit la notion de voisinage \'etale.

\begin{definition}
\label{definition-etale-neighbourhood}
Soit $S$ un schéma. Soit $X$ un espace algébrique sur $S$.
Soit $\overline{x}$ un point géométrique de $X$.
\begin{enumerate}
\item Un {\it voisinage \'etale} de $\overline{x}$
dans $X$ est un diagramme commutatif
$$
\xymatrix{
& U \ar[d]^\varphi \\
{\bar x} \ar[r]^{\bar x} \ar[ur]^{\bar u} & X
}
$$
où $\varphi$ est un morphisme \'etale d'espaces algébriques sur $S$.
Nous noterons $\varphi : (U, \overline{u}) \to (X, \overline{x})$
une telle situation.
\item Un {\it morphisme de voisinages \'etales}
$(U, \overline{u}) \to (U', \overline{u}')$
est un $X$-morphisme $h : U \to U'$
tel que $\overline{u}' = h \circ \overline{u}$.
\end{enumerate}
\end{definition}

\noindent
Notons que nous permettons à $U$ d'être un espace algébrique. Lorsque nous
prenons les fibres d'un faisceau sur $X_\etale$, nous devons nous restreindre aux $U$ qui
appartiennent à $X_\etale$ et ne considérons donc, dans ce cas,
que des schémas $U$. On peut aussi travailler avec le site
$X_{spaces, \etale}$ et considérer tous les voisinages \'etales. Cela
ne fait aucune différence, d'après la dernière assertion du lemme suivant.

\begin{lemma}
\label{lemma-cofinal-etale}
Soit $S$ un schéma. Soit $X$ un espace algébrique sur $S$.
Soit $\overline{x}$ un point géométrique de $X$.
La catégorie des voisinages \'etales est cofiltrante. Plus précisément :
\begin{enumerate}
\item Soient $(U_i, \overline{u}_i)_{i = 1, 2}$ deux voisinages \'etales de
$\overline{x}$ dans $X$. Il existe alors un troisième voisinage \'etale
$(U, \overline{u})$ et des morphismes
$(U, \overline{u}) \to (U_i, \overline{u}_i)$, $i = 1, 2$.
\item Soient $h_1, h_2: (U, \overline{u}) \to (U', \overline{u}')$ deux
morphismes entre voisinages \'etales de $\overline{x}$. Il existe alors un
voisinage \'etale $(U'', \overline{u}'')$ et un morphisme
$h : (U'', \overline{u}'') \to (U, \overline{u})$
qui égalise $h_1$ et $h_2$, c'est-à-dire tel que
$h_1 \circ h = h_2 \circ h$.
\end{enumerate}
De plus, pour tout voisinage \'etale
$(U, \overline{u}) \to (X, \overline{x})$
il existe un morphisme de voisinages \'etales
$(U', \overline{u}') \to (U, \overline{u})$
où $U'$ est un schéma.
\end{lemma}

\begin{proof}
Pour (1), considérons le produit fibré $U = U_1 \times_X U_2$.
Il est \'etale sur $U_1$ et sur $U_2$, car les morphismes \'etales sont
stables par changement de base et par composition, voir les
lemmes \ref{lemma-base-change-etale} et \ref{lemma-composition-etale}.
Le morphisme $\overline{u} \to U$ défini par $(\overline{u}_1, \overline{u}_2)$
en fait un voisinage \'etale muni de morphismes vers
$U_1$ et $U_2$.

\medskip\noindent
Pour (2), définissons $U''$ comme le produit fibré
$$
\xymatrix{
U'' \ar[r] \ar[d] & U \ar[d]^{(h_1, h_2)} \\
U' \ar[r]^-\Delta & U' \times_X U'.
}
$$
Comme les composés de $\overline{u}$ et $\overline{u}'$ vers $X$ valent $\overline{x}$,
$\overline{u}'' = (\overline{u}, \overline{u}')$ est un point géométrique
de $U''$. En particulier, $U'' \not = \emptyset$.
De plus, comme $U'$ est \'etale sur $X$, le produit fibré
$U'\times_X U'$ l'est aussi (comme on l'a vu ci-dessus pour $U_1 \times_X U_2$).
La flèche verticale $(h_1, h_2)$ est donc \'etale d'après le
lemme \ref{lemma-etale-permanence}.
Ainsi, $U''$ est \'etale sur $U'$ par changement de base, donc aussi
\'etale sur $X$ (puisque les composés de morphismes \'etales sont \'etales).
Par conséquent, $(U'', \overline{u}'')$ résout le problème posé en (2).

\medskip\noindent
Pour la dernière assertion, choisissons un morphisme \'etale surjectif
$U' \to U$, où $U'$ est un schéma. Alors
$U' \times_U \overline{u}$ est un schéma \'etale et surjectif sur
$\overline{u} = \Spec(k)$, où $k$ est algébriquement clos.
Il s'ensuit (voir
Morphismes, lemme \ref{morphisms-lemma-etale-over-field})
que $U' \times_U \overline{u} \to \overline{u}$ possède une section,
qui fournit le point $\overline{u}'$ voulu.
\end{proof}

\begin{lemma}
\label{lemma-geometric-lift-to-usual}
Soit $S$ un schéma. Soit $X$ un espace algébrique sur $S$.
Soit $\overline{x} : \Spec(k) \to X$ un point géométrique de $X$
au-dessus de $x \in |X|$. Soit $\varphi : U \to X$ un morphisme \'etale
d'espaces algébriques et soit $u \in |U|$ tel que $\varphi(u) = x$.
Il existe alors un point géométrique
$\overline{u} : \Spec(k) \to U$ au-dessus de $u$ tel que
$\overline{x} = \varphi \circ \overline{u}$.
\end{lemma}

\begin{proof}
Choisissons un schéma affine $U'$ muni de $u' \in U'$ et un morphisme \'etale
$U' \to U$ qui envoie $u'$ sur $u$. Il suffit de prouver le lemme pour
$(U', u') \to (X, x)$. Nous pouvons donc supposer que
$U$ est un schéma, et en particulier que $U \to X$ est représentable.
Considérons alors le diagramme cartésien
$$
\xymatrix{
\Spec(k) \times_{\overline{x}, X, \varphi} U
\ar[d]_{\text{pr}_1} \ar[r]_-{\text{pr}_2} & U
\ar[d]^\varphi \\
\Spec(k) \ar[r]^-{\overline{x}} & X
}
$$
La projection $\text{pr}_1$ est le changement de base d'un morphisme \'etale ;
elle est donc \'etale, voir le
lemme \ref{lemma-base-change-etale}.
Par conséquent, le schéma $\Spec(k) \times_{\overline{x}, X, \varphi} U$
est une somme disjointe de spectres d'extensions finies séparables de $k$, voir
Morphismes, lemme \ref{morphisms-lemma-etale-over-field}.
Mais $k$ est algébriquement clos, donc toutes ces extensions sont triviales.
Ainsi, $\Spec(k) \times_{\overline{x}, X, \varphi} U$
est une somme disjointe de copies de $\Spec(k)$, et chacune
correspond à un point géométrique $\overline{u}$ tel que
$\varphi \circ \overline{u} = \overline{x}$. D'après le
lemme \ref{lemma-points-cartesian},
l'application
$$
|\Spec(k) \times_{\overline{x}, X, \varphi} U|
\longrightarrow
|\Spec(k)| \times_{|X|} |U|
$$
est surjective ; nous pouvons donc choisir $\overline{u}$ au-dessus de $u$.
\end{proof}

\begin{lemma}
\label{lemma-geometric-lift-to-cover}
Soit $S$ un schéma. Soit $X$ un espace algébrique sur $S$.
Soit $\overline{x}$ un point géométrique de $X$.
Soit $(U, \overline{u})$ un voisinage \'etale de $\overline{x}$.
Soit $\{\varphi_i : U_i \to U\}_{i \in I}$ un recouvrement \'etale dans
$X_{spaces, \etale}$.
Il existe alors $i \in I$ et $\overline{u}_i : \overline{x} \to U_i$
tels que $\varphi_i : (U_i, \overline{u}_i) \to (U, \overline{u})$
soit un morphisme de voisinages \'etales.
\end{lemma}

\begin{proof}
Soit $u \in |U|$ l'image de $\overline{u}$.
Comme $|U| = \bigcup_{i \in I} \varphi_i(|U_i|)$, il existe
$i$ et un point $u_i \in U_i$ envoyé sur $u$. Appliquons le
lemme \ref{lemma-geometric-lift-to-usual}
à $(U_i, u_i) \to (U, u)$ et à $\overline{u}$ pour
obtenir le point géométrique voulu.
\end{proof}

\begin{definition}
\label{definition-stalk}
Soit $S$ un schéma. Soit $X$ un espace algébrique sur $S$.
Soit $\mathcal{F}$ un préfaisceau sur $X_\etale$.
Soit $\overline{x}$ un point géométrique de $X$.
La {\it fibre} de $\mathcal{F}$ en $\overline{x}$ est
$$
\mathcal{F}_{\bar x}
=
\colim_{(U, \overline{u})} \mathcal{F}(U)
$$
où $(U, \overline{u})$ parcourt tous les voisinages \'etales
de $\overline{x}$ dans $X$ tels que $U \in \Ob(X_\etale)$.
\end{definition}

\noindent
D'après le
lemme \ref{lemma-cofinal-etale},
cette limite inductive est indexée par une catégorie filtrante,
à savoir l'opposée de la catégorie des voisinages \'etales
dans $X_\etale$. Plus précisément, le
lemme \ref{lemma-cofinal-etale}
dit que l'opposée de la catégorie de tous les voisinages \'etales est filtrante,
et que la sous-catégorie pleine de ceux qui appartiennent à $X_\etale$ est
cofinale, donc elle aussi filtrante.

\medskip\noindent
Ainsi, un élément de $\mathcal{F}_{\overline{x}}$ peut être
représenté par un triplet $(U, \overline{u}, \sigma)$, où
$U \in \Ob(X_\etale)$ et $\sigma \in \mathcal{F}(U)$.
Deux triplets $(U, \overline{u}, \sigma)$, $(U', \overline{u}', \sigma')$
définissent le même élément de la fibre s'il existe un troisième
voisinage \'etale
$(U'', \overline{u}'')$, $U'' \in \Ob(X_\etale)$
et des morphismes de voisinages \'etales
$h : (U'', \overline{u}'') \to (U, \overline{u})$,
$h' : (U'', \overline{u}'') \to (U', \overline{u}')$ tels que
$h^*\sigma = (h')^*\sigma'$ dans $\mathcal{F}(U'')$. Voir
Catégories, section \ref{categories-section-directed-colimits}.

\medskip\noindent
Cela implique aussi que, si $\mathcal{F}'$ est le faisceau sur
$X_{spaces, \etale}$ correspondant à $\mathcal{F}$ sur
$X_\etale$, alors
\begin{equation}
\label{equation-stalk-spaces-etale}
\mathcal{F}_{\overline{x}} = \colim_{(U, \overline{u})} \mathcal{F}'(U)
\end{equation}
où la limite inductive porte maintenant sur tous les voisinages \'etales de $\overline{x}$.
Nous passerons souvent sans autre mention du point de vue de $X_\etale$
à celui de $X_{spaces, \etale}$, et réciproquement.

\medskip\noindent
En particulier, si $\mathcal{F}$ est un préfaisceau de
groupes abéliens, d'anneaux, etc., alors $\mathcal{F}_{\overline{x}}$ est
un groupe abélien, un anneau, etc., par la construction usuelle de la
structure algébrique sur une limite inductive filtrante de tels objets.

\begin{lemma}
\label{lemma-stalk-gives-point}
\begin{slogan}
Un point géométrique d'un espace algébrique donne un point de son topos \'etale.
\end{slogan}
Soit $S$ un schéma.
Soit $X$ un espace algébrique sur $S$.
Soit $\overline{x}$ un point géométrique de $X$.
Considérons le foncteur
$$
u : X_\etale \longrightarrow \textit{Ens}, \quad
U \longmapsto |U_{\overline{x}}|
$$
Alors $u$ définit un point $p$ du site $X_\etale$
(Sites, définition \ref{sites-definition-point}),
et le foncteur fibre associé $\mathcal{F} \mapsto \mathcal{F}_p$
(Sites, équation \ref{sites-equation-stalk})
est le foncteur $\mathcal{F} \mapsto \mathcal{F}_{\overline{x}}$
défini ci-dessus.
\end{lemma}

\begin{proof}
Dans la preuve du
lemme \ref{lemma-geometric-lift-to-usual},
nous avons vu que le schéma $U_{\overline{x}} = \overline{x} \times_X U$
est une somme disjointe de
schémas isomorphes à $\overline{x}$. On peut donc aussi voir
$|U_{\overline{x}}|$ comme l'ensemble des points géométriques de $U$
au-dessus de $\overline{x}$, c'est-à-dire comme l'ensemble des morphismes
$\overline{u} : \overline{x} \to U$ qui s'insèrent dans le diagramme de la
définition \ref{definition-etale-neighbourhood}.
Il s'ensuit que $u(X)$ est un singleton et que
$u(U \times_V W) = u(U) \times_{u(V)} u(W)$
dès que $U \to V$ et $W \to V$ sont des morphismes dans $X_\etale$.
De plus, pour un recouvrement $\{U_i \to U\}_{i \in I}$ de $X_\etale$,
l'application $\coprod u(U_i) \to u(U)$ est surjective d'après le
lemme \ref{lemma-geometric-lift-to-cover}.
La
proposition \ref{sites-proposition-point-limits} de Sites
s'applique donc, si bien que $p$ est un point du site $X_\etale$.
Enfin, notre foncteur $\mathcal{F} \mapsto \mathcal{F}_{\overline{x}}$
est donné par exactement la même limite inductive que le foncteur
$\mathcal{F} \mapsto \mathcal{F}_p$ associé à $p$ dans
Sites, équation \ref{sites-equation-stalk},
ce qui prouve la dernière assertion.
\end{proof}

\begin{lemma}
\label{lemma-stalk-exact}
Soit $S$ un schéma.
Soit $X$ un espace algébrique sur $S$.
Soit $\overline{x}$ un point géométrique de $X$.
\begin{enumerate}
\item Le foncteur fibre
$\textit{PAb}(X_\etale) \to \textit{Ab}$,
$\mathcal{F}  \mapsto  \mathcal{F}_{\overline{x}}$
est exact.
\item On a $(\mathcal{F}^\#)_{\overline{x}} = \mathcal{F}_{\overline{x}}$
pour tout préfaisceau d'ensembles $\mathcal{F}$ sur $X_\etale$.
\item Le foncteur
$\textit{Ab}(X_\etale) \to \textit{Ab}$,
$\mathcal{F} \mapsto \mathcal{F}_{\overline{x}}$ est exact.
\item De même, les foncteurs
$\textit{PSh}(X_\etale) \to \textit{Ens}$ et
$\Sh(X_\etale) \to \textit{Ens}$ donnés par le foncteur fibre
$\mathcal{F} \mapsto \mathcal{F}_{\overline{x}}$ sont exacts (voir
Catégories, définition \ref{categories-definition-exact})
et commutent aux colimites quelconques.
\end{enumerate}
\end{lemma}

\begin{proof}
Ce résultat découle des résultats généraux de
Modules sur les sites, section \ref{sites-modules-section-stalks}.
En effet, $\mathcal{F} \mapsto \mathcal{F}_{\overline{x}}$
provient d'un point du petit site \'etale de $X$, voir le
lemme \ref{lemma-stalk-gives-point}. Voir la preuve de
Cohomologie \'etale, lemme \ref{etale-cohomology-lemma-stalk-exact},
pour une preuve directe de certaines de ces assertions dans le cadre
du petit site \'etale d'un schéma.
\end{proof}

\noindent
Nous verrons plus bas que le foncteur fibre
$\mathcal{F} \mapsto \mathcal{F}_{\overline{x}}$
est en fait l'image réciproque par le morphisme $\overline{x}$. En ce sens,
le lemme suivant généralise le lemme précédent.

\begin{lemma}
\label{lemma-stalk-pullback}
Soit $S$ un schéma.
Soit $f : X \to Y$ un morphisme d'espaces algébriques sur $S$.
\begin{enumerate}
\item Le foncteur
$f_{small}^{-1} :
\textit{Ab}(Y_\etale)
\to
\textit{Ab}(X_\etale)$
est exact.
\item Le foncteur
$f_{small}^{-1} :
\Sh(Y_\etale)
\to
\Sh(X_\etale)$
est exact, c'est-à-dire qu'il commute aux limites et colimites finies, voir
Catégories, définition \ref{categories-definition-exact}.
\item Pour tout morphisme \'etale $V \to Y$ d'espaces algébriques,
on a $f_{small}^{-1}h_V = h_{X \times_Y V}$.
\item Soit $\overline{x} \to X$ un point géométrique.
Soit $\mathcal{G}$ un faisceau sur $Y_\etale$.
Il existe alors une identification canonique
$$
(f_{small}^{-1}\mathcal{G})_{\overline{x}} = \mathcal{G}_{\overline{y}}.
$$
où $\overline{y} = f \circ \overline{x}$.
\end{enumerate}
\end{lemma}

\begin{proof}
Rappelons que $f_{small}$ est défini au moyen de $f_{spaces, small}$ dans le
lemme \ref{lemma-functoriality-etale-site}.
Les assertions (1), (2) et (3) découlent en général du fait que
$f_{spaces, \etale} :
X_{spaces, \etale}
\to
Y_{spaces, \etale}$
est un morphisme de sites : voir
Sites, définition \ref{sites-definition-morphism-sites},
pour (2),
Modules sur les sites, lemme \ref{sites-modules-lemma-flat-pullback-exact},
pour (1), et
Sites, lemme \ref{sites-lemma-pullback-representable-sheaf},
pour (3).

\medskip\noindent
Preuve de (4). Cette assertion est un cas particulier de
Sites, lemme \ref{sites-lemma-point-morphism-sites},
via le
lemme \ref{lemma-stalk-gives-point}.
Nous en donnons aussi une preuve directe. D'après le
lemme \ref{lemma-stalk-exact},
le foncteur fibre commute au passage au faisceau associé.
Soit $\mathcal{G}'$ le faisceau sur $Y_{spaces, \etale}$ dont la restriction
à $Y_\etale$ est $\mathcal{G}$.
Rappelons que $f_{spaces, \etale}^{-1}\mathcal{G}'$ est le faisceau
associé au préfaisceau
$$
U \longrightarrow \colim_{U \to X \times_Y V} \mathcal{G}'(V),
$$
voir
Sites, sections \ref{sites-section-continuous-functors} et
\ref{sites-section-functoriality-PSh}.
On a donc
\begin{align*}
(f_{spaces, \etale}^{-1}\mathcal{G}')_{\overline{x}}
& =
\colim_{(U, \overline{u})} f_{spaces, \etale}^{-1}\mathcal{G}'(U)
\\
& = \colim_{(U, \overline{u})}
\colim_{a : U \to X \times_Y V} \mathcal{G}'(V) \\
& = \colim_{(V, \overline{v})} \mathcal{G}'(V) \\
& = \mathcal{G}'_{\overline{y}}
\end{align*}
où, dans la troisième égalité, le couple $(U, \overline{u})$ et le morphisme
$a : U \to X \times_Y V$ correspondent au couple $(V, a \circ \overline{u})$.
Comme la fibre de $\mathcal{G}'$
(resp.\ de $f_{spaces, \etale}^{-1}\mathcal{G}'$)
coïncide avec celle de $\mathcal{G}$ (resp.\ de $f_{small}^{-1}\mathcal{G}$),
voir
l'équation (\ref{equation-stalk-spaces-etale}),
le résultat en découle.
\end{proof}

\begin{remark}
\label{remark-stalk-pullback}
Cette remarque est l'analogue de
Cohomologie \'etale, remarque \ref{etale-cohomology-remark-stalk-pullback}.
Soit $S$ un schéma.
Soit $X$ un espace algébrique sur $S$.
Soit $\overline{x} : \Spec(k) \to X$ un point géométrique de $X$.
D'après
Cohomologie \'etale,
théorème \ref{etale-cohomology-theorem-equivalence-sheaves-point},
la catégorie des faisceaux sur $\Spec(k)_\etale$ est
équivalente à la catégorie des ensembles (en associant à un faisceau ses
sections globales). Il découle alors de l'assertion (4) du
lemme \ref{lemma-stalk-pullback},
appliquée au morphisme $\overline{x}$, que le foncteur
$$
\Sh(X_\etale) \longrightarrow \textit{Ens}, \quad
\mathcal{F} \longmapsto \mathcal{F}_{\overline{x}}
$$
est isomorphe au foncteur
$$
\Sh(X_\etale)
\longrightarrow
\Sh(\Spec(k)_\etale) = \textit{Ens},
\quad
\mathcal{F} \longmapsto \overline{x}^*\mathcal{F}
$$
Nous pouvons donc considérer les foncteurs fibres comme des foncteurs
d'image réciproque par des morphismes géométriques (et non simplement par
des morphismes abstraits de topos, comme dans le résultat du
lemme \ref{lemma-stalk-gives-point}).
\end{remark}

\begin{remark}
\label{remark-map-stalks}
Soit $S$ un schéma.
Soit $X$ un espace algébrique sur $S$.
Soit $x \in |X|$.
Nous affirmons que, pour toute paire de points géométriques $\overline{x}$ et
$\overline{x}'$ au-dessus de $x$, les foncteurs fibres sont isomorphes.
Par définition de $|X|$, on peut trouver un troisième point géométrique
$\overline{x}''$ tel qu'il existe un diagramme commutatif
$$
\xymatrix{
\overline{x}'' \ar[r] \ar[d] \ar[rd]^{\overline{x}''} &
\overline{x}' \ar[d]^{\overline{x}'} \\
\overline{x} \ar[r]^{\overline{x}} & X.
}
$$
Comme le foncteur fibre $\mathcal{F} \mapsto \mathcal{F}_{\overline{x}}$
est donné par l'image réciproque par le morphisme $\overline{x}$ (et de même
pour les autres), on conclut par fonctorialité des images réciproques.
\end{remark}




\noindent
Le théorème suivant affirme que le petit site \'etale d'un espace
algébrique a assez de points.

\begin{theorem}
\label{theorem-exactness-stalks}
Soit $S$ un schéma. Soit $X$ un espace algébrique sur $S$.
Un morphisme $a : \mathcal{F} \to \mathcal{G}$ de faisceaux d'ensembles est
injectif (resp.\ surjectif) si et seulement si le morphisme sur les fibres
$a_{\overline{x}} : \mathcal{F}_{\overline{x}} \to \mathcal{G}_{\overline{x}}$
est injectif (resp.\ surjectif) pour tout point géométrique de $X$.
Une suite de faisceaux abéliens sur $X_\etale$ est exacte
si et seulement si elle est exacte sur toutes les fibres aux points géométriques de $X$.
\end{theorem}

\begin{proof}
Nous savons que le théorème est vrai lorsque $X$ est un schéma, voir
Cohomologie \'etale, théorème \ref{etale-cohomology-theorem-exactness-stalks}.
Choisissons un morphisme \'etale surjectif $f : U \to X$, où $U$ est un schéma.
Comme $\{U \to X\}$ est un recouvrement (dans $X_{spaces, \etale}$), on peut
vérifier qu'un morphisme de faisceaux est injectif ou surjectif après
restriction à $U$. Or, si $\overline{u} : \Spec(k) \to U$ est un point
géométrique de $U$, alors
$(\mathcal{F}|_U)_{\overline{u}} = \mathcal{F}_{\overline{x}}$
où $\overline{x} = f \circ \overline{u}$. (Cela résulte directement des
limites inductives qui définissent les fibres en $\overline{u}$ et $\overline{x}$, mais
cela résulte aussi du
lemme \ref{lemma-stalk-pullback}.)
Le résultat pour $U$ implique donc celui pour $X$, ce qui conclut.
\end{proof}

\noindent
Le lemme suivant peut être omis en première lecture.

\begin{lemma}
\label{lemma-points-small-etale-site}
Soit $S$ un schéma.
Soit $X$ un espace algébrique sur $S$.
Soit $p : \Sh(pt) \to \Sh(X_\etale)$
un point du petit topos \'etale de $X$.
Il existe alors un point géométrique $\overline{x}$ de $X$
tel que le foncteur fibre $\mathcal{F} \mapsto \mathcal{F}_p$
soit isomorphe au foncteur fibre
$\mathcal{F} \mapsto \mathcal{F}_{\overline{x}}$.
\end{lemma}

\begin{proof}
D'après
Sites, lemme \ref{sites-lemma-point-site-topos},
il y a une correspondance biunivoque entre les points du site et ceux
du topos associé. Nous pouvons donc supposer que $p$ est donné par
un foncteur $u : X_\etale \to \textit{Ens}$ qui définit un point
du site $X_\etale$.
Soit $U \in \Ob(X_\etale)$ un objet dont le morphisme structural
$j : U \to X$ est surjectif. Notons que $h_U$ est un faisceau
muni d'une surjection sur le faisceau final. Comme le foncteur fibre est exact,
on voit que $(h_U)_p = u(U)$ n'est pas vide (utiliser
Sites, lemme \ref{sites-lemma-points-recover}).
Choisissons $x \in u(U)$. D'après
Sites, lemme \ref{sites-lemma-point-localize},
on obtient un point $q : \Sh(pt) \to \Sh(U_\etale)$
tel que $p = j_{small} \circ q$, et donc, fonctoriellement,
$\mathcal{F}_p = (\mathcal{F}|_U)_q$.
D'après
Cohomologie \'etale, lemme \ref{etale-cohomology-lemma-points-small-etale-site},
il existe un point géométrique $\overline{u}$ de $U$ et un isomorphisme
fonctoriel $\mathcal{G}_q = \mathcal{G}_{\overline{u}}$
pour $\mathcal{G} \in \Sh(U_\etale)$. Posons
$\overline{x} = j \circ \overline{u}$. On voit alors que
$\mathcal{F}_{\overline{x}} \cong (\mathcal{F}|_U)_{\overline{u}}$
fonctoriellement en $\mathcal{F}$ sur $X_\etale$, d'après le
lemme \ref{lemma-stalk-pullback},
ce qui conclut.
\end{proof}





\section{Supports de faisceaux abéliens}
\label{section-support}

\noindent
Commençons par les supports des sections locales.

\begin{lemma}
\label{lemma-support-subsheaf-final}
Soit $S$ un schéma. Soit $X$ un espace algébrique sur $S$.
Soit $\mathcal{F}$ un sous-faisceau de l'objet final du topos \'etale
de $X$ (voir
Sites, exemple \ref{sites-example-singleton-sheaf}).
Il existe alors un unique sous-espace ouvert
$W \subset X$ tel que $\mathcal{F} = h_W$.
\end{lemma}

\begin{proof}
La condition signifie que $\mathcal{F}(U)$ est un singleton ou
est vide pour tout $\varphi : U \to X$ dans $\Ob(X_{spaces, \etale})$.
En particulier, les sections locales se recollent toujours. Si
$\mathcal{F}(U) \not = \emptyset$, alors
$\mathcal{F}(\varphi(U)) \not = \emptyset$, car
$\varphi(U) \subset X$ est un sous-espace ouvert
(lemme \ref{lemma-etale-open})
et
$\{\varphi : U \to \varphi(U)\}$ est un recouvrement dans $X_{spaces, \etale}$.
Il suffit de prendre
$W = \bigcup_{\varphi : U \to X, \mathcal{F}(U) \not = \emptyset} \varphi(U)$
pour conclure.
\end{proof}

\begin{lemma}
\label{lemma-zero-over-image}
Soit $S$ un schéma.
Soit $X$ un espace algébrique sur $S$.
Soit $\mathcal{F}$ un faisceau abélien sur $X_{spaces, \etale}$.
Soit $\sigma \in \mathcal{F}(U)$ une section locale.
Il existe un sous-espace ouvert $W \subset U$ tel que
\begin{enumerate}
\item $W \subset U$ soit le plus grand sous-espace ouvert de $U$
tel que $\sigma|_W = 0$,
\item pour tout $\varphi : V \to U$ dans $X_{spaces, \etale}$, on ait
$$
\sigma|_V = 0 \Leftrightarrow \varphi(V) \subset W,
$$
\item pour tout point géométrique $\overline{u}$ de $U$, on ait
$$
(U, \overline{u}, \sigma) = 0\text{ dans }\mathcal{F}_{\overline{x}}
\Leftrightarrow
\overline{u} \in W
$$
où $\overline{x} = (U \to X) \circ \overline{u}$.
\end{enumerate}
\end{lemma}

\begin{proof}
Comme $\mathcal{F}$ est un faisceau pour la topologie \'etale, la restriction de
$\mathcal{F}$ à $U_{Zar}$ est un faisceau sur $U$ pour la topologie de Zariski.
Il existe donc un ouvert de Zariski $W$ ayant la propriété (1), voir
Modules, lemme \ref{modules-lemma-support-section-closed}. Soit
$\varphi : V \to U$ une flèche de $X_{spaces, \etale}$. Notons que
$\varphi(V) \subset U$ est un sous-espace ouvert
(lemme \ref{lemma-etale-open})
et que $\{V \to \varphi(V)\}$ est un recouvrement \'etale. Si donc
$\sigma|_V = 0$, la condition de faisceau de $\mathcal{F}$ montre que
$\sigma|_{\varphi(V)} = 0$. Cela prouve (2).
Pour prouver (3), il faut montrer que, si $(U, \overline{u}, \sigma)$
définit l'élément nul de $\mathcal{F}_{\overline{x}}$, alors
$\overline{u} \in W$. En effet, l'hypothèse signifie
qu'il existe un morphisme de voisinages \'etales
$(V, \overline{v}) \to (U, \overline{u})$ tel que
$\sigma|_V = 0$. D'après (2), le morphisme $V \to U$ se factorise par $W$,
et donc $\overline{u} \in W$.
\end{proof}

\noindent
Soit $S$ un schéma.
Soit $X$ un espace algébrique sur $S$.
Soit $x \in |X|$.
Soit $\mathcal{F}$ un faisceau sur $X_\etale$. D'après la
remarque \ref{remark-map-stalks},
la classe d'isomorphisme de la fibre du faisceau $\mathcal{F}$
en un point géométrique au-dessus de $x$ est bien définie.

\begin{definition}
\label{definition-support}
Soit $S$ un schéma.
Soit $X$ un espace algébrique sur $S$.
Soit $\mathcal{F}$ un faisceau abélien sur $X_\etale$.
\begin{enumerate}
\item Le {\it support de $\mathcal{F}$} est l'ensemble des
points $x \in |X|$ tels que $\mathcal{F}_{\overline{x}} \not = 0$
pour tout (ou, ce qui revient au même, pour un) point géométrique $\overline{x}$ au-dessus de $x$.
\item Soit $\sigma \in \mathcal{F}(U)$ une section.
Le {\it support de $\sigma$} est la partie fermée $U \setminus W$, où
$W \subset U$ est le plus grand ouvert de $U$ sur lequel la restriction de $\sigma$
est nulle (voir le
lemme \ref{lemma-zero-over-image}).
\end{enumerate}
\end{definition}

\begin{lemma}
\label{lemma-support-section-closed}
Soit $S$ un schéma.
Soit $X$ un espace algébrique sur $S$.
Soit $\mathcal{F}$ un faisceau abélien sur $X_\etale$.
Soient $U \in \Ob(X_\etale)$ et $\sigma \in \mathcal{F}(U)$.
\begin{enumerate}
\item Le support de $\sigma$ est fermé dans $|U|$.
\item Le support de $\sigma + \sigma'$ est contenu dans la réunion des
supports de $\sigma, \sigma' \in \mathcal{F}(U)$.
\item Si $\varphi : \mathcal{F} \to \mathcal{G}$ est un morphisme de
faisceaux abéliens sur $X_\etale$, alors le support de $\varphi(\sigma)$ est
contenu dans celui de $\sigma \in \mathcal{F}(U)$.
\item Le support de $\mathcal{F}$ est la réunion des images des
supports de toutes les sections locales de $\mathcal{F}$.
\item Si $\mathcal{F} \to \mathcal{G}$ est surjectif, alors le support
de $\mathcal{G}$ est contenu dans celui de $\mathcal{F}$.
\item Si $\mathcal{F} \to \mathcal{G}$ est injectif, alors le support
de $\mathcal{F}$ est contenu dans celui de $\mathcal{G}$.
\end{enumerate}
\end{lemma}

\begin{proof}
L'assertion (1) résulte de la définition.
Les assertions (2) et (3) sont vraies parce qu'elles le sont pour les
restrictions de $\mathcal{F}$ et $\mathcal{G}$ à $U_{Zar}$, voir
Modules, lemme \ref{modules-lemma-support-section-closed}.
L'assertion (4) découle directement de l'assertion (3) du
lemme \ref{lemma-zero-over-image}.
Les assertions (5) et (6) découlent des autres.
\end{proof}

\begin{lemma}
\label{lemma-support-sheaf-rings-closed}
Le support d'un faisceau d'anneaux sur le petit site \'etale d'un
espace algébrique est fermé.
\end{lemma}

\begin{proof}
Cela résulte du fait que (selon nos conventions)
un anneau est $0$ si et seulement si
$1 = 0$ ; le support d'un faisceau d'anneaux
est donc celui de sa section unité.
\end{proof}






\section{Le faisceau structural d'un espace algébrique}
\label{section-structure-sheaf}

\noindent
Le faisceau structural d'un espace algébrique est le faisceau d'anneaux
du lemme suivant.

\begin{lemma}
\label{lemma-sheaf-condition-holds}
Soit $S$ un schéma. Soit $X$ un espace algébrique sur $S$.
La règle $U \mapsto \Gamma(U, \mathcal{O}_U)$ définit
un faisceau d'anneaux sur $X_\etale$.
\end{lemma}

\begin{proof}
Cela résulte immédiatement de la définition d'un recouvrement et de
Descente, lemme \ref{descent-lemma-sheaf-condition-holds}.
\end{proof}

\begin{definition}
\label{definition-structure-sheaf}
Soit $S$ un schéma.
Soit $X$ un espace algébrique sur $S$.
Le {\it faisceau structural} de $X$
est le faisceau d'anneaux $\mathcal{O}_X$
sur le petit site \'etale $X_\etale$ décrit dans le
lemme \ref{lemma-sheaf-condition-holds}.
\end{definition}

\noindent
D'après le lemme \ref{lemma-characterize-sheaf-small-etale}, le faisceau
$\mathcal{O}_X$ correspond à un système de faisceaux \'etales $(\mathcal{O}_X)_U$,
où $U$ parcourt les objets de $X_\etale$. La preuve de ce lemme et
notre définition montrent clairement que l'on a simplement
$(\mathcal{O}_X)_U = \mathcal{O}_U$, où $\mathcal{O}_U$ est le faisceau
structural de $U_\etale$ introduit dans
Descente, définition \ref{descent-definition-structure-sheaf}.
En particulier, si $X$ est un schéma, on retrouve le faisceau $\mathcal{O}_X$
sur le petit site \'etale de $X$.

\medskip\noindent
Au moyen de l'équivalence
$\Sh(X_\etale) = \Sh(X_{spaces, \etale})$
du lemme \ref{lemma-compare-etale-sites}, on peut également considérer $\mathcal{O}_X$
comme un faisceau d'anneaux sur $X_{spaces, \etale}$. La
remarque \ref{remark-explain-equivalence}
explique comment calculer $\mathcal{O}_X(Y)$, et en particulier $\mathcal{O}_X(X)$, lorsque
$Y \to X$ est un objet de $X_{spaces, \etale}$.

\begin{lemma}
\label{lemma-morphism-ringed-topoi}
Soit $S$ un schéma.
Soit $f : X \to Y$ un morphisme d'espaces algébriques sur $S$.
Il existe alors un morphisme canonique
$f^\sharp : f_{small}^{-1}\mathcal{O}_Y \to \mathcal{O}_X$ tel que
$$
(f_{small}, f^\sharp) :
(\Sh(X_\etale), \mathcal{O}_X)
\longrightarrow
(\Sh(Y_\etale), \mathcal{O}_Y)
$$
soit un morphisme de topos annelés. De plus,
\begin{enumerate}
\item La construction $f \mapsto (f_{small}, f^\sharp)$ est compatible aux
compositions.
\item Si $f$ est un morphisme de schémas, alors $f^\sharp$ est le morphisme décrit dans
Descente, remarque \ref{descent-remark-change-topologies-ringed}.
\end{enumerate}
\end{lemma}

\begin{proof}
D'après le lemme \ref{lemma-f-map}, il suffit de donner une $f$-application de
$\mathcal{O}_Y$ vers $\mathcal{O}_X$. Autrement dit, pour tout
diagramme commutatif
$$
\xymatrix{
U \ar[d]_g \ar[r] & X \ar[d]^f \\
V \ar[r] & Y
}
$$
où $U \in X_\etale$, $V \in Y_\etale$, il faut donner un
morphisme d'anneaux
$
(f^\sharp)_{(U, V, g)} :
\Gamma(V, \mathcal{O}_V)
\to
\Gamma(U, \mathcal{O}_U).
$
On prend bien sûr $(f^\sharp)_{(U, V, g)} = g^\sharp$.
Il est clair que ce choix est compatible aux morphismes de restriction
et donne donc bien une $f$-application.
Nous omettons la vérification de la compatibilité aux compositions et de l'accord avec la
construction de
Descente, remarque \ref{descent-remark-change-topologies-ringed}.
\end{proof}

\begin{lemma}
\label{lemma-reduced-space}
Soit $S$ un schéma. Soit $X$ un espace algébrique sur $S$.
Les conditions suivantes sont équivalentes :
\begin{enumerate}
\item $X$ est réduit,
\item pour tout $x \in |X|$, l'anneau local de $X$ en $x$ est
réduit (remarque \ref{remark-list-properties-local-ring-local-etale-topology}).
\end{enumerate}
Dans ce cas, $\Gamma(X, \mathcal{O}_X)$ est un anneau réduit et,
si $f \in \Gamma(X, \mathcal{O}_X)$ vérifie $X = V(f)$, alors $f = 0$.
\end{lemma}

\begin{proof}
L'équivalence de (1) et (2) résulte de
Propriétés, lemme \ref{properties-lemma-characterize-reduced},
appliqué aux schémas affines \'etales sur $X$. Les dernières assertions
résultent du lemme cité et du fait que $\Gamma(X, \mathcal{O}_X)$ est
un sous-anneau de $\Gamma(U, \mathcal{O}_U)$ pour un certain
schéma réduit $U$ \'etale sur $X$.
\end{proof}







\section{Fibres du faisceau structural}
\label{section-stalks-structure-sheaf}

\noindent
Cette section est l'analogue de
Cohomologie \'etale, section
\ref{etale-cohomology-section-stalks-structure-sheaf}.

\begin{lemma}
\label{lemma-describe-etale-local-ring}
Soit $S$ un schéma.
Soit $X$ un espace algébrique sur $S$.
Soit $\overline{x}$ un point géométrique de $X$.
Soit $(U, \overline{u})$ un voisinage \'etale de $\overline{x}$,
où $U$ est un schéma. On a alors
$$
\mathcal{O}_{X, \overline{x}} =
\mathcal{O}_{U, \overline{u}} =
\mathcal{O}_{U, u}^{sh}
$$
où le membre de gauche est la fibre du faisceau structural de $X$,
et celui de droite l'hensélisé strict de l'anneau local
de $U$ en la localité $u$ de $\overline{u}$.
\end{lemma}

\begin{proof}
On sait que le faisceau structural $\mathcal{O}_U$ sur
$U_\etale$ est la restriction du faisceau structural de $X$.
La première égalité résulte donc du
lemme \ref{lemma-stalk-pullback}, partie (4).
La seconde égalité est expliquée dans
Cohomologie \'etale,
lemme \ref{etale-cohomology-lemma-describe-etale-local-ring}.
\end{proof}

\begin{definition}
\label{definition-etale-local-rings}
Soit $S$ un schéma.
Soit $X$ un espace algébrique sur $S$.
Soit $\overline{x}$ un point géométrique de $X$ au-dessus du point
$x \in |X|$.
\begin{enumerate}
\item L'{\it anneau local \'etale de $X$ en $\overline{x}$}
est la fibre du faisceau structural $\mathcal{O}_X$ sur $X_\etale$
en $\overline{x}$.
Notation : $\mathcal{O}_{X, \overline{x}}$.
\item L'{\it hensélisé strict de $X$ en $\overline{x}$}
est le schéma $\Spec(\mathcal{O}_{X, \overline{x}})$.
\end{enumerate}
\end{definition}

\noindent
À isomorphisme près, l'hensélisé strict de $X$ en $\overline{x}$
(comme schéma sur $X$) ne dépend que du point $x \in |X|$, et non
du choix du point géométrique au-dessus de $x$ ; voir la
remarque \ref{remark-map-stalks}.

\begin{lemma}
\label{lemma-etale-site-locally-ringed}
Soit $S$ un schéma.
Soit $X$ un espace algébrique sur $S$.
Le petit site \'etale $X_\etale$, muni de son
faisceau structural $\mathcal{O}_X$, est un site localement annelé ; voir
Modules sur les sites, définition \ref{sites-modules-definition-locally-ringed}.
\end{lemma}

\begin{proof}
Cela résulte du fait que les fibres
$\mathcal{O}_{X, \overline{x}}$ sont
des anneaux locaux et que $X_\etale$ a assez de points ; voir le
lemme \ref{lemma-describe-etale-local-ring} et le
théorème \ref{theorem-exactness-stalks}.
Voir
Modules sur les sites, lemme \ref{sites-modules-lemma-locally-ringed-stalk} et
\ref{sites-modules-lemma-ringed-stalk-not-zero}
pour le fait qu'il en résulte que le petit site \'etale est localement annelé.
\end{proof}

\begin{lemma}
\label{lemma-dimension-local-ring}
Soit $S$ un schéma. Soit $X$ un espace algébrique sur $S$.
Soit $x \in |X|$ un point. Soit $d \in \{0, 1, 2, \ldots, \infty\}$.
Les conditions suivantes sont équivalentes :
\begin{enumerate}
\item la dimension de l'anneau local de $X$ en $x$
(définition \ref{definition-dimension-local-ring}) est $d$ ;
\item $\dim(\mathcal{O}_{X, \overline{x}}) = d$ pour un point
géométrique $\overline{x}$ au-dessus de $x$ ;
\item $\dim(\mathcal{O}_{X, \overline{x}}) = d$ pour tout point
géométrique $\overline{x}$ au-dessus de $x$.
\end{enumerate}
\end{lemma}

\begin{proof}
L'équivalence de (2) et (3) résulte du fait que le
type d'isomorphie de $\mathcal{O}_{X, \overline{x}}$ ne dépend que
de $x \in |X|$ ; voir la remarque \ref{remark-map-stalks}.
En utilisant le lemme \ref{lemma-describe-etale-local-ring},
l'équivalence de (1) et (2)$+$(3) se ramène à
l'assertion suivante : pour tout anneau local $R$, on a
$\dim(R) = \dim(R^{sh})$. C'est le contenu de
Compléments d'algèbre, lemme \ref{more-algebra-lemma-henselization-dimension}.
\end{proof}

\begin{lemma}
\label{lemma-dimension-decent-invariant-under-etale}
Soit $S$ un schéma. Soit $f : X \to Y$ un morphisme \'etale
d'espaces algébriques sur $S$. Soit $x \in |X|$. Alors
(1) $\dim_x(X) = \dim_{f(x)}(Y)$ et (2) la dimension de
l'anneau local de $X$ en $x$ est égale à la dimension de
l'anneau local de $Y$ en $f(x)$. Si $f$ est surjectif, alors
(3) $\dim(X) = \dim(Y)$.
\end{lemma}

\begin{proof}
Choisissons un schéma $U$, un point $u \in U$ et un morphisme \'etale
$U \to X$ qui envoie $u$ sur $x$. Le composé $U \to Y$
est encore \'etale et envoie $u$ sur $f(x)$. Les assertions (1) et (2)
en résultent, car les nombres considérés sont définis par le comportement
du schéma $U$ en $u$. Voir la
définition \ref{definition-dimension-at-point} pour (1). La partie (3) est
une conséquence immédiate de (1) ; voir la définition \ref{definition-dimension}.
\end{proof}

\begin{lemma}
\label{lemma-reduced-local-ring}
Soit $S$ un schéma. Soit $X$ un espace algébrique sur $S$.
Soit $x \in |X|$ un point. Les conditions suivantes sont équivalentes :
\begin{enumerate}
\item l'anneau local de $X$ en $x$ est réduit
(remarque \ref{remark-list-properties-local-ring-local-etale-topology}) ;
\item $\mathcal{O}_{X, \overline{x}}$ est réduit pour un point
géométrique $\overline{x}$ au-dessus de $x$ ;
\item $\mathcal{O}_{X, \overline{x}}$ est réduit pour tout point
géométrique $\overline{x}$ au-dessus de $x$.
\end{enumerate}
\end{lemma}

\begin{proof}
L'équivalence de (2) et (3) résulte du fait que le
type d'isomorphie de $\mathcal{O}_{X, \overline{x}}$ ne dépend que
de $x \in |X|$ ; voir la remarque \ref{remark-map-stalks}.
En utilisant le lemme \ref{lemma-describe-etale-local-ring},
l'équivalence de (1) et (2)$+$(3) se ramène à
l'assertion suivante : un anneau local est réduit si et seulement si
son hensélisé strict est réduit. C'est le contenu de
Compléments d'algèbre, lemme \ref{more-algebra-lemma-henselization-reduced}.
\end{proof}











\section{Irréductibilité locale}
\label{section-irreducible-local-ring}

\noindent
Tout point d'un espace algébrique possède un anneau local \'etale bien défini, qui
correspond au hensélisé strict de l'anneau local dans le cas d'un
schéma. En général, il ne permet pas de voir combien de composantes
irréductibles d'un schéma ou d'un espace algébrique passent par le point donné :
l'anneau local \'etale ne permet de compter que le nombre de branches géométriques.

\begin{lemma}
\label{lemma-irreducible-local-ring}
Soit $S$ un schéma.
Soit $X$ un espace algébrique sur $S$.
Soit $x \in |X|$ un point.
Les conditions suivantes sont équivalentes :
\begin{enumerate}
\item pour tout schéma $U$, tout morphisme \'etale $a : U \to X$ et
tout $u \in U$ tel que $a(u) = x$, l'anneau local $\mathcal{O}_{U, u}$ possède un
unique idéal premier minimal ;
\item pour tout schéma $U$, tout morphisme \'etale $a : U \to X$ et
tout $u \in U$ tel que $a(u) = x$, il existe une unique composante irréductible de $U$
passant par $u$ ;
\item pour tout schéma $U$, tout morphisme \'etale $a : U \to X$ et
tout $u \in U$ tel que $a(u) = x$, l'anneau local $\mathcal{O}_{U, u}$
est unibranche ;
\item pour tout schéma $U$, tout morphisme \'etale $a : U \to X$ et
tout $u \in U$ tel que $a(u) = x$, l'anneau local $\mathcal{O}_{U, u}$
est géométriquement unibranche ;
\item $\mathcal{O}_{X, \overline{x}}$ possède un unique idéal premier minimal
pour tout point géométrique $\overline{x}$ au-dessus de $x$.
\end{enumerate}
\end{lemma}

\begin{proof}
L'équivalence de (1) et (2) résulte du fait que les composantes
irréductibles de $U$ passant par $u$ sont en correspondance biunivoque ($1$-$1$) avec les
idéaux premiers minimaux de l'anneau local de $U$ en $u$. Soient $a : U \to X$ et
$u \in U$ comme dans (1). Alors $\mathcal{O}_{X, \overline{x}}$ est
l'hensélisé strict de $\mathcal{O}_{U, u}$ d'après le
lemme \ref{lemma-describe-etale-local-ring}.
En particulier, (4) et (5) sont équivalentes d'après
Compléments d'algèbre, lemme \ref{more-algebra-lemma-geometrically-unibranch}.
L'équivalence de (2), (3) et (4) résulte de
Compléments sur les morphismes, lemme \ref{more-morphisms-lemma-nr-branches}.
\end{proof}

\begin{definition}
\label{definition-unibranch}
Soit $S$ un schéma. Soit $X$ un espace algébrique sur $S$.
Soit $x \in |X|$. On dit que $X$ est {\it géométriquement unibranche
en $x$} si les conditions équivalentes du
lemme \ref{lemma-irreducible-local-ring}
sont satisfaites. On dit que $X$ est {\it géométriquement unibranche} si $X$ l'est
en tout $x \in |X|$.
\end{definition}

\noindent
Cette définition est compatible avec celle donnée pour les schémas
(Propriétés, définition \ref{properties-definition-unibranch}).

\begin{lemma}
\label{lemma-nr-branches-local-ring}
Soit $S$ un schéma. Soit $X$ un espace algébrique sur $S$.
Soit $x \in |X|$ un point. Soit $n \in \{1, 2, \ldots\}$ un entier.
Les conditions suivantes sont équivalentes :
\begin{enumerate}
\item pour tout schéma $U$, tout morphisme \'etale $a : U \to X$ et
tout $u \in U$ tel que $a(u) = x$, le nombre d'idéaux premiers minimaux
de l'anneau local $\mathcal{O}_{U, u}$ est $\leq n$,
et, pour au moins un choix de $U, a, u$, il vaut $n$ ;
\item pour tout schéma $U$, tout morphisme \'etale $a : U \to X$ et
tout $u \in U$ tel que $a(u) = x$, le nombre de composantes irréductibles de
$U$ passant par $u$ est $\leq n$ et, pour au moins un choix
de $U, a, u$, il vaut $n$ ;
\item pour tout schéma $U$, tout morphisme \'etale $a : U \to X$ et tout $u \in U$
tel que $a(u) = x$, le nombre de branches de $U$ en $u$ est $\leq n$,
et, pour au moins un choix de $U, a, u$, il vaut $n$ ;
\item pour tout schéma $U$, tout morphisme \'etale $a : U \to X$ et tout $u \in U$
tel que $a(u) = x$, le nombre de branches géométriques de $U$ en $u$ est $n$ ;
\item le nombre d'idéaux premiers minimaux de
$\mathcal{O}_{X, \overline{x}}$ est $n$.
\end{enumerate}
\end{lemma}

\begin{proof}
L'équivalence de (1) et (2) résulte du fait que les composantes
irréductibles de $U$ passant par $u$ sont en correspondance biunivoque ($1$-$1$) avec les
idéaux premiers minimaux de l'anneau local de $U$ en $u$. Soient $a : U \to X$ et
$u \in U$ comme dans (1). Alors $\mathcal{O}_{X, \overline{x}}$ est
l'hensélisé strict de $\mathcal{O}_{U, u}$ d'après le
lemme \ref{lemma-describe-etale-local-ring}. Rappelons que
le nombre (géométrique) de branches de $U$ en $u$ est le nombre
d'idéaux premiers minimaux de l'hensélisé (strict) de $\mathcal{O}_{U, u}$.
En particulier, (4) et (5) sont équivalentes.
L'équivalence de (2), (3) et (4) résulte de
Compléments sur les morphismes, lemme \ref{more-morphisms-lemma-nr-branches}.
\end{proof}

\begin{definition}
\label{definition-number-of-geometric-branches}
Soit $S$ un schéma.
Soit $X$ un espace algébrique sur $S$.
Soit $x \in |X|$. Le {\it nombre de branches géométriques de $X$ en $x$} est
égal à $n \in \mathbf{N}$ si les conditions équivalentes du
lemme \ref{lemma-nr-branches-local-ring}
sont satisfaites, et à $\infty$ sinon.
\end{definition}








\section{Espaces noethériens}
\label{section-noetherian}

\noindent
Nous avons déjà défini les espaces algébriques localement noethériens dans la
section \ref{section-types-properties}.

\begin{definition}
\label{definition-noetherian}
Soit $S$ un schéma. Soit $X$ un espace algébrique sur $S$.
On dit que $X$ est {\it noethérien} si $X$ est quasi-compact, quasi-séparé
et localement noethérien.
\end{definition}

\noindent
Notons qu'un espace algébrique noethérien $X$ n'est pas seulement quasi-compact
et localement noethérien, mais aussi quasi-séparé. Cela ne contredit pas
la définition d'un schéma noethérien, car un schéma localement noethérien
est quasi-séparé ; voir
Propriétés, lemme \ref{properties-lemma-locally-Noetherian-quasi-separated}.
Cette propriété n'est pas vraie pour les espaces algébriques. En effet,
$X = \mathbf{A}^1_k/\mathbf{Z}$, voir
Espaces, exemple \ref{spaces-example-affine-line-translation},
est localement noethérien et quasi-compact, mais non quasi-séparé
(donc non noethérien selon nos définitions).

\medskip\noindent
Une conséquence du choix fait ci-dessus est qu'un espace algébrique
de type fini sur un espace algébrique noethérien n'est pas automatiquement
noethérien ; autrement dit, l'analogue de
Morphismes, lemme \ref{morphisms-lemma-finite-type-noetherian},
n'est pas vrai. L'énoncé correct est qu'un espace algébrique de
présentation finie sur un espace algébrique noethérien est noethérien
(voir
Morphismes d'espaces,
lemme \ref{spaces-morphisms-lemma-finite-presentation-noetherian}).

\medskip\noindent
Un espace algébrique noethérien $X$ est très proche d'être un schéma.
Dans la suite de cette section, nous rassemblons quelques lemmes qui l'illustrent.

\begin{lemma}
\label{lemma-Noetherian-topology}
Soit $S$ un schéma. Soit $X$ un espace algébrique sur $S$.
\begin{enumerate}
\item Si $X$ est localement noethérien, alors $|X|$ est un espace topologique
localement noethérien.
\item Si $X$ est quasi-compact et localement noethérien, alors $|X|$
est un espace topologique noethérien.
\end{enumerate}
\end{lemma}

\begin{proof}
Supposons $X$ localement noethérien.
Choisissons un schéma $U$ et un morphisme \'etale surjectif
$U \to X$. Puisque $X$ est localement noethérien, $U$ l'est aussi.
D'après
Propriétés, lemme \ref{properties-lemma-Noetherian-topology},
cela signifie que $|U|$ est un espace topologique localement noethérien.
Comme $|U| \to |X|$ est ouverte et surjective, on en conclut que
$|X|$ est localement noethérien d'après
Topologie, lemme \ref{topology-lemma-image-Noetherian}.
Cela prouve (1). Si $X$ est quasi-compact et localement noethérien,
alors $|X|$ est quasi-compact et localement noethérien. Ainsi, $|X|$
est noethérien d'après
Topologie,
lemme \ref{topology-lemma-quasi-compact-locally-Noetherian-Noetherian}.
\end{proof}

\begin{lemma}
\label{lemma-Noetherian-sober}
Soit $S$ un schéma. Soit $X$ un espace algébrique sur $S$.
Si $X$ est noethérien, alors $|X|$ est un espace topologique noethérien sobre.
\end{lemma}

\begin{proof}
L'espace topologique sous-jacent à un espace algébrique quasi-séparé
est sobre ; voir le
lemme \ref{lemma-quasi-separated-sober}.
Il est noethérien d'après le
lemme \ref{lemma-Noetherian-topology}.
\end{proof}

\begin{lemma}
\label{lemma-Noetherian-local-ring-Noetherian}
Soit $S$ un schéma. Soit $X$ un espace algébrique noethérien sur $S$.
Soit $\overline{x}$ un point géométrique de $X$. Alors
$\mathcal{O}_{X, \overline{x}}$ est un anneau local noethérien.
\end{lemma}

\begin{proof}
Choisissons un voisinage \'etale $(U, \overline{u})$ de $\overline{x}$,
où $U$ est un schéma. Alors $\mathcal{O}_{X, \overline{x}}$ est
l'hensélisé strict de l'anneau local de $U$ en $u$ ; voir le
lemme \ref{lemma-describe-etale-local-ring}.
D'après notre définition des espaces noethériens, le schéma $U$ est localement noethérien.
On conclut donc par
Compléments d'algèbre, lemme \ref{more-algebra-lemma-henselization-noetherian}.
\end{proof}








\section{Espaces algébriques réguliers}
\label{section-regular}

\noindent
Nous avons déjà défini les espaces algébriques réguliers à la
section \ref{section-types-properties}.

\begin{lemma}
\label{lemma-regular}
Soit $S$ un schéma.
Soit $X$ un espace algébrique localement noethérien sur $S$.
Les conditions suivantes sont équivalentes :
\begin{enumerate}
\item $X$ est régulier ;
\item tout anneau local \'etale $\mathcal{O}_{X, \overline{x}}$ est
régulier.
\end{enumerate}
\end{lemma}

\begin{proof}
Soit $U$ un schéma et soit $U \to X$ un morphisme \'etale surjectif.
Par hypothèse, $U$ est localement noethérien. De plus, tout anneau local \'etale
$\mathcal{O}_{X, \overline{x}}$ est le hensélisé strict d'un
anneau local de $U$, et réciproquement ; voir le
lemme \ref{lemma-describe-etale-local-ring}.
Ainsi, d'après le
lemme \ref{more-algebra-lemma-henselization-regular} de Compléments d'algèbre,
nous voyons que (2) équivaut à ce que tout anneau local de $U$ soit
régulier, c'est-à-dire à ce que $U$ soit un schéma régulier (voir le
lemme \ref{properties-lemma-characterize-regular} de Propriétés).
Cela équivaut à (1) d'après la
définition \ref{definition-type-property}.
\end{proof}

\noindent
Nous pouvons utiliser le lemme \ref{descent-lemma-regular-local-ring-local}
de Descente pour définir ce que signifie, pour un espace algébrique $X$, être
régulier en un point $x$.

\begin{definition}
\label{definition-regular-at-point}
Soit $S$ un schéma. Soit $X$ un espace algébrique sur $S$.
Soit $x \in |X|$ un point. Nous disons que {\it $X$ est régulier en $x$}
si $\mathcal{O}_{U, u}$ est un anneau local régulier pour tout
(ou, de manière équivalente, pour un certain) couple $(a : U \to X, u)$ formé d'un
morphisme \'etale $a : U \to X$ d'un schéma vers $X$ et d'un point
$u \in U$ tel que $a(u) = x$.
\end{definition}

\noindent
Voir la définition \ref{definition-property-at-point}, le
lemme \ref{lemma-local-source-target-at-point} et le
lemme \ref{descent-lemma-regular-local-ring-local} de Descente.

\begin{lemma}
\label{lemma-regular-at-x}
Soit $S$ un schéma. Soit $X$ un espace algébrique sur $S$.
Soit $x \in |X|$ un point. Les conditions suivantes sont équivalentes :
\begin{enumerate}
\item $X$ est régulier en $x$ ;
\item l'anneau local \'etale $\mathcal{O}_{X, \overline{x}}$ est
régulier pour tout point géométrique $\overline{x}$ au-dessus du point
$x$ (ou, de manière équivalente, pour un certain).
\end{enumerate}
\end{lemma}

\begin{proof}
Soient $U$ un schéma, $u \in U$ un point, et $a : U \to X$ un
morphisme \'etale envoyant $u$ sur $x$. Pour tout point géométrique
$\overline{x}$ de $X$ au-dessus de $x$, l'anneau local \'etale
$\mathcal{O}_{X, \overline{x}}$ est le hensélisé strict de
l'anneau local de $U$ en $u$ ; voir le
lemme \ref{lemma-describe-etale-local-ring}.
On conclut donc par le
lemme \ref{more-algebra-lemma-henselization-regular} de Compléments d'algèbre.
\end{proof}

\begin{lemma}
\label{lemma-regular-normal}
Tout espace algébrique régulier est normal.
\end{lemma}

\begin{proof}
Cela résulte des définitions et du cas des schémas.
Voir le lemme \ref{properties-lemma-regular-normal} de Propriétés.
\end{proof}

 








\section{Faisceaux de modules sur les espaces algébriques}
\label{section-modules}

\noindent
Si $X$ est un espace algébrique, un faisceau de modules sur $X$ est
un faisceau de $\mathcal{O}_X$-modules sur le petit site \'etale de $X$,
où $\mathcal{O}_X$ est le faisceau structural de $X$. La catégorie
des faisceaux de modules est notée $\textit{Mod}(\mathcal{O}_X)$.

\medskip\noindent
Étant donné un morphisme $f : X \to Y$ d'espaces algébriques, le
lemme \ref{lemma-morphism-ringed-topoi}
nous donne un morphisme de topos annelés et, par conséquent, la
définition \ref{sites-modules-definition-pushforward} de Modules sur les sites
nous donne des foncteurs image inverse et image directe bien définis
\begin{equation}
\label{equation-push-pull}
f^* :
\textit{Mod}(\mathcal{O}_Y)
\longrightarrow
\textit{Mod}(\mathcal{O}_X), \quad
f_* :
\textit{Mod}(\mathcal{O}_X)
\longrightarrow
\textit{Mod}(\mathcal{O}_Y)
\end{equation}
qui sont adjoints de la manière habituelle. Si $g : Y \to Z$ est un autre morphisme
d'espaces algébriques sur $S$, alors nous avons
$(g \circ f)^* = f^* \circ g^*$ et $(g \circ f)_* = g_* \circ f_*$
simplement parce que les morphismes de topos annelés se composent de la manière
correspondante (d'après le lemme).

\begin{lemma}
\label{lemma-etale-exact-pullback}
Soit $S$ un schéma.
Soit $f : X \to Y$ un morphisme \'etale d'espaces algébriques sur $S$.
Alors $f^{-1}\mathcal{O}_Y = \mathcal{O}_X$, et
$f^*\mathcal{G} = f_{small}^{-1}\mathcal{G}$ pour tout faisceau de
$\mathcal{O}_Y$-modules $\mathcal{G}$. En particulier,
$f^* : \textit{Mod}(\mathcal{O}_Y) \to \textit{Mod}(\mathcal{O}_X)$
est exact.
\end{lemma}

\begin{proof}
D'après la description de l'image réciproque dans le lemme \ref{lemma-etale-morphism-topoi}
et la définition des faisceaux structuraux, il est clair que
$f_{small}^{-1}\mathcal{O}_Y = \mathcal{O}_X$. Puisque, par définition, l'image inverse est
$$
f^*\mathcal{G} =
f_{small}^{-1}\mathcal{G} \otimes_{f_{small}^{-1}\mathcal{O}_Y}
\mathcal{O}_X
$$
nous concluons que $f^*\mathcal{G} = f_{small}^{-1}\mathcal{G}$.
L'exactitude est claire, car $f_{small}^{-1}$ est exact puisque $f_{small}$
est un morphisme de topos.
\end{proof}

\noindent
Nous poursuivons l'abus de notation introduit dans
l'équation (\ref{equation-restrict})
en écrivant
\begin{equation}
\label{equation-restrict-modules}
\mathcal{G}|_{X_\etale}
= f^*\mathcal{G}
= f_{small}^{-1}\mathcal{G}
\end{equation}
dans la situation du lemme ci-dessus. Nous en discuterons de manière plus
technique à la
section \ref{section-localize}.

\begin{lemma}
\label{lemma-pushforward-etale-base-change-modules}
Soit $S$ un schéma. Soit
$$
\xymatrix{
X' \ar[r] \ar[d]_{f'} & X \ar[d]^f \\
Y' \ar[r]^g & Y
}
$$
un carré cartésien d'espaces algébriques sur $S$. Soit
$\mathcal{F} \in \textit{Mod}(\mathcal{O}_X)$. Si $g$ est \'etale, alors
$f'_*(\mathcal{F}|_{X'}) = (f_*\mathcal{F})|_{Y'}$\footnote{Nous avons aussi
$(f')^*(\mathcal{G}|_{Y'}) = (f^*\mathcal{G})|_{X'}$
par commutativité du diagramme et par (\ref{equation-restrict-modules})} et
$R^if'_*(\mathcal{F}|_{X'}) = (R^if_*\mathcal{F})|_{Y'}$ dans
$\textit{Mod}(\mathcal{O}_{Y'})$.
\end{lemma}

\begin{proof}
C'est une reformulation du
lemme \ref{lemma-pushforward-etale-base-change}
dans le cas des modules.
\end{proof}

\begin{lemma}
\label{lemma-characterize-module-small-etale}
Soit $S$ un schéma. Soit $X$ un espace algébrique sur $S$.
Un faisceau $\mathcal{F}$ de $\mathcal{O}_X$-modules est donné par les
données suivantes :
\begin{enumerate}
\item pour tout $U \in \Ob(X_\etale)$, un faisceau
$\mathcal{F}_U$ de $\mathcal{O}_U$-modules sur $U_\etale$ ;
\item pour tout $f : U' \to U$ dans $X_\etale$, un isomorphisme
$c_f : f_{small}^*\mathcal{F}_U \to \mathcal{F}_{U'}$.
\end{enumerate}
Ces données sont soumises à la condition suivante : pour tous $f : U' \to U$
et $g : U'' \to U'$ dans $X_\etale$, le composé
$c_g \circ g_{small}^*c_f$ est égal à $c_{f \circ g}$.
\end{lemma}

\begin{proof}
Combiner les lemmes \ref{lemma-etale-exact-pullback}
et \ref{lemma-characterize-sheaf-small-etale}, et utiliser le fait que
tout morphisme entre objets de $X_\etale$ est un morphisme \'etale
de schémas.
\end{proof}







\section{Localisation \'etale}
\label{section-localize}

\noindent
Il faut éviter à tout prix de lire cette section.

\medskip\noindent
Soit $X \to Y$ un morphisme \'etale d'espaces algébriques.
Alors $X$ est un objet de $Y_{spaces, \etale}$ et il résulte
immédiatement des définitions, voir aussi la preuve du
lemme \ref{lemma-etale-morphism-topoi},
que
\begin{equation}
\label{equation-localize}
X_{spaces, \etale} = Y_{spaces, \etale}/X
\end{equation}
où le membre de droite est la localisation du site
$Y_{spaces, \etale}$ en l'objet $X$ ; voir la
définition \ref{sites-definition-localize} de Sites.
De plus, cette identification est compatible avec les faisceaux structuraux d'après le
lemme \ref{lemma-etale-exact-pullback}.
Par conséquent, le site annelé $(X_{spaces, \etale}, \mathcal{O}_X)$
s'identifie à la localisation du site annelé
$(Y_{spaces, \etale}, \mathcal{O}_Y)$ en l'objet $X$ :
\begin{equation}
\label{equation-localize-ringed}
(X_{spaces, \etale}, \mathcal{O}_X) =
(Y_{spaces, \etale}/X, \mathcal{O}_Y|_{Y_{spaces, \etale}/X})
\end{equation}
La localisation d'un site annelé utilisée dans le membre de droite est définie dans
Modules sur les sites, à la
définition \ref{sites-modules-definition-localize-ringed-site}.

\medskip\noindent
Supposons maintenant que $X \to Y$ soit un morphisme \'etale d'espaces algébriques et que $X$ soit
un schéma. Alors $X$ est un objet de $Y_\etale$ et il s'ensuit que
\begin{equation}
\label{equation-localize-at-scheme}
X_\etale = Y_\etale/X
\end{equation}
et
\begin{equation}
\label{equation-localize-at-scheme-ringed}
(X_\etale, \mathcal{O}_X) =
(Y_\etale/X, \mathcal{O}_Y|_{Y_\etale/X})
\end{equation}
comme ci-dessus.

\medskip\noindent
Enfin, si $X \to Y$ est un morphisme \'etale d'espaces algébriques et si $X$ est
un schéma affine, alors $X$ est un objet de $Y_{affine, \etale}$ et
\begin{equation}
\label{equation-localize-at-affine}
X_{affine, \etale} = Y_{affine, \etale}/X
\end{equation}
et
\begin{equation}
\label{equation-localize-at-affine-ringed}
(X_{affine, \etale}, \mathcal{O}_X) =
(Y_{affine, \etale}/X, \mathcal{O}_Y|_{Y_{affine, \etale}/X})
\end{equation}
comme ci-dessus.

\medskip\noindent
Montrons ensuite que ces localisations sont compatibles avec les morphismes.

\begin{lemma}
\label{lemma-relocalize-morphism}
Soit $S$ un schéma. Soit
$$
\xymatrix{
U \ar[d]_p \ar[r]_g & V \ar[d]^q \\
X \ar[r]^f & Y
}
$$
un diagramme commutatif d'espaces algébriques sur $S$ dans lequel $p$ et $q$ sont \'etales.
Par les identifications
(\ref{equation-localize-ringed}) pour $U \to X$ et $V \to Y$,
le morphisme de topos annelés
$$
(g_{spaces, \etale}, g^\sharp) :
(\Sh(U_{spaces, \etale}), \mathcal{O}_U)
\longrightarrow
(\Sh(V_{spaces, \etale}), \mathcal{O}_V)
$$
est $2$-isomorphe au morphisme $(f_{spaces, \etale, c}, f_c^\sharp)$
construit dans
Modules sur les sites, au
lemme \ref{sites-modules-lemma-relocalize-morphism-ringed-sites},
à partir du morphisme de sites annelés
$(f_{spaces, \etale}, f^\sharp)$ et de
l'application $c : U \to V \times_Y X$ correspondant à $g$.
\end{lemma}

\begin{proof}
Le morphisme $(f_{spaces, \etale, c}, f_c^\sharp)$ est défini comme le
composé $f' \circ j$
d'un morphisme de localisation et d'une application de changement de base. De même, $g$ est le composé
$U \to V \times_Y X \to V$. Il suffit donc de démontrer
le lemme dans les deux cas suivants : (1) $f = \text{id}$, et
(2) $U = X \times_Y V$. Dans le cas (1), le morphisme $g : U \to V$ est
\'etale ; voir le
lemme \ref{lemma-etale-permanence}.
Par conséquent, $(g_{spaces, \etale}, g^\sharp)$ est un morphisme de localisation
d'après la discussion autour des
équations (\ref{equation-localize}) et
(\ref{equation-localize-ringed}),
ce qui est exactement l'assertion du lemme dans ce cas.
Dans le cas (2), le morphisme $g_{spaces, \etale}$
provient du morphisme de sites annelés donné par le foncteur
$V_{spaces, \etale} \to U_{spaces, \etale}$,
$V'/V \mapsto V' \times_V U/U$
qui est aussi celui par lequel le morphisme $f'$ est défini ; voir le
lemme \ref{sites-lemma-localize-morphism} de Sites.
Nous omettons la vérification de $(f')^\sharp = g^\sharp$
dans ce cas (tous deux sont la restriction de $f^\sharp$
à $U_{spaces, \etale}$).
\end{proof}

\begin{lemma}
\label{lemma-relocalize-morphism-at-schemes}
Mêmes notations et hypothèses que dans le
lemme \ref{lemma-relocalize-morphism},
sauf que nous supposons aussi que $U$ et $V$ sont des schémas.
Par les identifications
(\ref{equation-localize-at-scheme-ringed})
pour $U \to X$ et $V \to Y$, le morphisme de topos annelés
$$
(g_{small}, g^\sharp) :
(\Sh(U_\etale), \mathcal{O}_U)
\longrightarrow
(\Sh(V_\etale), \mathcal{O}_V)
$$
est $2$-isomorphe au morphisme $(f_{small, s}, f_s^\sharp)$
construit dans
Modules sur les sites, au
lemme \ref{sites-modules-lemma-relocalize-morphism-ringed-topoi},
à partir de $(f_{small}, f^\sharp)$ et de
l'application $s : h_U \to f_{small}^{-1}h_V$ correspondant à $g$.
\end{lemma}

\begin{proof}
Notons que $(g_{small}, g^\sharp)$ est $2$-isomorphe, comme
morphisme de topos annelés, au morphisme de topos annelés
associé au morphisme de sites annelés
$(g_{spaces, \etale}, g^\sharp)$. Nous concluons donc par le
lemme \ref{lemma-relocalize-morphism}
et par le résultat de
Modules sur les sites donné au
lemme \ref{sites-modules-lemma-relocalize-morphism-compare}.
\end{proof}

\noindent
Enfin, nous examinons la relation entre les faisceaux d'ensembles sur le petit
site \'etale $Y_\etale$ d'un espace algébrique $Y$ et les espaces algébriques
\'etales sur $Y$. Soit $S$ un schéma et soit $Y$ un espace algébrique
sur $S$. Soit $\mathcal{F}$ un objet de $\Sh(Y_\etale)$. Considérons
le foncteur
$$
X : (\Sch/S)_{fppf}^{opp} \longrightarrow \textit{Ens}
$$
défini par la règle
$$
X(T) = \{(y, s) \mid y : T \to Y\text{ est un morphisme sur }S\text{ et }
s \in \Gamma(T, y_{small}^{-1}\mathcal{F})\}
$$
Étant donné un morphisme $g : T' \to T$, l'application de restriction envoie
$(y, s)$ sur $(y \circ g, g_{small}^{-1}s)$. Cela a un sens
puisque $y_{small} \circ g_{small} = (y \circ g)_{small}$ d'après le
lemme \ref{lemma-functoriality-etale-site}. Il existe une application canonique
$X \to Y$ qui envoie le couple $(y, s)$ sur $y$.

\begin{lemma}
\label{lemma-sheaf-gives-space}
Soit $S$ un schéma et soit $Y$ un espace algébrique sur $S$.
Soit $\mathcal{F}$ un faisceau d'ensembles sur $Y_\etale$.
Sous réserve d'une condition ensembliste (voir la preuve), nous avons :
\begin{enumerate}
\item le foncteur $X$ associé ci-dessus à $\mathcal{F}$ est un
espace algébrique ;
\item l'application $X \to Y$ est un morphisme \'etale d'espaces algébriques ;
\item par l'identification $\Sh(Y_\etale) = \Sh(Y_{spaces, \etale})$,
nous avons $\mathcal{F} \cong h_X$ ;
\item nous avons $\mathcal{F} \cong f_{small, !}*$. Ici, $*$ est l'objet final
de la catégorie $\Sh(X_\etale)$ et $f_{small, !}$ existe d'après le
lemme \ref{lemma-etale-morphism-topoi}.
\end{enumerate}
\end{lemma}

\begin{proof}
Montrons que $X$ est un faisceau pour la topologie fppf. Supposons donc
que $\{g_i : T_i \to T\}$ soit un recouvrement de $(\Sch/S)_{fppf}$ et que
les $(y_i, s_i) \in X(T_i)$ vérifient la condition de recollement, c'est-à-dire que les
restrictions de $(y_i, s_i)$ et de $(y_j, s_j)$ à $T_i \times_T T_j$ coïncident.
Puisque $Y$ est un faisceau pour la topologie fppf, nous voyons alors que
les $y_i$ donnent naissance à un unique morphisme $y : T \to Y$ tel que
$y_i = y \circ g_i$. Nous voyons ensuite que
$y_{i, small}^{-1}\mathcal{F} = g_{i, small}^{-1}y_{small}^{-1}\mathcal{F}$.
Les sections $s_i$ se recollent donc de manière unique en une section de
$y_{small}^{-1}\mathcal{F}$ d'après le résultat donné dans
Cohomologie \'etale, lemme \ref{etale-cohomology-lemma-describe-pullback}.

\medskip\noindent
La construction qui envoie $\mathcal{F} \in \Ob(\Sh(Y_\etale))$
sur $X \in \Ob(\Sh((\Sch/S)_{fppf}))$ commute aux limites finies et
aux colimites quelconques, puisque chacun des foncteurs $y_{small}^{-1}$ possède
ces propriétés. Bien entendu, si $V \in \Ob(Y_\etale)$, alors
la construction envoie le faisceau représentable $h_V$ de $Y_\etale$
sur le foncteur représentable représenté par $V$.

\medskip\noindent
D'après le lemme \ref{sites-lemma-sheaf-coequalizer-representable} de Sites,
nous pouvons trouver un ensemble $I$ et, pour chaque $i \in I$, un objet $V_i$ de $Y_\etale$
et une surjection de faisceaux
$$
\coprod h_{V_i} \longrightarrow \mathcal{F}
$$
sur $Y_\etale$. La condition ensembliste dont nous avons besoin est que l'ensemble
d'indices $I$ ne soit pas trop grand\footnote{Il suffit que la borne supérieure des
cardinalités des fibres de $\mathcal{F}$ aux points géométriques de $Y$
soit majorée par la taille d'un certain objet de $(\Sch/S)_{fppf}$.}.
Alors $V = \coprod V_i$ est un objet de $(\Sch/S)_{fppf}$ et
donc un objet de $Y_\etale$, et nous avons une surjection
$h_V \to \mathcal{F}$.

\medskip\noindent
Observons que le produit de $h_V$ par lui-même dans $\Sh(Y_\etale)$
est $h_{V \times_Y V}$. Considérons le produit fibré
$$
h_V \times_\mathcal{F} h_V \subset h_{V \times_Y V}
$$
Il existe un sous-schéma ouvert $R$ de $V \times_Y V$ tel que
$h_V \times_\mathcal{F} h_V = h_R$ ; voir le
lemme \ref{lemma-support-subsheaf-final} (un petit détail est omis).
Par le lemme de Yoneda, nous obtenons deux morphismes $s, t : R \to V$ dans $Y_\etale$
et nous trouvons un diagramme de conoyau du couple
$$
\xymatrix{
h_R \ar@<1ex>[r] \ar@<-1ex>[r] &
h_V \ar[r] &
\mathcal{F}
}
$$
dans $\Sh(Y_\etale)$. Bien entendu, les morphismes $s, t$ sont \'etales
et définissent une relation d'équivalence \'etale $(t, s) : R \to V \times_S V$.

\medskip\noindent
La discussion des deux paragraphes précédents nous donne un
diagramme de conoyau du couple
$$
\xymatrix{
R \ar@<1ex>[r] \ar@<-1ex>[r] &
V \ar[r] &
X
}
$$
dans $(\Sch/S)_{fppf}$. Ainsi, $X = V/R$ est un espace algébrique d'après le
théorème \ref{spaces-theorem-presentation} d'Espaces. Cela démontre (1).
L'assertion (2) résulte de ce que $V \to Y$ est \'etale. L'assertion (3) découle immédiatement
de la définition de $X$ et de $h_X$. Nous omettons la preuve de l'assertion (4) ;
elle résulte de l'identification du morphisme associé au foncteur cocontinu
$j$ du lemme \ref{lemma-etale-morphism-topoi}
avec la description de $X_{spaces, \etale}$ comme localisation
de $Y_{spaces, \etale}$ en $X$ donnée ci-dessus, et du
lemme \ref{sites-lemma-describe-j-shriek-representable} de Sites.
\end{proof}







\section{Reconstruction des morphismes}
\label{section-morphisms}

\noindent
Dans cette section, nous démontrons que le procédé qui associe à un espace algébrique
son petit topos \'etale localement annelé est pleinement fidèle en un sens
approprié ; voir le
théorème \ref{theorem-fully-faithful}.

\begin{lemma}
\label{lemma-morphism-locally-ringed}
Soit $S$ un schéma.
Soit $f : X \to Y$ un morphisme d'espaces algébriques sur $S$.
Le morphisme de topos annelés $(f_{small}, f^\sharp)$
associé à $f$ est un morphisme de topos localement annelés ; voir
Modules sur les sites,
la définition \ref{sites-modules-definition-morphism-locally-ringed-topoi}.
\end{lemma}

\begin{proof}
Remarquons que l'assertion a un sens puisque nous avons vu que
$(X_\etale, \mathcal{O}_{X_\etale})$ et
$(Y_\etale, \mathcal{O}_{Y_\etale})$
sont des sites localement annelés ; voir le
lemme \ref{lemma-etale-site-locally-ringed}.
De plus, nous savons que $X_\etale$ a assez de points ; voir le
théorème \ref{theorem-exactness-stalks}.
Il suffit donc de démontrer que $(f_{small}, f^\sharp)$
satisfait la condition (3) du
Modules sur les sites,
lemme \ref{sites-modules-lemma-locally-ringed-morphism}.
Pour le voir, prenons un point $p$ de $X_\etale$. D'après le
lemme \ref{lemma-points-small-etale-site},
$p$ correspond à un point géométrique $\overline{x}$ de $X$.
D'après le
lemme \ref{lemma-stalk-pullback},
le point $q = f_{small} \circ p$ correspond au
point géométrique $\overline{y} = f \circ \overline{x}$ de $Y$.
L'assertion à démontrer est donc que l'homomorphisme induit
d'anneaux locaux \'etales
$$
\mathcal{O}_{Y, \overline{y}} \longrightarrow \mathcal{O}_{X, \overline{x}}
$$
est un homomorphisme local d'anneaux locaux. On peut le démontrer directement, mais nous le déduisons plutôt
du résultat correspondant pour les schémas. Choisissons à cette fin un
diagramme commutatif
$$
\xymatrix{
U \ar[d] \ar[r]_\psi & V \ar[d] \\
X \ar[r] & Y
}
$$
où $U$ et $V$ sont des schémas et les flèches verticales sont surjectives
et \'etales (voir le
lemme \ref{spaces-lemma-lift-morphism-presentations} d'Espaces).
Choisissons un relèvement $\overline{u} : \overline{x} \to U$ (ce qui est possible d'après le
lemme \ref{lemma-geometric-lift-to-cover}).
Posons $\overline{v} = \psi \circ \overline{u}$. Nous obtenons un
diagramme commutatif d'anneaux locaux \'etales
$$
\xymatrix{
\mathcal{O}_{U, \overline{u}} &
\mathcal{O}_{V, \overline{v}} \ar[l] \\
\mathcal{O}_{X, \overline{x}} \ar[u] &
\mathcal{O}_{Y, \overline{y}}. \ar[l] \ar[u]
}
$$
D'après
Cohomologie \'etale, le lemme \ref{etale-cohomology-lemma-morphism-locally-ringed},
la flèche horizontale supérieure est un homomorphisme local d'anneaux locaux. Enfin, d'après le
lemme \ref{lemma-describe-etale-local-ring},
les flèches verticales sont des isomorphismes. Cela conclut.
\end{proof}

\begin{lemma}
\label{lemma-2-morphism}
Soit $S$ un schéma.
Soient $X$, $Y$ des espaces algébriques sur $S$.
Soit $f : X \to Y$ un morphisme d'espaces algébriques sur $S$.
Soit $t$ un $2$-morphisme de $(f_{small}, f^\sharp)$ vers lui-même ; voir
Modules sur les sites,
la définition \ref{sites-modules-definition-2-morphism-ringed-topoi}.
Alors $t = \text{id}$.
\end{lemma}

\begin{proof}
Soient $X'$, resp.\ $Y'$, les espaces $X$, resp.\ $Y$, considérés comme espaces algébriques sur
$\Spec(\mathbf{Z})$ ; voir la
définition \ref{spaces-definition-base-change} d'Espaces.
Il ressort de la construction que $(X_{small}, \mathcal{O})$
est égal à $(X'_{small}, \mathcal{O})$, et de même pour $Y$.
Nous pouvons donc travailler avec $X'$ et $Y'$. Autrement dit, nous pouvons
supposer que $S = \Spec(\mathbf{Z})$.

\medskip\noindent
Supposons que $S = \Spec(\mathbf{Z})$ et que $f : X \to Y$ et $t$ soient comme dans
le lemme. Autrement dit, $t : f^{-1}_{small} \to f^{-1}_{small}$
est une transformation de foncteurs telle que le diagramme
$$
\xymatrix{
f_{small}^{-1}\mathcal{O}_Y
\ar[rd]_{f^\sharp}  & &
f_{small}^{-1}\mathcal{O}_Y \ar[ll]^t \ar[ld]^{f^\sharp} \\
& \mathcal{O}_X
}
$$
soit commutatif. Supposons que $V \to Y$ soit \'etale, avec $V$ affine.
Écrivons $V = \Spec(B)$. Choisissons des générateurs $b_j \in B$, $j \in J$,
de $B$ comme $\mathbf{Z}$-algèbre. Posons
$T = \Spec(\mathbf{Z}[\{x_j\}_{j \in J}])$.
Dans la suite, nous utiliserons le fait que
$\Mor_{\Sch}(U, T) = \prod_{j \in J} \Gamma(U, \mathcal{O}_U)$
pour tout schéma $U$, sans le rappeler davantage.
L'homomorphisme surjectif d'anneaux $\mathbf{Z}[x_j] \to B$, $x_j \mapsto b_j$,
correspond à une immersion fermée $V \to T$.
Nous obtenons un monomorphisme
$$
i : V \longrightarrow T_Y = T \times Y
$$
d'espaces algébriques sur $Y$. En termes de faisceaux sur $Y_\etale$,
le morphisme $i$ induit une injection de faisceaux
$h_i : h_V \to \prod_{j \in J} \mathcal{O}_Y$.
Le changement de base $i' : X \times_Y V \to T_X$ de $i$ à $X$
est lui aussi un monomorphisme
(voir le
lemme \ref{spaces-lemma-base-change-representable-transformations-property} d'Espaces).
Ainsi, $i' : X \times_Y V \to T_X$ est un monomorphisme, ce qui
signifie à son tour que
$h_{i'} : h_{X \times_Y V} \to \prod_{j \in J} \mathcal{O}_X$
est une injection de faisceaux.
Par l'identification $f_{small}^{-1}h_V = h_{X \times_Y V}$ du
lemme \ref{lemma-stalk-pullback},
l'application $h_{i'}$ est égale à
$$
\xymatrix{
f_{small}^{-1}h_V \ar[r]^-{f^{-1}h_i} &
\prod_{j \in J} f_{small}^{-1}\mathcal{O}_Y
\ar[r]^{\prod f^\sharp} &
\prod_{j \in J} \mathcal{O}_X
}
$$
(vérification omise). Cela signifie que l'application
$t : f_{small}^{-1}h_V \to f_{small}^{-1}h_V$
s'insère dans le diagramme commutatif
$$
\xymatrix{
f_{small}^{-1}h_V \ar[r]^-{f^{-1}h_i} \ar[d]^t &
\prod_{j \in J} f_{small}^{-1}\mathcal{O}_Y
\ar[r]^-{\prod f^\sharp} \ar[d]^{\prod t} &
\prod_{j \in J} \mathcal{O}_X \ar[d]^{\text{id}}\\
f_{small}^{-1}h_V \ar[r]^-{f^{-1}h_i} &
\prod_{j \in J} f_{small}^{-1}\mathcal{O}_Y
\ar[r]^-{\prod f^\sharp} &
\prod_{j \in J} \mathcal{O}_X
}
$$
Le carré de droite est commutatif par l'hypothèse faite ci-dessus sur $t$.

Comme la composée des flèches horizontales est injective
d'après la discussion précédente, nous concluons que la flèche verticale de gauche
est elle aussi l'identité. Tout faisceau d'ensembles sur
$Y_\etale$ admet une surjection provenant d'une somme disjointe (énorme) de faisceaux
de la forme $h_V$, avec $V$ affine (combiner le
lemme \ref{lemma-alternative}
et le
lemme \ref{sites-lemma-sheaf-coequalizer-representable} de Sites).
Nous concluons ainsi que $t : f_{small}^{-1} \to f_{small}^{-1}$
est la transformation identique, comme souhaité.
\end{proof}

\begin{lemma}
\label{lemma-faithful}
Soit $S$ un schéma.
Soient $X$, $Y$ des espaces algébriques sur $S$.
Deux morphismes quelconques $a, b : X \to Y$ d'espaces algébriques sur $S$
pour lesquels il existe un $2$-isomorphisme
$(a_{small}, a^\sharp) \cong (b_{small}, b^\sharp)$
dans la $2$-catégorie des topos annelés sont égaux.
\end{lemma}

\begin{proof}
Soit $t : a_{small}^{-1} \to b_{small}^{-1}$ le $2$-isomorphisme.
Nous pouvons, de manière équivalente, considérer $t$ comme une transformation
$t : a_{spaces, \etale}^{-1} \to b_{spaces, \etale}^{-1}$
puisqu'il n'y a aucune différence entre les faisceaux sur $X_\etale$
et ceux sur $X_{spaces, \etale}$.
Choisissons un diagramme commutatif
$$
\xymatrix{
U \ar[d]_p \ar[r]_\alpha  & V \ar[d]^q \\
X \ar[r]^a & Y
}
$$
où $U$ et $V$ sont des schémas, et où $p$ et $q$ sont \'etales surjectifs.
Considérons le diagramme
$$
\xymatrix{
h_U \ar[r]_-\alpha \ar@{=}[d] & a_{spaces, \etale}^{-1}h_V \ar[d]^t \\
h_U \ar@{..>}[r] & b_{spaces, \etale}^{-1}h_V
}
$$
Puisque le faisceau $b_{spaces, \etale}^{-1}h_V$ est isomorphe à
$h_{V \times_{Y, b} X}$, nous voyons que la flèche pointillée provient d'un
morphisme de schémas
$\beta : U \to V$ qui s'insère dans un diagramme commutatif
$$
\xymatrix{
U \ar[d]_p \ar[r]_\beta  & V \ar[d]^q \\
X \ar[r]^b & Y
}
$$
Nous affirmons qu'il existe une suite de $2$-isomorphismes
\begin{align*}
(\alpha_{small}, \alpha^\sharp)
& \cong
(\alpha_{spaces, \etale}, \alpha^\sharp) \\
& \cong
(a_{spaces, \etale, c}, a_c^\sharp) \\
& \cong
(b_{spaces, \etale, d}, b_d^\sharp) \\
& \cong
(\beta_{spaces, \etale}, \beta^\sharp) \\
& \cong
(\beta_{small}, \beta^\sharp)
\end{align*}
Le premier et le dernier $2$-isomorphisme proviennent des identifications
entre les faisceaux sur $U_{spaces, \etale}$ et ceux sur
$U_\etale$, et de même pour $V$. Les deuxième et quatrième
$2$-isomorphismes sont ceux du
lemme \ref{lemma-relocalize-morphism},
où $c : U \to X \times_{a, Y} V$ est induit par $\alpha$ et
$d : U \to X \times_{b, Y} V$ est induit par $\beta$.
Le $2$-isomorphisme médian provient de la transformation $t$.
Plus précisément, le foncteur $a_{spaces, \etale, c}^{-1}$ correspond
au foncteur
$$
(\mathcal{H} \to h_V) \longmapsto
(a_{spaces, \etale}^{-1}\mathcal{H}
\times_{a_{spaces, \etale}^{-1}h_V, \alpha}
h_U \to
h_U)
$$
et de même pour $b_{spaces, \etale, d}^{-1}$ ; voir le
lemme \ref{sites-lemma-relocalize-morphism} de Sites.
On utilise ici l'identification des faisceaux sur $Y_{spaces, \etale}/V$
avec les flèches $(\mathcal{H} \to h_V)$ dans $\Sh(Y_{spaces, \etale})$,
et de même pour $U/X$ ; voir le
lemme \ref{sites-lemma-essential-image-j-shriek} de Sites.
Par cette identification, le faisceau structural $\mathcal{O}_V$ correspond à la
paire $(\mathcal{O}_Y \times h_V \to h_V)$, et de même
pour $\mathcal{O}_U$ ; voir le
Modules sur les sites,
lemme \ref{sites-modules-lemma-localize-compare}.
Puisque $t$ échange $\alpha$ et $\beta$, nous voyons que $t$ induit un isomorphisme
$$
t :
a_{spaces, \etale}^{-1}\mathcal{H}
\times_{a_{spaces, \etale}^{-1}h_V, \alpha}
h_U
\longrightarrow
b_{spaces, \etale}^{-1}\mathcal{H}
\times_{b_{spaces, \etale}^{-1}h_V, \beta}
h_U
$$
au-dessus de $h_U$, fonctoriellement en $(\mathcal{H} \to h_V)$. De plus, $t$ est compatible
avec $a_c^\sharp$ et $b_d^\sharp$, puisque $t$ est
compatible avec $a^\sharp$ et $b^\sharp$ d'après notre description
ci-dessus des faisceaux structuraux $\mathcal{O}_U$ et $\mathcal{O}_V$.
Par conséquent, les morphismes de topos annelés
$(\alpha_{small}, \alpha^\sharp)$ et $(\beta_{small}, \beta^\sharp)$
sont $2$-isomorphes. D'après
Cohomologie \'etale, le lemme \ref{etale-cohomology-lemma-faithful},
nous concluons que $\alpha = \beta$ ! Puisque $p : U \to X$ est une surjection
de faisceaux, il s'ensuit que $a = b$.
\end{proof}

\noindent
Voici le résultat principal de cette section.

\begin{theorem}
\label{theorem-fully-faithful}
Soient $X$, $Y$ des espaces algébriques sur $\Spec(\mathbf{Z})$.
Soit
$$
(g, g^\sharp) :
(\Sh(X_\etale), \mathcal{O}_X)
\longrightarrow
(\Sh(Y_\etale), \mathcal{O}_Y)
$$
un morphisme de topos localement annelés. Alors il existe un
unique morphisme d'espaces algébriques $f : X \to Y$ tel que
$(g, g^\sharp)$ soit isomorphe à $(f_{small}, f^\sharp)$.
Autrement dit, la construction
$$
\textit{Espaces}/\Spec(\mathbf{Z})
\longrightarrow \textit{Topos localement annelés},
\quad
X \longrightarrow (X_\etale, \mathcal{O}_X)
$$
est pleinement fidèle (les morphismes du membre de droite étant pris à $2$-isomorphisme près).
\end{theorem}

\begin{proof}
L'unicité a été démontrée au
lemme \ref{lemma-faithful}.
Il suffit donc de démontrer l'existence.
Dans cette preuve, nous utiliserons librement les identifications de
l'équation (\ref{equation-localize-at-scheme-ringed})
ainsi que le résultat du
lemme \ref{lemma-relocalize-morphism-at-schemes}.

\medskip\noindent
Soit $U \in \Ob(X_\etale)$, soit
$V \in \Ob(Y_\etale)$,
et soit $s \in g^{-1}h_V(U)$ une section. Nous pouvons considérer
$s$ comme une application de faisceaux $s : h_U \to g^{-1}h_V$. D'après le
Modules sur les sites,
lemme \ref{sites-modules-lemma-relocalize-morphism-ringed-topoi},
nous obtenons un diagramme commutatif de morphismes de topos annelés
$$
\xymatrix{
(\Sh(X_\etale/U), \mathcal{O}_U)
\ar[rr]_-{(j, j^\sharp)} \ar[d]_{(g_s, g_s^\sharp)} & &
(\Sh(X_\etale), \mathcal{O}_X) \ar[d]^{(g, g^\sharp)} \\
(\Sh(V_\etale), \mathcal{O}_V) \ar[rr] & &
(\Sh(Y_\etale), \mathcal{O}_Y).
}
$$
D'après
Cohomologie \'etale, le théorème \ref{etale-cohomology-theorem-fully-faithful},
nous obtenons un unique morphisme de schémas $f_s : U \to V$ tel que
$(g_s, g_s^\sharp)$ soit $2$-isomorphe à $(f_{s, small}, f_s^\sharp)$.
La construction $(U, V, s) \leadsto f_s$ que nous venons de décrire satisfait
la propriété de fonctorialité suivante. Supposons donnés des morphismes
$a : U' \to U$ dans $X_\etale$ et $b : V' \to V$ dans $Y_\etale$,
ainsi qu'une application $s' : h_{U'} \to g^{-1}h_{V'}$ telle que le diagramme
$$
\xymatrix{
h_{U'} \ar[d]_a \ar[r]_{s'} & g^{-1}h_{V'} \ar[d]^{g^{-1}b} \\
h_U \ar[r]^s & g^{-1}h_V
}
$$
soit commutatif. Alors le diagramme
$$
\xymatrix{
U' \ar[r]_-{f_{s'}} \ar[d]_a & u(V') \ar[d]^{u(b)} \\
U \ar[r]^-{f_s} & u(V)
}
$$
de schémas est commutatif. En effet, la même condition
vaut pour les morphismes $(g_s, g_s^\sharp)$ construits dans
Modules sur les sites,
au lemme \ref{sites-modules-lemma-relocalize-morphism-ringed-topoi},
et l'on applique l'unicité de
Cohomologie \'etale, théorème \ref{etale-cohomology-theorem-fully-faithful}.

\medskip\noindent
Il s'agit de recoller les morphismes $f_s$ en un morphisme d'espaces
algébriques. Pour ce faire, choisissons d'abord un schéma $V$ et un morphisme \'etale
surjectif $V \to Y$. Ainsi, $h_V \to *$ est surjectif et, par conséquent,
$g^{-1}h_V \to *$ est lui aussi surjectif. Il existe donc un schéma $U$,
un morphisme \'etale surjectif $U \to X$ et un morphisme
$s : h_U \to g^{-1}h_V$. Posons ensuite $R = V \times_Y V$ et
$R' = U \times_X U$. Nous obtenons alors
$g^{-1}h_R = g^{-1}h_V \times g^{-1}h_V$, car $g^{-1}$ est exact.
Ainsi, $s$ induit un morphisme $s \times s : h_{R'} \to g^{-1}h_R$.
En appliquant les constructions précédentes, nous obtenons un
diagramme commutatif de morphismes de schémas
$$
\xymatrix{
R' \ar@<1ex>[d] \ar@<-1ex>[d] \ar[rr]_{f_{s \times s}} & &
R \ar@<1ex>[d] \ar@<-1ex>[d] \\
U \ar[rr]^{f_s} & &
V
}
$$
Comme $X = U/R'$ et $Y = V/R$ (voir le
lemme \ref{spaces-lemma-space-presentation} d'Espaces),
nous concluons que ce diagramme
définit un morphisme d'espaces algébriques $f : X \to Y$ qui s'insère
dans le diagramme commutatif évident.
Il reste à montrer que $(f_{small}, f^\sharp)$ est
$2$-isomorphe à $(g, g^\sharp)$.
Soient $t_V : f_{s, small}^{-1} \to g_s^{-1}$ et
$t_R : f_{s \times s, small}^{-1} \to g_{s \times s}^{-1}$ les
$2$-isomorphismes fournis par la construction précédente.
Soit $\mathcal{G}$ un faisceau sur $Y_\etale$. Nous voyons alors que
$t_V$ définit un isomorphisme
$$
f_{small}^{-1}\mathcal{G}|_{U_\etale}
=
f_{s, small}^{-1}\mathcal{G}|_{V_\etale}
\xrightarrow{t_V}
g_s^{-1}\mathcal{G}|_{V_\etale}
=
g^{-1}\mathcal{G}|_{U_\etale}.
$$
De plus, l'image inverse de cet isomorphisme sur $R'$ par l'une ou l'autre projection
$R' \to U$ est l'isomorphisme
$$
f_{small}^{-1}\mathcal{G}|_{R'_\etale}
=
f_{s \times s, small}^{-1}\mathcal{G}|_{R_\etale}
\xrightarrow{t_R}
g_{s \times s}^{-1}\mathcal{G}|_{R_\etale}
=
g^{-1}\mathcal{G}|_{R'_\etale}.
$$
Comme $\{U \to X\}$ est un recouvrement du site $X_{spaces, \etale}$,
cela signifie que le premier isomorphisme affiché descend en un isomorphisme
$t : f_{small}^{-1}\mathcal{G} \to g^{-1}\mathcal{G}$
de faisceaux (un petit détail est omis). L'isomorphisme est fonctoriel
en $\mathcal{G}$, puisque $t_V$ et $t_R$ sont des transformations de foncteurs.
Enfin, $t$ est compatible avec $f^\sharp$ et $g^\sharp$, puisque
$t_V$ et $t_R$ le sont (certains détails sont omis).
Cela achève la preuve du théorème.
\end{proof}

\begin{lemma}
\label{lemma-isomorphism-ringed-topoi}
Soient $X$, $Y$ des espaces algébriques sur $\mathbf{Z}$. Si
$$
(g, g^\sharp) :
(\Sh(X_\etale), \mathcal{O}_X)
\longrightarrow
(\Sh(Y_\etale), \mathcal{O}_Y)
$$
est un isomorphisme de topos annelés, alors il existe un unique
morphisme d'espaces algébriques $f : X \to Y$ tel que
$(g, g^\sharp)$ soit isomorphe à $(f_{small}, f^\sharp)$ ;
de plus, $f$ est un isomorphisme d'espaces algébriques.
\end{lemma}

\begin{proof}
D'après le
théorème \ref{theorem-fully-faithful},
il suffit de montrer que $(g, g^\sharp)$ est un morphisme de
topos localement annelés. D'après le
lemme \ref{sites-modules-lemma-locally-ringed-morphism} de Modules sur les sites
(et puisque le site $X_\etale$ a assez de points),
il suffit de vérifier que l'homomorphisme
$\mathcal{O}_{Y, q} \to \mathcal{O}_{X, p}$ induit par $g^\sharp$
est un homomorphisme local d'anneaux locaux, où $q = g \circ p$ et $p$ est un point quelconque de
$X_\etale$. Comme c'est un isomorphisme, cela est clair.
\end{proof}














\section{Faisceaux quasi-cohérents sur les espaces algébriques}
\label{section-quasi-coherent}

\noindent
Dans les
sections \ref{descent-section-quasi-coherent-sheaves},
\ref{descent-section-quasi-coherent-cohomology} et
\ref{descent-section-quasi-coherent-sheaves-bis}
de Descente, nous avons vu que, pour un schéma $U$, il n'y a aucune différence entre un
$\mathcal{O}_U$-module quasi-cohérent sur $U$ et un
$\mathcal{O}$-module quasi-cohérent sur le petit site \'etale de $U$. La définition
suivante est donc compatible avec notre notion initiale de faisceau quasi-cohérent
sur un schéma
(Schémas, section \ref{schemes-section-quasi-coherent})
lorsqu'on l'applique à un espace algébrique représentable.

\begin{definition}
\label{definition-quasi-coherent}
Soit $S$ un schéma. Soit $X$ un espace algébrique sur $S$.
Un {\it quasi-cohérent} $\mathcal{O}_X$-module
est un module quasi-cohérent sur le site annelé
$(X_\etale, \mathcal{O}_X)$ au sens de
Modules sur les sites,
la définition \ref{sites-modules-definition-site-local}.
La catégorie des faisceaux quasi-cohérents sur $X$ se note
$\QCoh(\mathcal{O}_X)$.
\end{definition}

\noindent
Remarquons que, la quasi-cohérence étant une notion intrinsèque (voir le
lemme \ref{sites-modules-lemma-special-locally-free} de Modules sur les sites),
cela revient à dire que le $\mathcal{O}_X$-module correspondant
sur $X_{spaces, \etale}$ est quasi-cohérent.

\medskip\noindent
Comme d'habitude, les faisceaux quasi-cohérents se comportent bien par image inverse.

\begin{lemma}
\label{lemma-pullback-quasi-coherent}
Soit $S$ un schéma.
Soit $f : X \to Y$ un morphisme d'espaces algébriques sur $S$.
Le foncteur image inverse
$f^* : \textit{Mod}(\mathcal{O}_Y) \to \textit{Mod}(\mathcal{O}_X)$
préserve les faisceaux quasi-cohérents.
\end{lemma}

\begin{proof}
C'est un fait général ; voir le
lemme \ref{sites-modules-lemma-local-pullback} de Modules sur les sites.
\end{proof}

\noindent
Remarquons que ce foncteur image inverse coïncide avec le foncteur image inverse usuel
entre faisceaux de modules quasi-cohérents lorsque $X$ et $Y$ sont des
schémas ; voir la
proposition de Descente
\ref{descent-proposition-equivalence-quasi-coherent-functorial}.
Voici le lemme incontournable qui compare cette notion aux faisceaux quasi-cohérents
sur les objets du petit site \'etale de $X$.

\begin{lemma}
\label{lemma-characterize-quasi-coherent-small-etale}
Soit $S$ un schéma. Soit $X$ un espace algébrique sur $S$.
Un $\mathcal{O}_X$-module quasi-cohérent $\mathcal{F}$
est donné par les données suivantes :
\begin{enumerate}
\item pour tout $U \in \Ob(X_\etale)$, un
$\mathcal{O}_U$-module quasi-cohérent $\mathcal{F}_U$ sur $U_\etale$ ;
\item pour tout $f : U' \to U$ dans $X_\etale$, un isomorphisme
$c_f : f_{small}^*\mathcal{F}_U \to \mathcal{F}_{U'}$.
\end{enumerate}
Ces données sont soumises à la condition que, pour tous $f : U' \to U$
et $g : U'' \to U'$ dans $X_\etale$, la composée
$c_g \circ g_{small}^*c_f$ soit égale à $c_{f \circ g}$.
\end{lemma}

\begin{proof}
Combiner les lemmes \ref{lemma-pullback-quasi-coherent} et
\ref{lemma-characterize-module-small-etale}.
\end{proof}

\begin{lemma}
\label{lemma-stalk-quasi-coherent}
Soit $S$ un schéma.
Soit $X$ un espace algébrique sur $S$.
Soit $\mathcal{F}$ un $\mathcal{O}_X$-module quasi-cohérent.
Soit $x \in |X|$ un point et soit $\overline{x}$ un point
géométrique situé au-dessus de $x$. Enfin, soit
$\varphi : (U, \overline{u}) \to (X, \overline{x})$
un voisinage \'etale, où $U$ est un schéma.
Alors
$$
(\varphi^*\mathcal{F})_u \otimes_{\mathcal{O}_{U, u}}
\mathcal{O}_{X, \overline{x}} =
\mathcal{F}_{\overline{x}}
$$
où $u \in U$ est l'image de $\overline{u}$.
\end{lemma}

\begin{proof}
Remarquons que $\mathcal{O}_{X, \overline{x}} = \mathcal{O}_{U, u}^{sh}$ d'après le
lemme \ref{lemma-describe-etale-local-ring} ;
le produit tensoriel a donc un sens. De plus, il ressort de la
définition \ref{definition-stalk}
que
$$
\mathcal{F}_{\overline{x}} = \colim (\varphi^*\mathcal{F})_u
$$
où la colimite porte sur les $\varphi : (U, \overline{u}) \to (X, \overline{x})$
comme dans le lemme. Il existe donc un homomorphisme canonique du membre de gauche vers celui de droite dans
l'énoncé du lemme. Nous disposons d'une description colimite analogue pour
$\mathcal{O}_{X, \overline{x}}$
et, d'après le
lemme \ref{lemma-characterize-quasi-coherent-small-etale},
nous avons
$$
((\varphi')^*\mathcal{F})_{u'} =
(\varphi^*\mathcal{F})_u \otimes_{\mathcal{O}_{U, u}} \mathcal{O}_{U', u'}
$$
dès que $(U', \overline{u}') \to (U, \overline{u})$ est un morphisme de
voisinages \'etales. Pour achever la preuve, nous utilisons le fait que
$\otimes$ commute aux colimites.
\end{proof}

\begin{lemma}
\label{lemma-stalk-pullback-quasi-coherent}
Soit $S$ un schéma. Soit $f : X \to Y$ un morphisme d'espaces algébriques
sur $S$. Soit $\mathcal{G}$ un $\mathcal{O}_Y$-module quasi-cohérent.
Soit $\overline{x}$ un point géométrique de $X$ et soit
$\overline{y} = f \circ \overline{x}$ son image dans $Y$.
Alors il existe un isomorphisme canonique
$$
(f^*\mathcal{G})_{\overline{x}} =
\mathcal{G}_{\overline{y}} \otimes_{\mathcal{O}_{Y, \overline{y}}}
\mathcal{O}_{X, \overline{x}}
$$
de la fibre de l'image inverse avec le produit tensoriel de la fibre
par l'anneau local de $X$ en $\overline{x}$.
\end{lemma}

\begin{proof}
Comme $f^*\mathcal{G} =
f_{small}^{-1}\mathcal{G} \otimes_{f_{small}^{-1}\mathcal{O}_Y} \mathcal{O}_X$
cela résulte de la description des fibres des images inverses donnée au
lemme \ref{lemma-stalk-pullback}
et du fait que la prise des fibres commute aux produits tensoriels.
Voici une démonstration plus directe.
Choisissons un diagramme commutatif
$$
\xymatrix{
U \ar[d]_p \ar[r]_\alpha  & V \ar[d]^q \\
X \ar[r]^f & Y
}
$$
où $U$ et $V$ sont des schémas, et où $p$ et $q$ sont \'etales surjectifs.
D'après le
lemme \ref{lemma-geometric-lift-to-usual},
nous pouvons choisir un point géométrique $\overline{u}$ de $U$ tel que
$\overline{x} = p \circ \overline{u}$. Posons
$\overline{v} = \alpha \circ \overline{u}$.
Nous voyons alors que
\begin{align*}
(f^*\mathcal{G})_{\overline{x}} & =
(p^*f^*\mathcal{G})_u \otimes_{\mathcal{O}_{U, u}}
\mathcal{O}_{X, \overline{x}} \\
& = (\alpha^*q^*\mathcal{G})_u \otimes_{\mathcal{O}_{U, u}}
\mathcal{O}_{X, \overline{x}} \\
& = (q^*\mathcal{G})_v \otimes_{\mathcal{O}_{V, v}}
\mathcal{O}_{U, u} \otimes_{\mathcal{O}_{U, u}}
\mathcal{O}_{X, \overline{x}} \\
& = (q^*\mathcal{G})_v \otimes_{\mathcal{O}_{V, v}}
\mathcal{O}_{X, \overline{x}} \\
& = (q^*\mathcal{G})_v \otimes_{\mathcal{O}_{V, v}}
\mathcal{O}_{Y, \overline{y}} \otimes_{\mathcal{O}_{Y, \overline{y}}}
\mathcal{O}_{X, \overline{x}} \\
& = \mathcal{G}_{\overline{y}} \otimes_{\mathcal{O}_{Y, \overline{y}}}
\mathcal{O}_{X, \overline{x}}
\end{align*}
Nous avons utilisé ici
le lemme \ref{lemma-stalk-quasi-coherent} (deux fois)
et le résultat correspondant pour les images inverses de faisceaux quasi-cohérents
sur les schémas ; voir le
lemme \ref{sheaves-lemma-stalk-pullback-modules} de Faisceaux.
\end{proof}

\begin{lemma}
\label{lemma-characterize-quasi-coherent}
Soit $S$ un schéma. Soit $X$ un espace algébrique sur $S$.
Soit $\mathcal{F}$ un faisceau de $\mathcal{O}_X$-modules.
Les conditions suivantes sont équivalentes :
\begin{enumerate}
\item $\mathcal{F}$ est un $\mathcal{O}_X$-module quasi-cohérent ;
\item il existe un morphisme \'etale $f : Y \to X$
d'espaces algébriques sur $S$, avec $|f| : |Y| \to |X|$ surjectif,
tel que $f^*\mathcal{F}$ soit quasi-cohérent sur $Y$ ;
\item il existe un schéma $U$ et un morphisme \'etale surjectif
$\varphi : U \to X$ tel que $\varphi^*\mathcal{F}$ soit un
$\mathcal{O}_U$-module quasi-cohérent ;
\item pour tout schéma affine $U$ et tout morphisme \'etale $\varphi : U \to X$, la
restriction $\varphi^*\mathcal{F}$ est un $\mathcal{O}_U$-module quasi-cohérent.
\end{enumerate}
\end{lemma}

\begin{proof}
Il est clair que (1) implique (2) en considérant $\text{id}_X$.
Supposons que $f : Y \to X$ soit comme dans (2), et soit $V \to Y$ un morphisme
\'etale surjectif d'un schéma vers $Y$. Alors la composée $V \to X$ est elle aussi
\'etale surjective
et, d'après le lemme \ref{lemma-pullback-quasi-coherent}, l'image inverse de $\mathcal{F}$
sur $V$ est elle aussi quasi-cohérente. Ainsi, (2) implique (3).

\medskip\noindent
Soit $U \to X$ comme dans (3). Employons l'abus de notation introduit
dans l'équation (\ref{equation-restrict-modules}).
Comme $\mathcal{F}|_{U_\etale}$ est quasi-cohérent, il existe un
recouvrement \'etale $\{U_i \to U\}$ tel que
$\mathcal{F}|_{U_{i, \etale}}$ admette une présentation globale ; voir la
définition \ref{sites-modules-definition-global} de Modules sur les sites et le
lemme \ref{sites-modules-lemma-local-final-object}.
Soit $V \to X$ un objet de $X_\etale$. Puisque $U \to X$ est
surjectif et \'etale, la famille de morphismes $\{U_i \times_X V \to V\}$ est un
recouvrement \'etale
de $V$. Par les morphismes $U_i \times_X V \to U_i$, nous pouvons restreindre les
présentations globales de $\mathcal{F}|_{U_{i, \etale}}$ pour obtenir une présentation
globale de $\mathcal{F}|_{(U_i \times_X V)_\etale}$.
Le faisceau $\mathcal{F}$ sur $X_\etale$ satisfait donc la condition de la
définition \ref{sites-modules-definition-site-local} de Modules sur les sites
et est par conséquent quasi-cohérent.

\medskip\noindent
L'équivalence de (3) et (4) résulte du fait que tout schéma possède
un recouvrement par des ouverts affines.
\end{proof}

\begin{lemma}
\label{lemma-properties-quasi-coherent}
Soit $S$ un schéma. Soit $X$ un espace algébrique sur $S$.
La catégorie $\QCoh(\mathcal{O}_X)$ des faisceaux quasi-cohérents sur $X$
possède les propriétés suivantes :
\begin{enumerate}
\item Toute somme directe de faisceaux quasi-cohérents est quasi-cohérente.
\item Toute colimite de faisceaux quasi-cohérents est quasi-cohérente.
\item Le noyau et le conoyau d'un morphisme de faisceaux quasi-cohérents
sont quasi-cohérents.
\item Dans une suite exacte courte de $\mathcal{O}_X$-modules
$0 \to \mathcal{F}_1 \to \mathcal{F}_2 \to \mathcal{F}_3 \to 0$
si deux termes sur trois sont quasi-cohérents, le troisième l'est aussi.
\item Le produit tensoriel de deux $\mathcal{O}_X$-modules quasi-cohérents
est quasi-cohérent.
\item Étant donnés deux $\mathcal{O}_X$-modules quasi-cohérents
$\mathcal{F}$, $\mathcal{G}$, dont $\mathcal{F}$
est de présentation finie (voir la
section \ref{section-properties-modules}),
le Hom interne
$\SheafHom_{\mathcal{O}_X}(\mathcal{F}, \mathcal{G})$
est quasi-cohérent.
\end{enumerate}
\end{lemma}

\begin{proof}
Si $X$ est un schéma, c'est le
lemme \ref{descent-lemma-properties-quasi-coherent} de Descente.
Nous allons réduire le lemme à ce cas par localisation \'etale.

\medskip\noindent
Choisissons un schéma $U$ et un morphisme \'etale surjectif $\varphi : U \to X$.
Nous adopterons les notations $\textit{Mod}(\mathcal{O}_U) =
\textit{Mod}(U_\etale, \mathcal{O}_U)$ et
$\QCoh(\mathcal{O}_U) = \QCoh(U_\etale, \mathcal{O}_U)$ ; autrement
dit, bien que $U$ soit un schéma, nous considérons les modules quasi-cohérents
sur $U$ comme des modules sur le petit site \'etale de $U$.
D'après le lemme \ref{lemma-pullback-quasi-coherent}, nous avons un diagramme commutatif
$$
\xymatrix{
\QCoh(\mathcal{O}_X) \ar[r]_{\varphi^*} \ar[d] &
\QCoh(\mathcal{O}_U) \ar[d] \\
\textit{Mod}(\mathcal{O}_X) \ar[r]^{\varphi^*} &
\textit{Mod}(\mathcal{O}_U)
}
$$
La flèche horizontale inférieure est le foncteur de restriction
(\ref{equation-restrict-modules})
$\mathcal{G} \mapsto \mathcal{G}|_{U_\etale}$.
Ce foncteur possède à la fois un adjoint à gauche et un adjoint à droite ; voir la
section \ref{sites-modules-section-localize} de Modules sur les sites ;
il commute donc à toutes les limites et colimites.
De plus, nous savons qu'un objet de $\textit{Mod}(\mathcal{O}_X)$ appartient à
$\QCoh(\mathcal{O}_X)$ si et seulement si sa restriction à $U$ appartient à
$\QCoh(\mathcal{O}_U)$ ; voir le
lemme \ref{lemma-characterize-quasi-coherent}.
Ces préliminaires étant établis, nous pouvons commencer la preuve.

\medskip\noindent
Preuve de (1). Soit $\mathcal{F}_i$, $i \in I$, une famille de
$\mathcal{O}_X$-modules quasi-cohérents. D'après la discussion précédente, nous avons
$$
\Big(\bigoplus \mathcal{F}_i\Big)|_{U_\etale} =
\bigoplus \mathcal{F}_i|_{U_\etale}
$$
Chacun des modules $\mathcal{F}_i|_{U_\etale}$ est quasi-cohérent.
La somme directe est donc quasi-cohérente d'après le cas des schémas.
Par conséquent, $\bigoplus \mathcal{F}_i$ est quasi-cohérent, puisque sa
restriction à $U$ est un module quasi-cohérent.

\medskip\noindent
Preuve de (2). Soit $\mathcal{I} \to \QCoh(\mathcal{O}_X)$,
$i \mapsto \mathcal{F}_i$, un diagramme. Alors
$$
(\colim \mathcal{F}_i)|_{U_\etale} = \colim \mathcal{F}_i|_{U_\etale}
$$
d'après la discussion précédente, et nous concluons de la même manière.

\medskip\noindent
Preuve de (3). Soit $a : \mathcal{F} \to \mathcal{F}'$
une flèche de $\QCoh(\mathcal{O}_X)$. Alors
nous avons $\Ker(a)|_{U_\etale} = \Ker(a|_{U_\etale})$
et $\Coker(a)|_{U_\etale} = \Coker(a|_{U_\etale})$,
et nous concluons de la même manière.

\medskip\noindent
Preuve de (4). La restriction
$0 \to \mathcal{F}_1|_{U_\etale} \to \mathcal{F}_2|_{U_\etale}
\to \mathcal{F}_3|_{U_\etale} \to 0$
est une suite exacte courte. Elle possède donc la propriété des deux sur trois,
et nous concluons comme précédemment.

\medskip\noindent
Preuve de (5). Soient $\mathcal{F}$ et $\mathcal{G}$ dans $\QCoh(\mathcal{O}_X)$.
Alors nous avons
$$
(\mathcal{F} \otimes_{\mathcal{O}_X} \mathcal{G})_{U_\etale} =
\mathcal{F}|_{U_\etale} \otimes_{\mathcal{O}_U} \mathcal{G}|_{U_\etale}
$$
et nous concluons comme précédemment.

\medskip\noindent
Preuve de (6). Soient $\mathcal{F}$ et $\mathcal{G}$
dans $\QCoh(\mathcal{O}_X)$, avec $\mathcal{F}$
de présentation finie. Nous avons
$$
\SheafHom_{\mathcal{O}_X}(\mathcal{F}, \mathcal{G})|_{U_\etale} =
\SheafHom_{\mathcal{O}_U}(\mathcal{F}|_{U_\etale}, \mathcal{G}|_{U_\etale})
$$
En effet, la restriction est une localisation ; voir la
section \ref{section-localize}, en particulier la formule
(\ref{equation-localize-at-scheme-ringed}), et la formation du Hom
interne commute à la localisation ; voir le
lemme \ref{sites-modules-lemma-internal-hom-restriction} de Modules sur les sites.
Nous concluons donc comme précédemment.
\end{proof}

\noindent
En général, l'image directe d'un faisceau quasi-cohérent
par un morphisme d'espaces algébriques n'est pas quasi-cohérente. Nous reviendrons sur cette
question dans
Morphismes d'espaces, section \ref{spaces-morphisms-section-pushforward}.




\section{Propriétés des modules}
\label{section-properties-modules}

\noindent
Dans
Modules sur les sites, les sections
\ref{sites-modules-section-global},
\ref{sites-modules-section-local} et la
définition \ref{sites-modules-definition-flat},
nous avons défini plusieurs propriétés intrinsèques des
$\mathcal{O}$-modules sur tout topos annelé. Si $X$ est un espace
algébrique, nous appliquerons librement ces notions aux modules sur le site
annelé $(X_\etale, \mathcal{O}_X)$, ou, de manière équivalente, sur le site annelé
$(X_{spaces, \etale}, \mathcal{O}_X)$.

\medskip\noindent
Propriétés globales $\mathcal{P}$ :
\begin{enumerate}
\item[(a)] {\it libre} ;
\item[(b)] {\it libre de rang fini} ;
\item[(c)] {\it engendré par ses sections globales} ;
\item[(d)] {\it engendré par un nombre fini de sections globales} ;
\item[(e)] admettre une {\it présentation globale} ;
\item[(f)] admettre une {\it présentation globale finie}.
\end{enumerate}
Propriétés locales $\mathcal{P}$ :
\begin{enumerate}
\item[(g)] {\it localement libre} ;
\item[(h)] {\it localement libre de rang fini} ;
\item[(i)] {\it localement engendré par des sections} ;
\item[(j)] {\it localement engendré par $r$ sections} ;
\item[(k)] {\it de type fini} ;
\item[(l)] {\it quasi-cohérent} (voir la section \ref{section-quasi-coherent}) ;
\item[(m)] {\it de présentation finie} ;
\item[(n)] {\it cohérent} ;
\item[(o)] {\it plat}.
\end{enumerate}
Voici quelques résultats qui découlent immédiatement des définitions :
\begin{enumerate}
\item Dans chaque cas, sauf pour $\mathcal{P}=$``cohérent'', la propriété
est préservée par image inverse ; voir les
lemmes \ref{sites-modules-lemma-global-pullback},
\ref{sites-modules-lemma-local-pullback} et
\ref{sites-modules-lemma-pullback-flat}.
\item Chacune des propriétés ci-dessus (y compris la cohérence) est préservée par
image inverse le long d'un morphisme \'etale d'espaces algébriques (car, dans ce
cas, l'image inverse est donnée par restriction ; voir le
lemme \ref{lemma-etale-morphism-topoi}).
\item Supposons que $f : Y \to X$ soit un morphisme \'etale surjectif d'espaces
algébriques. Pour chacune des propriétés locales (g) -- (o), si
$f^*\mathcal{F}$ possède $\mathcal{P}$, alors $\mathcal{F}$ possède
$\mathcal{P}$. Cela résulte du fait que $\{Y \to X\}$ est un recouvrement de
$X_{spaces, \etale}$ et du
lemme \ref{sites-modules-lemma-local-final-object} de Modules sur les sites.
\item Si $X$ est un schéma, si $\mathcal{F}$ est un module quasi-cohérent
sur $X_\etale$ et si $\mathcal{P}$ est une propriété autre que ``cohérent'' ou
``localement libre'', alors $\mathcal{P}$ pour $\mathcal{F}$ sur $X_\etale$
équivaut à la propriété correspondante pour
$\mathcal{F}|_{X_{Zar}}$ ; autrement dit, elle correspond à $\mathcal{P}$
pour $\mathcal{F}$ considéré comme faisceau quasi-cohérent
sur le schéma $X$. Voir le
lemme \ref{descent-lemma-equivalence-quasi-coherent-properties} de Descente.
\item Si $X$ est un schéma localement noethérien et si $\mathcal{F}$ est un
module quasi-cohérent sur $X_\etale$, alors $\mathcal{F}$
est cohérent sur $X_\etale$ si et seulement si
$\mathcal{F}|_{X_{Zar}}$ est cohérent ; autrement dit, cela équivaut à ce que
le faisceau quasi-cohérent usuel sur le schéma $X$ soit
cohérent. Voir le
lemme \ref{descent-lemma-equivalence-quasi-coherent-properties} de Descente.
\end{enumerate}



\section{Modules localement projectifs}
\label{section-locally-projective}

\noindent
Rappelons que, dans
Propriétés, section \ref{properties-section-locally-projective},
nous avons défini la notion de module quasi-cohérent
localement projectif.

\begin{lemma}
\label{lemma-locally-projective}
Soit $S$ un schéma. Soit $X$ un espace algébrique sur $S$.
Soit $\mathcal{F}$ un $\mathcal{O}_X$-module quasi-cohérent.
Les conditions suivantes sont équivalentes :
\begin{enumerate}
\item il existe un schéma $U$ et un morphisme \'etale surjectif
$U \to X$ tels que la restriction $\mathcal{F}|_U$ soit localement projective
sur $U$ ;
\item pour tout schéma $U$ et tout morphisme \'etale
$U \to X$, la restriction $\mathcal{F}|_U$ est localement projective
sur $U$.
\end{enumerate}
\end{lemma}

\begin{proof}
Soit $U \to X$ comme dans (1), et soit $V \to X$ \'etale, où
$V$ est un schéma. Alors $\{U \times_X V \to V\}$ est un recouvrement fppf
de schémas. Ainsi, si $\mathcal{F}|_U$ est localement projectif, alors
$\mathcal{F}|_{U \times_X V}$ est localement projectif (voir le
lemme \ref{properties-lemma-locally-projective-pullback} de Propriétés)
et $\mathcal{F}|_V$ est donc localement projectif ; voir le
lemme \ref{descent-lemma-locally-projective-descends} de Descente.
\end{proof}

\begin{definition}
\label{definition-locally-projective}
Soit $S$ un schéma. Soit $X$ un espace algébrique sur $S$.
Soit $\mathcal{F}$ un $\mathcal{O}_X$-module quasi-cohérent.
On dit que $\mathcal{F}$ est {\it localement projectif}
si les conditions équivalentes du
lemme \ref{lemma-locally-projective}
sont satisfaites.
\end{definition}

\begin{lemma}
\label{lemma-locally-projective-pullback}
Soit $S$ un schéma.
Soit $f : X \to Y$ un morphisme d'espaces algébriques sur $S$.
Soit $\mathcal{G}$ un $\mathcal{O}_Y$-module quasi-cohérent.
Si $\mathcal{G}$ est localement projectif sur $Y$, alors $f^*\mathcal{G}$
est localement projectif sur $X$.
\end{lemma}

\begin{proof}
Choisissons un morphisme \'etale surjectif $V \to Y$, où $V$ est un schéma.
Choisissons un morphisme \'etale surjectif $U \to V \times_Y X$, où
$U$ est un schéma. Notons $\psi : U \to V$ le morphisme induit.
Alors
$$
f^*\mathcal{G}|_U = \psi^*(\mathcal{G}|_V)
$$
Le lemme résulte donc de la définition et du résultat dans le
cas des schémas ; voir le
lemme \ref{properties-lemma-locally-projective-pullback} de Propriétés.
\end{proof}





\section{Faisceaux quasi-cohérents et présentations}
\label{section-quasi-coherent-presentation}

\noindent
Soit $S$ un schéma. Soit $X$ un espace algébrique sur $S$.
Soit $X = U/R$ une présentation de $X$ provenant d'un morphisme
\'etale surjectif quelconque $\varphi : U \to X$ ; voir la
définition \ref{spaces-definition-presentation} d'Espaces.
Nous obtenons en particulier un groupoïde $(U, R, s, t, c)$ tel que
$j = (t, s) : R \to U \times_S U$ ; voir le
lemme \ref{groupoids-lemma-equivalence-groupoid} de Groupoïdes.
Dans la
définition \ref{groupoids-definition-groupoid-module} de Groupoïdes,
nous avons défini la notion de faisceau quasi-cohérent
sur un groupoïde arbitraire. Ces notions étant posées, nous avons
l'observation suivante.

\begin{proposition}
\label{proposition-quasi-coherent}
Conservons $S$, $\varphi : U \to X$ et $(U, R, s, t, c)$ comme ci-dessus.
Pour tout $\mathcal{O}_X$-module quasi-cohérent $\mathcal{F}$, le
faisceau $\varphi^*\mathcal{F}$ est muni d'un
isomorphisme canonique
$$
\alpha : t^*\varphi^*\mathcal{F} \longrightarrow s^*\varphi^*\mathcal{F}
$$
qui satisfait aux conditions de la
définition \ref{groupoids-definition-groupoid-module} de Groupoïdes
et définit donc un faisceau quasi-cohérent sur $(U, R, s, t, c)$.
Le foncteur $\mathcal{F} \mapsto (\varphi^*\mathcal{F}, \alpha)$
définit une équivalence de catégories
$$
\begin{matrix}
\text{Catégorie des} \\
\mathcal{O}_X\text{-Modules quasi-cohérents}
\end{matrix}
\longleftrightarrow
\begin{matrix}
\text{Catégorie des Modules}\\
\text{quasi-cohérents sur }(U, R, s, t, c)
\end{matrix}
$$
\end{proposition}

\begin{proof}
Dans l'énoncé de la proposition et dans cette preuve, nous considérons un
faisceau quasi-cohérent sur un schéma comme un faisceau quasi-cohérent sur le petit
site \'etale de ce schéma. Cela est permis par les résultats de
Descente, sections \ref{descent-section-quasi-coherent-sheaves},
\ref{descent-section-quasi-coherent-cohomology} et
\ref{descent-section-quasi-coherent-sheaves-bis}.

\medskip\noindent
L'existence de $\alpha$ résulte du fait que
$\varphi \circ t = \varphi \circ s$ et de la fonctorialité de l'image inverse
par rapport au morphisme ; voir la discussion autour de
l'équation (\ref{equation-push-pull}). De la même manière, c'est-à-dire par
fonctorialité de l'image inverse, nous voyons que l'isomorphisme $\alpha$ satisfait
à la condition (1) de la
définition \ref{groupoids-definition-groupoid-module} de Groupoïdes.
Pour la condition (2) de la définition, il suffit de constater que $\alpha$
est un isomorphisme, ce qui est clair. La construction
$\mathcal{F} \mapsto (\varphi^*\mathcal{F}, \alpha)$
est manifestement fonctorielle par rapport au faisceau quasi-cohérent $\mathcal{F}$.
Nous obtenons donc le foncteur de gauche à droite dans la formule
affichée de la proposition.

\medskip\noindent
Réciproquement, supposons que $(\mathcal{F}, \alpha)$ soit un faisceau
quasi-cohérent sur $(U, R, s, t, c)$. Soit $V \to X$ un objet de $X_\etale$.
Alors le morphisme $V' = U \times_X V \to V$ est un morphisme \'etale
surjectif de schémas, et $\{V' \to V\}$ est donc un recouvrement \'etale
de $V$. En outre, l'image inverse du faisceau quasi-cohérent $\mathcal{F}$
est un faisceau quasi-cohérent $\mathcal{F}'$ sur $V'$.
Comme $R = U \times_X U$, avec $t = \text{pr}_0$ et $s = \text{pr}_1$,
nous avons $V' \times_V V' = R \times_X V$, les morphismes de projection
$V' \times_V V' \to V'$ étant obtenus par changement de base à partir de $t$ et $s$. Ainsi,
l'image inverse de $\alpha$ est un isomorphisme
$\alpha' : \text{pr}_0^*\mathcal{F}' \to \text{pr}_1^*\mathcal{F}'$, et
le couple $(\mathcal{F}', \alpha')$ est une donnée de descente pour les faisceaux
quasi-cohérents relativement à $\{V' \to V\}$. D'après la
proposition
\ref{descent-proposition-fpqc-descent-quasi-coherent} de Descente,
cette donnée de descente est effective, et nous obtenons un
$\mathcal{O}_V$-module quasi-cohérent $\mathcal{F}_V$ sur $V_\etale$.
Pour voir que cela définit un faisceau quasi-cohérent sur $X_\etale$, il faut
montrer (d'après le
lemme \ref{lemma-characterize-quasi-coherent-small-etale})
que, pour tout morphisme $f : V_1 \to V_2$ dans $X_\etale$,
il existe un isomorphisme canonique
$c_f : f^*\mathcal{F}_{V_2} \to \mathcal{F}_{V_1}$
compatible aux compositions de morphismes. Nous omettons la vérification.
Nous omettons aussi de vérifier que cela définit un foncteur de la
catégorie de droite vers celle de gauche, inverse
du foncteur décrit ci-dessus.
\end{proof}

\begin{proposition}
\label{proposition-coherator}
Soit $S$ un schéma. Soit $X$ un espace algébrique sur $S$.
\begin{enumerate}
\item La catégorie $\QCoh(\mathcal{O}_X)$ est une catégorie
abélienne de Grothendieck. Par conséquent, $\QCoh(\mathcal{O}_X)$
possède suffisamment d'injectifs et toutes les limites.
\item Le foncteur d'inclusion
$\QCoh(\mathcal{O}_X) \to \textit{Mod}(\mathcal{O}_X)$
possède un adjoint à droite\footnote{Ce foncteur est parfois appelé
le {\it cohérateur}.}
$$
Q : \textit{Mod}(\mathcal{O}_X) \longrightarrow \QCoh(\mathcal{O}_X)
$$
tel que, pour tout faisceau quasi-cohérent $\mathcal{F}$, le morphisme d'adjonction
$Q(\mathcal{F}) \to \mathcal{F}$ soit un isomorphisme.
\end{enumerate}
\end{proposition}

\begin{proof}
Cette preuve reprend celle du cas des schémas ; voir la
proposition \ref{properties-proposition-coherator} de Propriétés.
Nous conseillons au lecteur de lire d'abord cette preuve.

\medskip\noindent
L'assertion (1) signifie que $\QCoh(\mathcal{O}_X)$ (a) admet toutes les colimites,
(b) que les colimites filtrantes y sont exactes et (c) qu'elle admet un générateur ; voir la
section \ref{injectives-section-grothendieck-conditions} d'Injectifs.
D'après le lemme \ref{lemma-properties-quasi-coherent},
les colimites dans $\QCoh(\mathcal{O}_X)$ existent et coïncident
avec celles de $\textit{Mod}(\mathcal{O}_X)$. D'après le
lemme \ref{sites-modules-lemma-limits-colimits} de Modules sur les sites,
les colimites filtrantes sont exactes. Les conditions (a) et (b) sont donc satisfaites.

\medskip\noindent
Pour construire un générateur, choisissons une présentation $X = U/R$ telle que
$(U, R, s, t, c)$ soit un
groupoïde en schémas \'etale ; en particulier, $s$ et $t$ sont des morphismes plats
de schémas. Choisissons un cardinal $\kappa$ comme dans le
lemme \ref{groupoids-lemma-colimit-kappa} de Groupoïdes.
Choisissons une famille $(\mathcal{E}_t, \alpha_t)_{t \in T}$ de
modules quasi-cohérents $\kappa$-engendrés sur
$(U, R, s, t, c)$ comme dans le
lemme \ref{groupoids-lemma-set-of-iso-classes} de Groupoïdes.
Soit $\mathcal{F}_t$ le module quasi-cohérent sur $X$ qui
correspond au module quasi-cohérent $(\mathcal{E}_t, \alpha_t)$ par
l'équivalence de catégories de la
proposition \ref{proposition-quasi-coherent}.
Nous voyons alors que tout module quasi-cohérent $\mathcal{H}$ est la
colimite filtrante de ses sous-modules quasi-cohérents qui sont isomorphes
à l'un des $\mathcal{F}_t$. Ainsi, $\bigoplus_t \mathcal{F}_t$ est
un générateur de $\QCoh(\mathcal{O}_X)$, ce qui établit (c).
Les assertions sur les limites et les injectifs valent dans toute catégorie
abélienne de Grothendieck ; voir le
théorème
\ref{injectives-theorem-injective-embedding-grothendieck} et le
lemme \ref{injectives-lemma-grothendieck-products} d'Injectifs.

\medskip\noindent
Preuve de (2). Pour construire $Q$, nous employons le procédé général suivant.
Étant donné un objet $\mathcal{F}$ de $\textit{Mod}(\mathcal{O}_X)$,
nous considérons le foncteur
$$
\QCoh(\mathcal{O}_X)^{opp} \longrightarrow \textit{Ens},\quad
\mathcal{G} \longmapsto \Hom_X(\mathcal{G}, \mathcal{F})
$$
Ce foncteur transforme les colimites en limites ;
il est donc représentable ; voir le
lemme \ref{injectives-lemma-grothendieck-brown} d'Injectifs.
Il existe ainsi un faisceau quasi-cohérent $Q(\mathcal{F})$
et un isomorphisme fonctoriel
$\Hom_X(\mathcal{G}, \mathcal{F}) = \Hom_X(\mathcal{G}, Q(\mathcal{F}))$
pour $\mathcal{G}$ dans $\QCoh(\mathcal{O}_X)$. D'après le lemme de Yoneda
(lemme \ref{categories-lemma-yoneda} de Catégories),
la construction $\mathcal{F} \leadsto Q(\mathcal{F})$ est
fonctorielle en $\mathcal{F}$. Par construction, $Q$ est un adjoint à
droite du foncteur d'inclusion.
Le fait que $Q(\mathcal{F}) \to \mathcal{F}$ soit un isomorphisme
lorsque $\mathcal{F}$ est quasi-cohérent résulte formellement du fait
que le foncteur d'inclusion
$\QCoh(\mathcal{O}_X) \to \textit{Mod}(\mathcal{O}_X)$
est pleinement fidèle.
\end{proof}






\section{Morphismes à valeurs dans des schémas}
\label{section-morphisms-to-schemes}

\noindent
Voici l'analogue du
lemme \ref{schemes-lemma-morphism-into-affine} de Schémas.

\begin{lemma}
\label{lemma-morphism-to-affine-scheme}
Soit $X$ un espace algébrique sur $\mathbf{Z}$.
Soit $T$ un schéma affine.
L'application
$$
\Mor(X, T)
\longrightarrow
\Hom(\Gamma(T, \mathcal{O}_T), \Gamma(X, \mathcal{O}_X))
$$
qui envoie $f$ sur $f^\sharp$ (sur les sections globales) est bijective.
\end{lemma}

\begin{proof}
Nous construisons l'application inverse.
Soit $\varphi : \Gamma(T, \mathcal{O}_T) \to \Gamma(X, \mathcal{O}_X)$
un homomorphisme d'anneaux. Choisissons une présentation $X = U/R$ ; voir la
définition \ref{spaces-definition-presentation} d'Espaces.
D'après le
lemme \ref{schemes-lemma-morphism-into-affine} de Schémas,
la composée
$$
\Gamma(T, \mathcal{O}_T) \to \Gamma(X, \mathcal{O}_X) \to
\Gamma(U, \mathcal{O}_U)
$$
correspond à un unique morphisme de schémas $g : U \to T$. Par le même lemme,
les deux composées $R \to U \to T$ sont égales. Nous obtenons donc un morphisme
$f : X = U/R \to T$ tel que $U \to X \to T$ soit égal à $g$. Par construction,
le diagramme
$$
\xymatrix{
\Gamma(U, \mathcal{O}_U) & \Gamma(X, \mathcal{O}_X) \ar[l] \\
& \Gamma(T, \mathcal{O}_T) \ar[lu]^{g^\sharp} \ar[u]^{\varphi}_{f^\sharp}
}
$$
est commutatif. Ainsi, $f^\sharp$ est égal à $\varphi$, car $U \to X$ est un
recouvrement \'etale et $\mathcal{O}_X$ est un faisceau sur $X_\etale$.
L'unicité de $f$ résulte de celle de $g$.
\end{proof}





\section{Quotients par des actions libres}
\label{section-quotient-by-free}

\noindent
Soit $S$ un schéma.
Soit $X$ un espace algébrique sur $S$.
Soit $G$ un groupe abstrait.
Soit $a : G \to \text{Aut}(X)$ un homomorphisme, c'est-à-dire que $a$ est une
{\it action} de $G$ sur $X$. Nous dirons que l'action est {\it libre}
si, pour tout schéma $T$ sur $S$, l'action
$$
G \times X(T) \longrightarrow X(T)
$$
est libre. (Nous ne pouvons pas employer un critère comme celui du
lemme \ref{spaces-lemma-quotient} d'Espaces,
car les points peuvent ne pas avoir de corps résiduels bien définis.)
Lorsque l'action est libre, nous allons construire le quotient $X/G$
comme espace algébrique. C'est un cas particulier du
lemme général \ref{bootstrap-lemma-quotient-free-action} d'Amorçage
que nous démontrerons plus tard.

\begin{lemma}
\label{lemma-quotient}
Soit $S$ un schéma.
Soit $X$ un espace algébrique sur $S$.
Soit $G$ un groupe abstrait opérant librement sur $X$.
Alors le faisceau quotient $X/G$ est un espace algébrique.
\end{lemma}

\begin{proof}
L'énoncé signifie que le faisceau $F$ associé au préfaisceau
$$
T \longmapsto X(T)/G
$$
est un espace algébrique. Pour le voir, nous allons construire une présentation.
Choisissons un schéma $U$ et un morphisme \'etale surjectif
$\varphi : U \to X$. Posons $V = \coprod_{g \in G} U$ et définissons
$\psi : V \to X$ comme $a(g) \circ \varphi$ sur la composante correspondant
à $g \in G$. Faisons agir $G$ sur $V$ en permutant les composantes : ainsi,
$g_0 \in G$ envoie la composante correspondant à $g$ sur celle correspondant
à $g_0g$ par le morphisme identité de $U$.
Alors $\psi$ est un morphisme $G$-équivariant ; nous nous ramenons donc au
cas traité dans le paragraphe suivant.

\medskip\noindent
Supposons qu'il existe une action de $G$ sur $U$ et que $U \to X$ soit surjectif,
\'etale et $G$-équivariant. Il existe alors une
action induite de $G$ sur $R = U \times_X U$, compatible avec les morphismes de
projection $t, s : R \to U$. Nous affirmons que
$$
X/G = U/\coprod\nolimits_{g \in G} R
$$
où l'application
$$
j : \coprod\nolimits_{g \in G} R
\longrightarrow
U \times_S U
$$
est donnée par $(r, g) \mapsto (t(r), g(s(r)))$. Remarquons que $j$ est un monomorphisme :
si $(t(r), g(s(r))) = (t(r'), g'(s(r')))$, alors
$t(r) = t(r')$ ; par conséquent, $r$ et $r'$ ont la même image dans $X$ par
$s$ et $t$, puis $g = g'$ (car $G$ opère librement sur $X$), puis
$s(r) = s(r')$, et enfin $r = r'$ (car $R$ est une relation d'équivalence sur $U$).
De plus, $j$ est une relation d'équivalence (détails omis).
Les deux projections $\coprod\nolimits_{g \in G} R \to U$ sont \'etales, puisque
$s$ et $t$ sont \'etales. Ainsi, $j$ est une relation d'équivalence \'etale
et $U/\coprod\nolimits_{g \in G} R$ est un espace algébrique d'après le
théorème \ref{spaces-theorem-presentation} d'Espaces.
Il existe une application
$$
U/\coprod\nolimits_{g \in G} R \longrightarrow X/G
$$
induite par l'application $U \to X$. Nous omettons la preuve qu'il s'agit d'un
isomorphisme de faisceaux.
\end{proof}






\begin{multicols}{2}[\section{Autres chapitres}]
\noindent
Préliminaires
\begin{enumerate}
\item \hyperref[introduction-section-phantom]{Introduction}
\item \hyperref[conventions-section-phantom]{Conventions}
\item \hyperref[sets-section-phantom]{Théorie des ensembles}
\item \hyperref[categories-section-phantom]{Catégories}
\item \hyperref[topology-section-phantom]{Topologie}
\item \hyperref[sheaves-section-phantom]{Faisceaux sur les espaces}
\item \hyperref[sites-section-phantom]{Sites et faisceaux}
\item \hyperref[stacks-section-phantom]{Champs}
\item \hyperref[fields-section-phantom]{Corps}
\item \hyperref[algebra-section-phantom]{Algèbre commutative}
\item \hyperref[brauer-section-phantom]{Groupes de Brauer}
\item \hyperref[homology-section-phantom]{Algèbre homologique}
\item \hyperref[derived-section-phantom]{Catégories dérivées}
\item \hyperref[simplicial-section-phantom]{Méthodes simpliciales}
\item \hyperref[more-algebra-section-phantom]{Compléments d'algèbre}
\item \hyperref[smoothing-section-phantom]{Lissage des morphismes d'anneaux}
\item \hyperref[modules-section-phantom]{Faisceaux de modules}
\item \hyperref[sites-modules-section-phantom]{Modules sur les sites}
\item \hyperref[injectives-section-phantom]{Objets injectifs}
\item \hyperref[cohomology-section-phantom]{Cohomologie des faisceaux}
\item \hyperref[sites-cohomology-section-phantom]{Cohomologie sur les sites}
\item \hyperref[dga-section-phantom]{Algèbre différentielle graduée}
\item \hyperref[dpa-section-phantom]{Algèbre à puissances divisées}
\item \hyperref[sdga-section-phantom]{Faisceaux différentiels gradués}
\item \hyperref[hypercovering-section-phantom]{Hyperrecouvrements}
\end{enumerate}
Schémas
\begin{enumerate}
\setcounter{enumi}{25}
\item \hyperref[schemes-section-phantom]{Schémas}
\item \hyperref[constructions-section-phantom]{Constructions de schémas}
\item \hyperref[properties-section-phantom]{Propriétés des schémas}
\item \hyperref[morphisms-section-phantom]{Morphismes de schémas}
\item \hyperref[coherent-section-phantom]{Cohomologie des schémas}
\item \hyperref[divisors-section-phantom]{Diviseurs}
\item \hyperref[limits-section-phantom]{Limites de schémas}
\item \hyperref[varieties-section-phantom]{Variétés}
\item \hyperref[topologies-section-phantom]{Topologies sur les schémas}
\item \hyperref[descent-section-phantom]{Descente}
\item \hyperref[perfect-section-phantom]{Catégories dérivées de schémas}
\item \hyperref[more-morphisms-section-phantom]{Compléments sur les morphismes}
\item \hyperref[flat-section-phantom]{Compléments sur la platitude}
\item \hyperref[groupoids-section-phantom]{Schémas en groupoïdes}
\item \hyperref[more-groupoids-section-phantom]{Compléments sur les schémas en groupoïdes}
\item \hyperref[etale-section-phantom]{Morphismes étales de schémas}
\end{enumerate}
Thèmes de théorie des schémas
\begin{enumerate}
\setcounter{enumi}{41}
\item \hyperref[chow-section-phantom]{Homologie de Chow}
\item \hyperref[intersection-section-phantom]{Théorie de l'intersection}
\item \hyperref[pic-section-phantom]{Schémas de Picard des courbes}
\item \hyperref[weil-section-phantom]{Théories cohomologiques de Weil}
\item \hyperref[adequate-section-phantom]{Modules adéquats}
\item \hyperref[dualizing-section-phantom]{Complexes dualisants}
\item \hyperref[duality-section-phantom]{Dualité pour les schémas}
\item \hyperref[discriminant-section-phantom]{Discriminants et différentes}
\item \hyperref[derham-section-phantom]{Cohomologie de de Rham}
\item \hyperref[local-cohomology-section-phantom]{Cohomologie locale}
\item \hyperref[algebraization-section-phantom]{Géométrie algébrique et formelle}
\item \hyperref[curves-section-phantom]{Courbes algébriques}
\item \hyperref[resolve-section-phantom]{Résolution des surfaces}
\item \hyperref[models-section-phantom]{Réduction semi-stable}
\item \hyperref[functors-section-phantom]{Foncteurs et morphismes}
\item \hyperref[equiv-section-phantom]{Catégories dérivées de variétés}
\item \hyperref[pione-section-phantom]{Groupes fondamentaux de schémas}
\item \hyperref[etale-cohomology-section-phantom]{Cohomologie étale}
\item \hyperref[crystalline-section-phantom]{Cohomologie cristalline}
\item \hyperref[proetale-section-phantom]{Cohomologie pro-étale}
\item \hyperref[relative-cycles-section-phantom]{Cycles relatifs}
\item \hyperref[more-etale-section-phantom]{Compléments de cohomologie étale}
\item \hyperref[trace-section-phantom]{La formule des traces}
\end{enumerate}
Espaces algébriques
\begin{enumerate}
\setcounter{enumi}{64}
\item \hyperref[spaces-section-phantom]{Espaces algébriques}
\item \hyperref[spaces-properties-section-phantom]{Propriétés des espaces algébriques}
\item \hyperref[spaces-morphisms-section-phantom]{Morphismes d'espaces algébriques}
\item \hyperref[decent-spaces-section-phantom]{Espaces algébriques décents}
\item \hyperref[spaces-cohomology-section-phantom]{Cohomologie des espaces algébriques}
\item \hyperref[spaces-limits-section-phantom]{Limites d'espaces algébriques}
\item \hyperref[spaces-divisors-section-phantom]{Diviseurs sur les espaces algébriques}
\item \hyperref[spaces-over-fields-section-phantom]{Espaces algébriques sur des corps}
\item \hyperref[spaces-topologies-section-phantom]{Topologies sur les espaces algébriques}
\item \hyperref[spaces-descent-section-phantom]{Descente et espaces algébriques}
\item \hyperref[spaces-perfect-section-phantom]{Catégories dérivées d'espaces}
\item \hyperref[spaces-more-morphisms-section-phantom]{Compléments sur les morphismes d'espaces}
\item \hyperref[spaces-flat-section-phantom]{Platitude sur les espaces algébriques}
\item \hyperref[spaces-groupoids-section-phantom]{Groupoïdes dans les espaces algébriques}
\item \hyperref[spaces-more-groupoids-section-phantom]{Compléments sur les groupoïdes dans les espaces}
\item \hyperref[bootstrap-section-phantom]{Amorçage}
\item \hyperref[spaces-pushouts-section-phantom]{Sommes amalgamées d'espaces algébriques}
\end{enumerate}
Thèmes de géométrie
\begin{enumerate}
\setcounter{enumi}{81}
\item \hyperref[spaces-chow-section-phantom]{Groupes de Chow des espaces}
\item \hyperref[groupoids-quotients-section-phantom]{Quotients de groupoïdes}
\item \hyperref[spaces-more-cohomology-section-phantom]{Compléments sur la cohomologie des espaces}
\item \hyperref[spaces-simplicial-section-phantom]{Espaces simpliciaux}
\item \hyperref[spaces-duality-section-phantom]{Dualité pour les espaces}
\item \hyperref[formal-spaces-section-phantom]{Espaces algébriques formels}
\item \hyperref[restricted-section-phantom]{Algébrisation des espaces formels}
\item \hyperref[spaces-resolve-section-phantom]{Résolution des surfaces revisitée}
\end{enumerate}
Théorie des déformations
\begin{enumerate}
\setcounter{enumi}{89}
\item \hyperref[formal-defos-section-phantom]{Théorie formelle des déformations}
\item \hyperref[defos-section-phantom]{Théorie des déformations}
\item \hyperref[cotangent-section-phantom]{Le complexe cotangent}
\item \hyperref[examples-defos-section-phantom]{Problèmes de déformation}
\end{enumerate}
Champs algébriques
\begin{enumerate}
\setcounter{enumi}{93}
\item \hyperref[algebraic-section-phantom]{Champs algébriques}
\item \hyperref[examples-stacks-section-phantom]{Exemples de champs}
\item \hyperref[stacks-sheaves-section-phantom]{Faisceaux sur les champs algébriques}
\item \hyperref[criteria-section-phantom]{Critères de représentabilité}
\item \hyperref[artin-section-phantom]{Axiomes d'Artin}
\item \hyperref[quot-section-phantom]{Espaces de Quot et de Hilbert}
\item \hyperref[stacks-properties-section-phantom]{Propriétés des champs algébriques}
\item \hyperref[stacks-morphisms-section-phantom]{Morphismes de champs algébriques}
\item \hyperref[stacks-limits-section-phantom]{Limites de champs algébriques}
\item \hyperref[stacks-cohomology-section-phantom]{Cohomologie des champs algébriques}
\item \hyperref[stacks-perfect-section-phantom]{Catégories dérivées de champs}
\item \hyperref[stacks-introduction-section-phantom]{Introduction aux champs algébriques}
\item \hyperref[stacks-more-morphisms-section-phantom]{Compléments sur les morphismes de champs}
\item \hyperref[stacks-geometry-section-phantom]{La géométrie des champs}
\end{enumerate}
Thèmes de théorie des espaces de modules
\begin{enumerate}
\setcounter{enumi}{107}
\item \hyperref[moduli-section-phantom]{Champs de modules}
\item \hyperref[moduli-curves-section-phantom]{Espaces de modules de courbes}
\end{enumerate}
Divers
\begin{enumerate}
\setcounter{enumi}{109}
\item \hyperref[examples-section-phantom]{Exemples}
\item \hyperref[exercises-section-phantom]{Exercices}
\item \hyperref[guide-section-phantom]{Guide bibliographique}
\item \hyperref[desirables-section-phantom]{Desiderata}
\item \hyperref[coding-section-phantom]{Style de programmation}
\item \hyperref[obsolete-section-phantom]{Obsolète}
\item \hyperref[fdl-section-phantom]{Licence de documentation libre GNU}
\item \hyperref[index-section-phantom]{Index généré automatiquement}
\end{enumerate}
\end{multicols}


\bibliography{my}
\bibliographystyle{amsalpha}

\end{document}
